% Redacción : artículo de prensa (120 palabras) ; traduction : la condition réalisable et irréalisable — Espagnol Tle A — Cours signé OnBuch+
% Programme officiel MINESEC. Compiler avec : tectonic E29-articulo-de-prensa-condicion-realizable-irrealizable.tex
\newcommand{\DOCMATIERE}{Espagnol}
\newcommand{\DOCNIVEAU}{Tle A}
\newcommand{\DOCMODULE}{Módulo 5 : Médias, communication et culture de la paix}
\newcommand{\DOCLECON}{E29}
\newcommand{\DOCTITRE}{Redacción : artículo de prensa (120 palabras) ; traduction : la condition réalisable et irréalisable}
% OnBuch+ — préambule « cours signés » ENRICHI (compilé avec tectonic / XeLaTeX)
% Système visuel « Pop » d'OnBuch+ : fond crème (jamais blanc), bordures encre
% épaisses, ombres « dures » décalées, titres Archivo Black en capitales.
% Version MATHÉMATIQUES : graphes (pgfplots/tikz), boîtes pédagogiques
% (définition, propriété, méthode, attention, le savais-tu, exercice résolu,
% exercice, TP, bilan), page de garde et signature OnBuch+ ancrée en bas.
%
% Macros attendues dans main.tex AVANT \input{preamble} :
%   \DOCMATIERE  \DOCNIVEAU  \DOCMODULE  \DOCLECON (numéro)  \DOCTITRE
\documentclass[11pt]{article}

\usepackage{fontspec}
\usepackage[spanish,french]{babel} % français = langue principale ; l'espagnol via \esp{..} / espagnol
\usepackage[a4paper,margin=2cm,top=3.4cm,bottom=2.8cm,headheight=1.4cm,footskip=1.3cm]{geometry}
\usepackage{amsmath,amssymb,amsfonts}
\usepackage{mathtools} % \xmapsto, \coloneqq... souvent utilisés par les modèles
\usepackage{cancel} % simplifications barrées (bilans de forces)
% \abs{z} (module, valeur absolue) et \norm{u} (norme), utilisés par les modèles
\providecommand{\abs}{}\let\abs\relax\DeclarePairedDelimiter{\abs}{\lvert}{\rvert}
\providecommand{\norm}{}\let\norm\relax\DeclarePairedDelimiter{\norm}{\lVert}{\rVert}
\usepackage[table]{xcolor} % [table] : fournit aussi \rowcolors (alternance de lignes)
\usepackage{pagecolor}
\usepackage{tikz}
\usetikzlibrary{arrows.meta,positioning,calc,shapes.geometric,shapes.misc,shapes.arrows,decorations.pathmorphing,decorations.pathreplacing,decorations.markings,patterns,shadows,mindmap,trees,fit,backgrounds,matrix,intersections,angles,quotes}
\usepackage{pgfplots}
\usepackage[version=4]{mhchem} % équations nucléaires \ce{^{235}_{92}U}
\usepackage[european,siunitx]{circuitikz} % schémas électriques
\pgfplotsset{compat=1.18}
\usepgfplotslibrary{fillbetween,groupplots} % aires entre courbes (calcul intégral)
\usepackage{enumitem}
\usepackage{fancyhdr}
% [most] charge toutes les bibliothèques tcolorbox (listings, minted, raster,
% documentation...) et gonfle inutilement la mémoire TeX sur un document long
% (~300 boîtes) : on ne charge que ce qui est réellement utilisé.
\usepackage[skins,breakable]{tcolorbox}
\usepackage{graphicx}
\usepackage{colortbl}
\usepackage{booktabs}
\usepackage{tabularx}
\usepackage{diagbox} % case d'en-tête diagonale des tableaux à double entrée
% Colonnes extensibles centrée / gauche / droite (C, L, R), souvent utilisées par les modèles
\newcolumntype{C}{>{\centering\arraybackslash}X}
\newcolumntype{L}{>{\raggedright\arraybackslash}X}
\newcolumntype{R}{>{\raggedleft\arraybackslash}X}
\usepackage{array}
\usepackage{multicol}
\usepackage{siunitx}
% parse-numbers=false : accepte aussi \SI{2,5 \times 10^{-3}}{...}
\sisetup{locale=FR,output-decimal-marker={,},group-minimum-digits=4,parse-numbers=false}
\DeclareSIUnit\ohm{\ensuremath{\symup{\Omega}}} % U+2126 absent de la police
\DeclareSIUnit\division{div} % réglages d'oscilloscope (ms/div, V/div)
\DeclareSIUnit\tr{tr} % tours (vitesse de rotation, ex. \SI{1500}{\tr\per\minute})
\DeclareSIUnit\molar{M} % concentration molaire (ex. \SI{3}{\molar}, \SI{1}{\micro\molar})
\DeclareSIUnit\an{an} % année (le modèle confond parfois avec \year, un compteur TeX) — ex. \SI{5}{\kilo\meter\per\an}
\DeclareSIUnit\FCFA{FCFA} % franc CFA (ex. \SI{500}{\FCFA\per\kilo\gram})
\DeclareSIUnit\degreeBrix{\textdegree Brix} % teneur en sucre d'un sirop/jus (réfractométrie)
\DeclareSIUnit\month{mois} % le modèle utilise parfois \month, un compteur TeX, comme unité de durée
\usepackage{qrcode}
% \degree : commande fréquemment utilisée par les modèles pour un angle ou une
% température en dehors de \SI{}{} (non fournie par les paquets chargés).
\providecommand{\degree}{^{\circ}}
% \qed (fin de démonstration), utilisé par les modèles sans amsthm
\providecommand{\qed}{\hfill$\square$}
% \barre : double barre des valeurs interdites dans un tableau de variations (tkz-tab)
\providecommand{\barre}{\ensuremath{\|}}
% environnement proof (amsthm non chargé) utilisé par les modèles
\ifdefined\proof\else\newenvironment{proof}[1][Démonstration]{\par\noindent\textit{#1.}\ }{\qed\par}\fi
\usepackage{lastpage}
\usepackage{hyperref}
\hypersetup{hidelinks}

% ============================================================
% Palette « Pop » OnBuch+ — mêmes hex que la classe P (plus_ui.dart)
% ============================================================
\definecolor{poporange}{HTML}{F59321}
\definecolor{popdark}{HTML}{E07A0C}
\definecolor{popcream}{HTML}{FFF6E8}
\definecolor{popink}{HTML}{1C1714}
\definecolor{poppurple}{HTML}{7A5AE0}
\definecolor{popgreen}{HTML}{2E9E6A}
\definecolor{poppink}{HTML}{E0457B}
\definecolor{popblue}{HTML}{2F7DE1}
\definecolor{popgold}{HTML}{C98A12}
\definecolor{popmuted}{HTML}{8A7F76}
\definecolor{popline}{HTML}{EADFD2}
\definecolor{popgray}{HTML}{8A7F76}
\colorlet{popgrey}{popgray}
% Alias tolérants (le modèle nomme parfois les couleurs différemment)
\colorlet{popwhite}{white}
\colorlet{popblack}{popink}
\colorlet{popred}{poppink}
\colorlet{popyellow}{popgold}
% Teintes claires (fonds de tableaux/encadrés) — le modèle nomme parfois
% les couleurs avec un suffixe L (ex. popblueL) sans les avoir définies.
\colorlet{poporangeL}{poporange!20}
\colorlet{popdarkL}{popdark!20}
\colorlet{poppurpleL}{poppurple!20}
\colorlet{popgreenL}{popgreen!20}
\colorlet{poppinkL}{poppink!20}
\colorlet{popblueL}{popblue!20}
\colorlet{popgoldL}{popgold!20}
\colorlet{popredL}{popred!20}
\colorlet{popgrayL}{popgray!20}
% Teintes claires pour fonds de tableaux / zones
\colorlet{popblueL}{popblue!10!white}
\colorlet{poporangeL}{poporange!14!white}
\colorlet{popgreenL}{popgreen!12!white}
\colorlet{poppurpleL}{poppurple!12!white}
\colorlet{poppinkL}{poppink!12!white}
\colorlet{popgoldL}{popgold!14!white}

\pagecolor{popcream}

% ============================================================
% Typographie
% ============================================================
\defaultfontfeatures{Ligatures=TeX}
\setmainfont{PlusJakartaSans-Medium.ttf}[Path=../../../fonts/,BoldFont=PlusJakartaSans-Bold.ttf]
% Maths en sans-serif (Fira Math) avec chiffres et lettres droites de Plus
% Jakarta, pour que les formules se fondent dans le texte.
\usepackage[math-style=ISO,bold-style=ISO]{unicode-math}
\setmathfont{FiraMath-Regular.otf}[Path=../../../fonts/]
\setmathfont{PlusJakartaSans-Medium.ttf}[Path=../../../fonts/,range={up/num,up/{Latin,latin},\mathup}]
\usepackage{icomma} % virgule décimale sans espace en mode math
% Caractères mathématiques Unicode absents des polices de texte (Plus Jakarta,
% Archivo Black) : rendus par leur équivalent mathématique, en texte comme en
% formule — sinon ils s'affichent en carré vide (ex. « Arithmétique dans ℤ »).
\usepackage{newunicodechar}
\newunicodechar{ℝ}{\ensuremath{\mathbb{R}}}
\newunicodechar{ℤ}{\ensuremath{\mathbb{Z}}}
\newunicodechar{ℂ}{\ensuremath{\mathbb{C}}}
\newunicodechar{ℕ}{\ensuremath{\mathbb{N}}}
\newunicodechar{ℚ}{\ensuremath{\mathbb{Q}}}
\newunicodechar{𝔻}{\ensuremath{\mathbb{D}}}
\newunicodechar{α}{\ensuremath{\alpha}}
\newunicodechar{β}{\ensuremath{\beta}}
\newunicodechar{γ}{\ensuremath{\gamma}}
\newunicodechar{δ}{\ensuremath{\delta}}
\newunicodechar{ε}{\ensuremath{\varepsilon}}
\newunicodechar{θ}{\ensuremath{\theta}}
\newunicodechar{λ}{\ensuremath{\lambda}}
\newunicodechar{μ}{\ensuremath{\mu}}
\newunicodechar{σ}{\ensuremath{\sigma}}
\newunicodechar{τ}{\ensuremath{\tau}}
\newunicodechar{φ}{\ensuremath{\varphi}}
\newunicodechar{ψ}{\ensuremath{\psi}}
\newunicodechar{ω}{\ensuremath{\omega}}
\newunicodechar{Σ}{\ensuremath{\Sigma}}
% majuscules produites par \MakeUppercase dans les titres (α-aminés -> Α)
\newunicodechar{Α}{\ensuremath{\alpha}}
\newunicodechar{Β}{\ensuremath{\beta}}
\newunicodechar{Γ}{\ensuremath{\Gamma}}
\newunicodechar{Δ}{\ensuremath{\Delta}}
\newunicodechar{Ε}{\ensuremath{\varepsilon}}
\newunicodechar{Θ}{\ensuremath{\Theta}}
\newunicodechar{Λ}{\ensuremath{\Lambda}}
\newunicodechar{Μ}{\ensuremath{\mu}}
\newunicodechar{Τ}{\ensuremath{\tau}}
\newunicodechar{Φ}{\ensuremath{\Phi}}
\newunicodechar{Ψ}{\ensuremath{\Psi}}
\newunicodechar{Ω}{\ensuremath{\Omega}}
\newunicodechar{↦}{\ensuremath{\mapsto}}
\newunicodechar{∈}{\ensuremath{\in}}
\newunicodechar{∉}{\ensuremath{\notin}}
\newunicodechar{∧}{\ensuremath{\wedge}}
\newunicodechar{⇔}{\ensuremath{\Leftrightarrow}}
\newunicodechar{⟺}{\ensuremath{\Longleftrightarrow}}
\newunicodechar{⇒}{\ensuremath{\Rightarrow}}
\newunicodechar{⟹}{\ensuremath{\Longrightarrow}}
\newunicodechar{∘}{\ensuremath{\circ}}
\newunicodechar{∪}{\ensuremath{\cup}}
\newunicodechar{∩}{\ensuremath{\cap}}
\newunicodechar{≡}{\ensuremath{\equiv}}
\newunicodechar{⊂}{\ensuremath{\subset}}
\newunicodechar{⊥}{\ensuremath{\perp}}
\newunicodechar{∃}{\ensuremath{\exists}}
\newunicodechar{∀}{\ensuremath{\forall}}
\newunicodechar{∛}{\ensuremath{\sqrt[3]{\,}}}
\newunicodechar{∜}{\ensuremath{\sqrt[4]{\,}}}
\newunicodechar{⃗}{\ensuremath{{}^{\to}}}
\newunicodechar{ⁿ}{\ifmmode^{n}\else\textsuperscript{\ensuremath{n}}\fi}
\newunicodechar{ˣ}{\ifmmode^{x}\else\textsuperscript{\ensuremath{x}}\fi}
\newunicodechar{ᵇ}{\ifmmode^{b}\else\textsuperscript{\ensuremath{b}}\fi}
\newunicodechar{ʳ}{\ifmmode^{r}\else\textsuperscript{\ensuremath{r}}\fi}
\newunicodechar{ᵘ}{\ifmmode^{u}\else\textsuperscript{\ensuremath{u}}\fi}
\newunicodechar{ʸ}{\ifmmode^{y}\else\textsuperscript{\ensuremath{y}}\fi}
\newunicodechar{ᵃ}{\ifmmode^{a}\else\textsuperscript{\ensuremath{a}}\fi}
\newunicodechar{ᵉ}{\ifmmode^{e}\else\textsuperscript{\ensuremath{e}}\fi}
\newunicodechar{ᵏ}{\ifmmode^{k}\else\textsuperscript{\ensuremath{k}}\fi}
\newunicodechar{ᵖ}{\ifmmode^{p}\else\textsuperscript{\ensuremath{p}}\fi}
\newunicodechar{ᴺ}{\ifmmode^{N}\else\textsuperscript{\ensuremath{N}}\fi}
\newunicodechar{ᶜ}{\ifmmode^{c}\else\textsuperscript{\ensuremath{c}}\fi}
\newunicodechar{ʰ}{\ifmmode^{h}\else\textsuperscript{\ensuremath{h}}\fi}
\newunicodechar{ᵠ}{\ifmmode^{q}\else\textsuperscript{\ensuremath{q}}\fi}
\newunicodechar{ₙ}{\ifmmode_{n}\else\textsubscript{\ensuremath{n}}\fi}
\newunicodechar{ₐ}{\ifmmode_{a}\else\textsubscript{\ensuremath{a}}\fi}
\newunicodechar{ᵢ}{\ifmmode_{i}\else\textsubscript{\ensuremath{i}}\fi}
\newunicodechar{ᵣ}{\ifmmode_{r}\else\textsubscript{\ensuremath{r}}\fi}
\newunicodechar{ₖ}{\ifmmode_{k}\else\textsubscript{\ensuremath{k}}\fi}
\newunicodechar{ᵦ}{\ifmmode_{\beta}\else\textsubscript{\ensuremath{\beta}}\fi}
\newunicodechar{ₓ}{\ifmmode_{x}\else\textsubscript{\ensuremath{x}}\fi}
\newfontfamily\popdisplay{ArchivoBlack-Regular.ttf}[Path=../../../fonts/]
\newfontfamily\poplogo{SpaceGrotesk-Bold.ttf}[Path=../../../fonts/]
\newfontfamily\popsemi{PlusJakartaSans-SemiBold.ttf}[Path=../../../fonts/]
\newfontfamily\popxbold{PlusJakartaSans-ExtraBold.ttf}[Path=../../../fonts/]
\newfontfamily\popxboldls{PlusJakartaSans-ExtraBold.ttf}[Path=../../../fonts/,LetterSpace=3.0]

\setlist[enumerate]{leftmargin=1.7em,itemsep=5pt,topsep=4pt}
\setlist[itemize]{leftmargin=1.7em,itemsep=4pt,topsep=4pt,label={\color{poporange}\textbullet}}
\setlength{\parindent}{0pt}
\setlength{\parskip}{5pt}
\renewcommand{\arraystretch}{1.25}

% pgfplots : style Pop par défaut
\pgfplotsset{
  every axis/.append style={axis line style={popink,line width=1pt},tick style={popink},
    grid style={popline,line width=0.5pt},label style={font=\small},tick label style={font=\footnotesize}},
  popcurve/.style={poporange,line width=1.8pt,no markers,smooth},
}
% Même style disponible dans un \draw TikZ simple (hors pgfplots)
\tikzset{popcurve/.style={poporange,line width=1.8pt}}
\tikzset{popgrid/.style={popline,very thin}}
\colorlet{popcurve}{poporange} % le nom du style est parfois utilisé comme couleur

% ============================================================
% Pill / squiggle
% ============================================================
\newcommand{\poppill}[3]{%
  \tikz[baseline=-0.55ex]\node[draw=popink,line width=1.2pt,rounded corners=2.6pt,%
    fill=#2,inner xsep=6.5pt,inner ysep=3pt]{%
    \popxboldls\fontsize{7}{7}\selectfont\color{#3}\MakeUppercase{#1}};%
}
\newcommand{\popsquiggle}[1]{%
  \tikz[baseline=-0.4ex]{
    \draw[#1,line width=2.6pt,line cap=round]
      plot [smooth] coordinates {(0,0.09) (0.34,0.01) (0.68,0.21) (1.02,0.01) (1.36,0.10)};
  }%
}
% Mot-clé surligné (marqueur orange sous le texte)
\newcommand{\cle}[1]{{\popxbold\color{popdark}#1}}

% ============================================================
% En-tête / pied
% ============================================================
\pagestyle{fancy}
\fancyhf{}
\renewcommand{\headrulewidth}{0pt}
\renewcommand{\footrulewidth}{0pt}
\lhead{\raisebox{-0.15ex}{\poplogo\fontsize{13}{13}\selectfont\color{popink}OnBuch\color{poporange}+}}
\rhead{\poppill{\DOCMATIERE\ \textbullet\ \DOCNIVEAU\ \textbullet\ Leçon \DOCLECON}{white}{popink}}
\lfoot{\popxbold\fontsize{7.6}{7.6}\selectfont\color{popmuted}\MakeUppercase{Cours signé OnBuch+}\ \textbullet\ {\popsemi\DOCTITRE}}
\rfoot{\tikz[baseline=-0.6ex]\node[rounded corners=8.5pt,fill=popink,minimum height=17pt,minimum width=17pt,inner xsep=5pt,inner sep=0]{\popxbold\fontsize{8}{8}\selectfont\color{popcream}\thepage\,/\,\pageref*{LastPage}};}

% ============================================================
% Titres
% ============================================================
\newcommand{\chaptitre}[2]{%
  \begin{tcolorbox}[enhanced,colback=poporange,colframe=popink,boxrule=2.6pt,arc=11pt,
      drop shadow=popink,left=15pt,right=15pt,top=12pt,bottom=11pt,before skip=3pt,after skip=18pt]
    \poppill{#2}{popink}{popcream}\par\vspace{9pt}
    {\raggedright\hyphenpenalty=10000\exhyphenpenalty=10000\popdisplay\fontsize{22}{24}\selectfont\color{popcream}\MakeUppercase{#1}\par}
  \end{tcolorbox}
}
% Section numérotée : pastille orange avec le numéro + titre Archivo
\newcounter{popsec}
\newcommand{\coursec}[1]{%
  \stepcounter{popsec}\setcounter{popsub}{0}%
  \par\vspace{16pt}\noindent
  \begin{minipage}{\linewidth}
  \tikz[baseline=-0.7ex]\node[draw=popink,line width=1.6pt,fill=poporange,rounded corners=4pt,minimum size=22pt,inner sep=0]{\popdisplay\fontsize{12}{12}\selectfont\color{popcream}\thepopsec};%
  \hspace{8pt}{\popdisplay\fontsize{15}{17}\selectfont\color{popink}\MakeUppercase{#1}}\par
  \vspace{4pt}\noindent{\color{poporange}\rule{36pt}{3.6pt}}
  \end{minipage}\par\nopagebreak\vspace{8pt}
}
\newcounter{popsub}
\newcommand{\courssub}[1]{%
  \stepcounter{popsub}%
  \par\vspace{9pt}\noindent{\popxbold\fontsize{11.5}{13}\selectfont\color{popink}{\color{poporange}\thepopsec.\thepopsub}\hspace{6pt}#1}\par\nopagebreak\vspace{4pt}
}

% ============================================================
% Boîtes pédagogiques — même recette : fond blanc, bordure épaisse colorée,
% ombre dure encre, badge plein. Titre optionnel [#1].
% ============================================================
\tcbset{popbase/.style={breakable,enhanced,colback=white,boxrule=2.3pt,arc=9pt,
  shadow={3pt}{-3pt}{0pt}{popink},left=14pt,right=14pt,top=11pt,bottom=11pt,before skip=11pt,after skip=15pt,
  fontupper=\popsemi\fontsize{10.6}{14.5}\selectfont\color{popink},
  fontlower=\popsemi\fontsize{10.6}{14.5}\selectfont\color{popink}}}
% Coche dessinée (le glyphe ✓ n'existe pas dans les polices du document)
\newcommand{\popcheck}[1]{\tikz[baseline=-0.5ex]\draw[#1,line width=1.6pt,line cap=round,line join=round](-0.13,0.0)--(-0.04,-0.09)--(0.13,0.1);}
\providecommand{\coche}{\popcheck{popgreen}}
\providecommand{\resultat}[1]{\par\smallskip\noindent\fcolorbox{popgreen}{popgreenL}{\parbox{\dimexpr\linewidth-2\fboxsep-2\fboxrule\relax}{\textbf{Résultat.} #1}}\par\smallskip}
\newcommand{\popboxhead}[3]{\poppill{#1}{#2}{white}\ifx\relax#3\relax\else\hspace{6pt}{\popxbold\fontsize{10.6}{12}\selectfont\color{popink}#3}\fi\par\vspace{7pt}}

\newtcolorbox{aretenir}[1][]{popbase,colframe=popgreen,before upper={\popboxhead{À retenir}{popgreen}{#1}}}
% Texte en espagnol (document de lecture, dialogue, poème, extrait) : cadre
% d'étude. Un titre optionnel (source, auteur) s'affiche dans l'en-tête.
\newtcolorbox{textebox}[1][]{popbase,colframe=popblue,before upper={\popboxhead{Texto}{popblue}{#1}\begin{otherlanguage}{spanish}},after upper={\end{otherlanguage}}}
% Marques ✓ / ✗ (forme correcte / fautive) souvent utilisées par les modèles
\providecommand{\cross}{\textcolor{poppink}{\ensuremath{\times}}}
\providecommand{\xmark}{\cross}\providecommand{\crossmark}{\cross}\providecommand{\cmark}{\ensuremath{\checkmark}}
% Ligne de réponse / justification à compléter (inventée par les modèles)
\providecommand{\justification}[1]{\par\noindent\textit{Justification~:} \dotfill\par}
% \textipa (package tipa non chargé) : la police ne couvre pas l'API -> simple passage
\providecommand{\textipa}[1]{#1}
% Soulignement (ulem non chargé) : \uline, \ul
\providecommand{\uline}[1]{\underline{#1}}\providecommand{\ul}[1]{\underline{#1}}
% Mot ou phrase en espagnol à l'intérieur d'un paragraphe français
\newcommand{\esp}[1]{\ifmmode\text{\textit{#1}}\else\begingroup\selectlanguage{spanish}\itshape #1\endgroup\fi}
\newtcolorbox{exemplebox}[1][]{popbase,colframe=poporange,before upper={\popboxhead{Exemple}{poporange}{#1}}}
\newtcolorbox{definition}[1][]{popbase,colframe=popblue,before upper={\popboxhead{Définition}{popblue}{#1}}}
\newtcolorbox{propriete}[1][]{popbase,colframe=poppurple,before upper={\popboxhead{Propriété}{poppurple}{#1}}}
\newtcolorbox{methode}[1][]{popbase,colframe=popgold,before upper={\popboxhead{Méthode}{popgold}{#1}}}
\newtcolorbox{attention}[1][]{popbase,colframe=poppink,before upper={\popboxhead{Attention}{poppink}{#1}}}
\newtcolorbox{experience}[1][]{popbase,colframe=popblue,colback=popblueL,before upper={\popboxhead{Expérience}{popblue}{#1}}}
% « Le savais-tu ? » : carte violette pleine (contexte, culture, Cameroun)
\newtcolorbox{savaistu}[1][]{popbase,colback=poppurple,colframe=popink,
  fontupper=\popsemi\fontsize{10.4}{14.2}\selectfont\color{white},
  before upper={\poppill{Le savais-tu\,?}{white}{poppurple}\ifx\relax#1\relax\else\hspace{6pt}{\popxbold\color{white}#1}\fi\par\vspace{7pt}}}
% Exercice résolu : énoncé en haut, \tcblower puis la solution sur fond crème
\newcounter{popexo}\newcounter{popexr}
\newtcolorbox{exoresolu}[1][]{popbase,colframe=poporange,colbacklower=popcream,
  segmentation style={popink,line width=1.2pt,dashed},
  before upper={\stepcounter{popexr}\popboxhead{Exercice résolu \thepopexr}{poporange}{#1}},
  before lower={\poppill{Solution}{popink}{popcream}\par\vspace{6pt}}}
% Exercice (non résolu en place) — la correction est dans la partie Corrigés
\newtcolorbox{exercice}[1][]{popbase,colframe=popink,
  before upper={\stepcounter{popexo}\popboxhead{Exercice \thepopexo}{popink}{#1}}}
% Corrigé
\newtcolorbox{corrige}[1][]{popbase,colframe=popgreen,colback=popgreenL,
  before upper={\popboxhead{Corrigé}{popgreen}{#1}}}
% Objectifs / prérequis / bilan
\newtcolorbox{objectifs}{popbase,colframe=popink,colback=poporangeL,
  before upper={\popboxhead{Objectifs de la leçon}{popink}{}}}
\newtcolorbox{prerequis}{popbase,colframe=popink,colback=popblueL,
  before upper={\popboxhead{Prérequis}{popblue}{}}}
\newtcolorbox{bilan}{popbase,colframe=popink,colback=white,boxrule=2.8pt,
  before upper={\popboxhead{Fiche bilan}{poporange}{L'essentiel à connaître pour l'examen}}}
% Figure encadrée avec légende
\newtcolorbox{popfigure}[1][]{enhanced,breakable=false,colback=white,colframe=popline,boxrule=1.4pt,arc=8pt,
  left=10pt,right=10pt,top=10pt,bottom=8pt,before skip=10pt,after skip=12pt,halign=center,
  lower separated=false,halign lower=center,
  fontlower=\popsemi\fontsize{9}{11.5}\selectfont\color{popmuted},
  #1}
\newcommand{\legende}[1]{\par\vspace{6pt}{\popsemi\fontsize{9}{11.5}\selectfont\color{popmuted}#1}}

% ============================================================
% Signature OnBuch+
% ============================================================
\newcommand{\poplogoinline}[1]{{\poplogo\fontsize{#1}{#1}\selectfont\color{popink}OnBuch\color{poporange}+}}
\newcommand{\signatureonbuch}{%
  \begin{tcolorbox}[enhanced,colback=popink,colframe=popink,boxrule=2.2pt,arc=10pt,
      drop shadow=poporange,left=14pt,right=14pt,top=10pt,bottom=10pt,before skip=0pt,after skip=0pt]
    \tikz[baseline=-0.7ex]\node[circle,fill=popgreen,draw=popcream,line width=1.3pt,minimum size=22pt,inner sep=0]{\popcheck{white}};%
    \hspace{9pt}\begin{minipage}[c]{0.62\linewidth}
      {\popxbold\fontsize{12}{13}\selectfont\color{popcream}Signé\ \textbullet\ {\poplogo OnBuch\color{poporange}+}}\par\vspace{2pt}
      {\raggedright\popsemi\fontsize{8}{10.5}\selectfont\color{popline}\MakeUppercase{Équipe pédagogique OnBuch+}\\\MakeUppercase{Conforme au programme officiel MINESEC}\par}
    \end{minipage}\hfill
    {\poplogo\fontsize{22}{22}\selectfont\color{popcream}OnBuch\color{poporange}+}
  \end{tcolorbox}%
}

% ============================================================
% Page de garde
% #1 titre · #2 matière · #3 niveau · #4 module · #5 n° leçon · #6 durée ·
% #7 description / objectifs (texte ou itemize)
% ============================================================
\newcommand{\pagegarde}[7]{%
  \thispagestyle{empty}
  \newgeometry{a4paper,margin=1.7cm,top=1.6cm,bottom=1.4cm}%
  \begin{tikzpicture}[remember picture,overlay]
    \fill[poporange!18] ([xshift=-3.2cm,yshift=-3.6cm]current page.north east) circle (4.6cm);
    \fill[poppurple!14] ([xshift=2.4cm,yshift=7.6cm]current page.south west) circle (3.8cm);
  \end{tikzpicture}%
  {\poplogo\fontsize{30}{30}\selectfont\color{popink}OnBuch\color{poporange}+}\hfill
  \raisebox{0.6ex}{\poppill{Cours signé \textbullet\ Premium}{popink}{popcream}}\par
  \vspace{1pt}\popsquiggle{poppurple}\par
  \vspace{8pt}
  \begin{tcolorbox}[enhanced,colback=poporange,colframe=popink,boxrule=2.8pt,arc=13pt,
      drop shadow=popink,left=18pt,right=18pt,top=12pt,bottom=13pt,after skip=10pt]
    \poppill{#2 \textbullet\ #3}{popink}{popcream}\hspace{5pt}\poppill{#4}{white}{popink}\par\vspace{8pt}
    {\popdisplay\fontsize{14}{14}\selectfont\color{popink}LEÇON #5}\par\vspace{4pt}
    {\raggedright\hyphenpenalty=10000\exhyphenpenalty=10000\popdisplay\fontsize{25}{27}\selectfont\color{popcream}\MakeUppercase{#1}\par}
    \vspace{8pt}
    {\popxbold\fontsize{9.5}{9.5}\selectfont\color{popink}\MakeUppercase{Durée conseillée : #6}}
  \end{tcolorbox}
  \begin{tcolorbox}[enhanced,colback=white,colframe=popink,boxrule=2.2pt,arc=10pt,
      drop shadow=popink,left=16pt,right=16pt,top=10pt,bottom=10pt,after skip=8pt]
    \poppill{Dans cette leçon}{poppurple}{white}\par\vspace{5pt}
    \popsemi\fontsize{9.3}{12.6}\selectfont\color{popink}#7
  \end{tcolorbox}
  \noindent\begin{minipage}[t]{0.487\linewidth}
    \begin{tcolorbox}[enhanced,colback=popgreen,colframe=popink,boxrule=2.2pt,arc=10pt,
        drop shadow=popink,left=11pt,right=11pt,top=10pt,bottom=10pt,height=3.1cm,valign=center]
      \tcbox[colback=white,colframe=popink,boxrule=1.3pt,arc=5pt,boxsep=0pt,nobeforeafter,
        left=3pt,right=3pt,top=3pt,bottom=3pt,valign=center]{\qrcode[height=1.9cm]{https://play.google.com/store/apps/details?id=com.luvvix.onbuch&hl=fr}}%
      \hspace{8pt}\begin{minipage}[c]{0.5\linewidth}
        {\popdisplay\fontsize{11}{12.5}\selectfont\color{white}\MakeUppercase{Télécharge OnBuch}}\par\vspace{4pt}
        {\raggedright\popsemi\fontsize{8.4}{11}\selectfont\color{white}Gratuit \textbullet\ Google Play\\Tuteur IA, annales, concours, résultats\par}
      \end{minipage}
    \end{tcolorbox}
  \end{minipage}\hfill
  \begin{minipage}[t]{0.487\linewidth}
    \begin{tcolorbox}[enhanced,colback=poppurple,colframe=popink,boxrule=2.2pt,arc=10pt,
        drop shadow=popink,left=11pt,right=11pt,top=10pt,bottom=10pt,height=3.1cm,valign=center]
      \tikz[baseline=-0.7ex]\node[circle,fill=white,draw=popink,line width=1.3pt,minimum size=1.6cm,inner sep=0]{\popdisplay\fontsize{24}{24}\selectfont\color{poppurple}+};%
      \hspace{8pt}\begin{minipage}[c]{0.6\linewidth}
        {\popdisplay\fontsize{11}{12.5}\selectfont\color{white}\MakeUppercase{Passe à OnBuch+}}\par\vspace{4pt}
        {\raggedright\popsemi\fontsize{8.4}{11}\selectfont\color{white}Tous les cours signés, Tuteur IA illimité, fiches PDF — onglet OnBuch+ de l'app\par}
      \end{minipage}
    \end{tcolorbox}
  \end{minipage}
  \vspace{8pt}
  \popcontenu
  \vfill
  \signatureonbuch
  \restoregeometry
  \newpage
}

% Rangée « Ce cours contient » (page de garde)
\newcommand{\poptile}[3]{% #1 couleur #2 gros texte #3 légende
  \begin{tcolorbox}[enhanced,colback=#1,colframe=popink,boxrule=2pt,arc=9pt,shadow={2.5pt}{-2.5pt}{0pt}{popink},
      left=6pt,right=6pt,top=9pt,bottom=9pt,halign=center,valign=center,height=1.75cm,width=\linewidth,nobeforeafter]
    {\popdisplay\fontsize{12.5}{13}\selectfont\color{white}\MakeUppercase{#2}}\par\vspace{3pt}
    {\centering\popsemi\fontsize{7.8}{9.5}\selectfont\color{white}#3\par}
  \end{tcolorbox}}
\newcommand{\popcontenu}{%
  \par\noindent{\popxboldls\fontsize{7.5}{7.5}\selectfont\color{popmuted}CE COURS SIGNÉ CONTIENT}\par\vspace{7pt}
  \noindent\begin{minipage}[t]{0.235\linewidth}\poptile{popblue}{Cours}{complet, illustré, pas à pas}\end{minipage}\hfill
  \begin{minipage}[t]{0.235\linewidth}\poptile{poporange}{Méthodes}{et exercices résolus}\end{minipage}\hfill
  \begin{minipage}[t]{0.235\linewidth}\poptile{poppink}{Exercices}{type examen + corrigés}\end{minipage}\hfill
  \begin{minipage}[t]{0.235\linewidth}\poptile{popgreen}{Fiche bilan}{l'essentiel à retenir}\end{minipage}\par}

% Sommaire
\newtcolorbox{sommaire}{enhanced,colback=white,colframe=popink,boxrule=2.2pt,arc=10pt,
  drop shadow=popink,left=18pt,right=18pt,top=13pt,bottom=13pt,before skip=6pt,after skip=20pt,
  before upper={\poppill{Sommaire}{poppurple}{white}\par\vspace{9pt}\popsemi\fontsize{10.6}{14.5}\selectfont\color{popink}}}

% Signature de fin de cours — poussée tout en bas de la dernière page
\newcommand{\findecours}{%
  \par\vfill\nopagebreak
  \begin{center}\popsquiggle{poporange}\end{center}\vspace{4pt}
  \signatureonbuch
}
\AtBeginDocument{\def\llbracket{\mathopen{[\mkern-3.5mu[}}\def\rrbracket{\mathclose{]\mkern-3.5mu]}}} % intervalles d'entiers (glyphe ⟦ absent de la police)
% Glyphes absents de Fira Math : redéfinis à partir de glyphes disponibles
\AtBeginDocument{%
  \def\setminus{\mathbin{\backslash}}\let\smallsetminus\setminus
  \def\vdots{\mathord{\vbox{\baselineskip3.2pt\lineskiplimit0pt\kern4pt\hbox{.}\hbox{.}\hbox{.}}}}%
  \def\ddots{\mathinner{\mkern1mu\raise6.5pt\hbox{.}\mkern2mu\raise3.6pt\hbox{.}\mkern2mu\raise0.7pt\hbox{.}\mkern1mu}}%
  \def\triangle{\mathord{\scalebox{0.9}{$\symup{\Delta}$}}}%
  \def\nabla{\mathord{\raisebox{\depth}{\scalebox{1}[-1]{$\symup{\Delta}$}}}}%
  \def\checkmark{\popcheck{popdark}}%
  \def\star{\mathbin{\ast}}\def\bigstar{\mathord{\ast}}%
  \def\diamond{\mathbin{\scalebox{0.6}{$\Diamond$}}}%
  \def\bigcup{\mathop{\mathchoice{\raisebox{-0.3em}{\scalebox{1.7}{$\cup$}}}{\scalebox{1.25}{$\cup$}}{\cup}{\cup}}\displaylimits}%
  \def\bigcap{\mathop{\mathchoice{\raisebox{-0.3em}{\scalebox{1.7}{$\cap$}}}{\scalebox{1.25}{$\cap$}}{\cap}{\cap}}\displaylimits}%
  \def\lesseqqgtr{\mathrel{\lessgtr}}\def\bigoplus{\mathop{\oplus}\displaylimits}%
}
\providecommand{\ssi}{si et seulement si} % abréviation fréquente
\AtBeginDocument{\providecommand{\wideparen}{\overparen}} % arc de cercle

\begin{document}

\pagegarde{Redacción : artículo de prensa (120 palabras) ; traduction : la condition réalisable et irréalisable}{Espagnol}{Tle A}{Módulo 5 : Médias, communication et culture de la paix}{E29}{2 h}{Cette leçon mixte de 2 heures apprend à rédiger un article de presse structuré (titre, chapeau, corps, conclusion) en répondant aux 5 W + H, et à traduire les trois types de condition (réalisable, irréalisable présent, irréalisable passé) avec les modes et temps corrects en espagnol.
\vspace{4pt}
\begin{multicols}{2}\small
\begin{itemize}
\item À la fin de cette leçon, je sais identifier les quatre parties d'un article de presse et les six questions clés (qui, quoi, où, quand, pourquoi, comment).
\item Je sais analyser un article modèle en relevant son lexique thématique (médias, paix, société) et sa structure.
\item Je sais rédiger un article de presse de 120 mots environ sur un sujet d'actualité camerounaise ou hispano-américaine.
\item Je sais distinguer la condition réalisable de l'irréalisable (présent et passé) à partir d'exemples authentiques.
\item Je sais conjuguer correctement les verbes après *si* : indicatif présent/futur, subjonctif imparfait/plus-que-parfait, conditionnel simple/composé.
\item Je sais traduire du français vers l'espagnol (thème) et de l'espagnol vers le français (version) des phrases à condition sans erreur de concordance des temps.
\end{itemize}\end{multicols}}

\begin{sommaire}
\begin{multicols}{2}
\begin{enumerate}
\item 1. L'article de presse : structure et questions clés
\item 2. Compréhension et lexique thématique de l'article modèle
\item 3. Production guidée : rédiger son article de presse (120 mots)
\item 4. Observation et découverte : la condition en espagnol (si-clauses)
\item 5. Règle complète, tableaux de formes et comparaison FR/ES
\item 6. Entraînement traduction (thème/version) et pièges du Bac
\item Activité d'intégration
\item Méthodes et exercices résolus
\item Exercices
\item Corrigés des exercices
\item Fiche bilan
\end{enumerate}
\end{multicols}
\end{sommaire}

\begin{objectifs}
\begin{itemize}
\item À la fin de cette leçon, je sais identifier les quatre parties d'un article de presse et les six questions clés (qui, quoi, où, quand, pourquoi, comment).
\item Je sais analyser un article modèle en relevant son lexique thématique (médias, paix, société) et sa structure.
\item Je sais rédiger un article de presse de 120 mots environ sur un sujet d'actualité camerounaise ou hispano-américaine.
\item Je sais distinguer la condition réalisable de l'irréalisable (présent et passé) à partir d'exemples authentiques.
\item Je sais conjuguer correctement les verbes après *si* : indicatif présent/futur, subjonctif imparfait/plus-que-parfait, conditionnel simple/composé.
\item Je sais traduire du français vers l'espagnol (thème) et de l'espagnol vers le français (version) des phrases à condition sans erreur de concordance des temps.
\item Je sais repérer et éviter les pièges fréquents : *si* + indicatif au lieu de subjonctif, confusion conditionnel/futur, omission du *se* impersonnel.
\end{itemize}
\end{objectifs}

\begin{prerequis}
\begin{itemize}
\item Conjugaison de l'indicatif présent, futur simple, conditionnel présent (Première).
\item Formation et emploi du subjonctif imparfait et plus-que-parfait (début Tle).
\item Vocabulaire de base sur les médias (periódico, titular, redactor, fuente) et la paix (conflicto, diálogo, reconciliación).
\item Structure de base d'un paragraphe argumentatif (idée directrice, connecteurs logiques).
\item Notion de phrase complexe avec subordonnée (conjonctions *que*, *cuando*, *porque*).
\end{itemize}
\end{prerequis}

\newpage

\coursec{L'article de presse : structure et questions clés}

Au Cameroun, comme dans tout le monde hispanophone, les journaux tels que \emph{El País}, \emph{La Nación} ou \emph{Cameroon Tribune} façonnent l'opinion publique chaque matin. Vous avez sans doute lu un article sur la CAN ou sur les défis de la jeunesse à Yaoundé. Mais comment structurer l'information pour qu'elle soit claire, complète et percutante en 120 mots ? Et comment exprimer en espagnol : « Si j'avais su, je serais venu » ou « Si tu viens, nous irons ensemble » ? Cette leçon vous donne les clés pour écrire comme un journaliste et maîtriser les phrases conditionnelles avec \esp{si}, indispensables au Bac. Dans cette première section, nous découvrons l'article de presse : sa définition, ses quatre parties obligatoires et la fameuse règle des questions clés.

\courssub{Qu'est-ce qu'un article de presse ?}

\begin{definition}[El artículo de prensa]
Un \cle{artículo de prensa} est un texte \textbf{informatif}, \textbf{objectif} et \textbf{bref}, publié dans un média écrit (journal, magazine) ou numérique (site d'information, blog d'actualité). Son but est de rapporter un fait réel et récent de manière claire, sans donner l'opinion personnelle du rédacteur. En espagnol, on distingue : \esp{la noticia} (la simple information), \esp{el reportaje} (le reportage approfondi) et \esp{el artículo de opinión} (l'article d'opinion).
\end{definition}

Observons d'abord un modèle avant d'en tirer les règles. Lisez attentivement cet article d'environ 120 mots, rédigé selon toutes les normes du genre journalistique : les quatre parties y sont signalées entre crochets pour vous aider à les repérer.

\begin{textebox}[Artículo modelo : « La juventud de Douala organiza un festival por la paz »]
\textbf{[TÍTULO]} La juventud de Douala organiza un festival por la paz

\textbf{[COPETE]} Más de dos mil jóvenes participaron ayer en el barrio de Bonabéri, en Douala, en un festival cultural dedicado a la cultura de paz.

\textbf{[CUERPO]} El evento, organizado por la asociación \emph{Juventud Unida}, reunió a artistas, estudiantes y líderes comunitarios. Según su presidenta, Amina Mbarga, « queremos demostrar que los jóvenes son actores de la reconciliación ». El festival incluyó conciertos, debates y un partido de fútbol entre barrios rivales. Las autoridades locales felicitaron a los organizadores por la iniciativa y prometieron apoyarla cada año.

\textbf{[CONCLUSIÓN]} El próximo festival tendrá lugar en Yaundé. Se espera una participación aún mayor.
\end{textebox}

\textbf{Glossaire :} \esp{el barrio} : le quartier ; \esp{jóvenes} : jeunes ; \esp{reunió a} : a rassemblé ; \esp{líderes comunitarios} : responsables communautaires ; \esp{la reconciliación} : la réconciliation ; \esp{incluyó} : a inclus, a comporté ; \esp{felicitaron a los organizadores} : ont félicité les organisateurs ; \esp{apoyarla} : la soutenir ; \esp{tendrá lugar} : aura lieu ; \esp{se espera} : on attend, on espère.

\textbf{Questions de compréhension :}
\begin{enumerate}
\item ¿Quién organiza el festival y dónde tiene lugar?
\item ¿Cuántos jóvenes participaron y cuándo?
\item ¿Qué actividades incluyó el evento?
\item ¿Qué prometieron las autoridades locales?
\end{enumerate}

\courssub{Les quatre parties obligatoires de l'article}

Tout article de presse espagnol — comme l'article modèle ci-dessus — repose sur une architecture fixe en quatre blocs. Le journaliste construit son texte comme une \textbf{pyramide inversée} : l'essentiel d'abord, les détails ensuite.

\begin{propriete}[La estructura del artículo de prensa]
\begin{enumerate}
\item \esp{El título} : le titre. Court, frappant, il annonce le sujet.
\item \esp{El copete} (ou \esp{la entradilla}) : le chapeau. Premier paragraphe qui résume toute la nouvelle.
\item \esp{El cuerpo} : le corps. Il développe les détails, les chiffres, les citations.
\item \esp{La conclusión} : la conclusion. Elle ouvre sur les conséquences ou l'avenir.
\end{enumerate}
\end{propriete}

\begin{definition}[El copete o entradilla]
Le \cle{copete} (ou \cle{entradilla}) est le premier paragraphe de l'article : il résume la nouvelle en répondant aux questions essentielles (qui, quoi, où, quand) en deux ou trois lignes seulement. Un lecteur pressé doit comprendre l'essentiel en lisant uniquement le titre et le copete.
\end{definition}

\begin{popfigure}
\begin{tikzpicture}[node distance=1.2cm, every node/.style={font=\small}]
\node[draw,rounded corners,fill=popblueL,align=center] (t) {TÍTULO};
\node[draw,rounded corners,fill=popgreenL,align=center,below of=t] (c) {COPETE / ENTRADILLA \\ (¿Quién? ¿Qué? ¿Dónde? ¿Cuándo?)};
\node[draw,rounded corners,fill=popgoldL,align=center,below of=c] (cu) {CUERPO \\ (¿Por qué? ¿Cómo? Detalles, citas, datos)};
\node[draw,rounded corners,fill=poporangeL,align=center,below of=cu] (co) {CONCLUSIÓN \\ (Consecuencias, futuro, cita final)};
\draw[->,thick] (t)--(c);
\draw[->,thick] (c)--(cu);
\draw[->,thick] (cu)--(co);
\end{tikzpicture}
\legende{La pyramide inversée de l'article de presse : de l'essentiel vers le détail.}
\end{popfigure}

\courssub{La règle des 5 W + H : les six questions clés}

Les journalistes anglo-saxons parlent des « 5 W » (\emph{Who, What, Where, When, Why}) auxquelles on ajoute le « H » (\emph{How}). En espagnol, cette règle devient un outil formidable pour vérifier qu'un article est complet.

\begin{center}
\begin{tabularx}{\linewidth}{|X|X|X|}
\hline
\rowcolor{popblueL} \textbf{Pregunta} & \textbf{Français} & \textbf{Respuesta en el artículo modelo} \\
\hline
\esp{¿Quién?} & Qui ? & \esp{la juventud de Douala, la asociación Juventud Unida} \\
\hline
\esp{¿Qué?} & Quoi ? & \esp{un festival cultural por la paz} \\
\hline
\esp{¿Dónde?} & Où ? & \esp{en el barrio de Bonabéri, en Douala} \\
\hline
\esp{¿Cuándo?} & Quand ? & \esp{ayer} \\
\hline
\esp{¿Por qué?} & Pourquoi ? & \esp{para demostrar que los jóvenes son actores de la reconciliación} \\
\hline
\esp{¿Cómo?} & Comment ? & \esp{con conciertos, debates y un partido de fútbol} \\
\hline
\end{tabularx}
\end{center}

\begin{aretenir}[Répartition des questions dans l'article]
Les quatre premières questions (\esp{¿Quién? ¿Qué? ¿Dónde? ¿Cuándo?}) trouvent leur réponse dans le \textbf{copete}. Les deux dernières (\esp{¿Por qué? ¿Cómo?}) sont développées dans le \textbf{cuerpo}, avec des données chiffrées (\esp{datos}) et des citations entre guillemets (\esp{citas}). La conclusion répond implicitement à la question « et maintenant ? ».
\end{aretenir}

\courssub{Noticia, reportaje, artículo de opinión : ne pas confondre}

\begin{center}
\begin{tabularx}{\linewidth}{|X|X|X|}
\hline
\rowcolor{popblueL} \textbf{Género} & \textbf{Características} & \textbf{Ejemplo de título} \\
\hline
\esp{La noticia} & Brève, factuelle, objective, sans opinion & \esp{« Incendio en el mercado de Mokolo »} \\
\hline
\esp{El reportaje} & Long, enquête approfondie, témoignages multiples & \esp{« La vida nocturna de los bendskins de Douala »} \\
\hline
\esp{El editorial / la columna} & Article d'opinion signé, subjectif, argumenté & \esp{« ¿Por qué el fútbol une a Camerún? »} \\
\hline
\end{tabularx}
\end{center}

Au Bac, on vous demande de rédiger une \esp{noticia} ou un court \esp{artículo de prensa} : restez donc objectif, sans « je » ni jugement personnel.

\courssub{Le style journalistique : registre formel et phrases efficaces}

L'article de presse exige un \textbf{registre formel} : phrases courtes (sujet–verbe–complément), vocabulaire précis, pas de familiarités. Un trait typique de l'espagnol journalistique est l'emploi de la \textbf{tournure impersonnelle} avec \esp{se} :

\esp{Se dice que el festival fue un éxito.} « On dit que le festival fut un succès. »

\esp{Se informó de que las autoridades apoyarán el evento.} « On a annoncé que les autorités soutiendront l'événement. »

\esp{Se espera una gran participación.} « On attend une forte participation. »

\begin{propriete}[Le titre : nominal ou au présent]
En espagnol, le titre de presse est souvent \textbf{nominal} (sans verbe conjugué) ou construit au \textbf{présent historique} : \esp{« Llega la paz a la región »}, \esp{« La juventud camerunesa impulsa la paz »}, \esp{« Camerún acoge la CAN »}. Évitez les phrases longues et les temps du passé dans le titre : le présent donne de l'immédiateté et du dynamisme.
\end{propriete}

\begin{exemplebox}[Títulos y copetes eficaces]
\textbf{Título :} \esp{« Camerún acoge la CAN »} « Le Cameroun accueille la CAN »

\textbf{Copete :} \esp{« La ceremonia de apertura tuvo lugar ayer en Yaundé ante miles de espectadores. »} « La cérémonie d'ouverture a eu lieu hier à Yaoundé devant des milliers de spectateurs. »

\textbf{Título nominal :} \esp{« Nuevo hospital en Bafoussam »} « Un nouvel hôpital à Bafoussam » (aucun verbe conjugué, efficacité maximale).
\end{exemplebox}

\begin{attention}[Erreurs fréquentes des francophones]
\begin{itemize}
\item \textbf{Le copete-littérature} : ne commencez jamais votre article par une introduction lyrique du type \esp{« Desde tiempos inmemoriales, el hombre... »}. Le journaliste donne l'information clé \textbf{immédiatement}.
\item \textbf{Le calque du français} : « avoir lieu » se traduit par \esp{tener lugar} (et non \esp{haber lugar}) ; « selon » se dit \esp{según} ; « plus de » devant un nombre se dit \esp{más de} (\esp{más de dos mil jóvenes}).
\item \textbf{Féliciter} : \esp{felicitar} se construit avec la personne félicitée, précédée de la préposition \esp{a} : \esp{felicitaron a los organizadores} « ils ont félicité les organisateurs ». On ne dit pas \esp{felicitaron la iniciativa}.
\item \textbf{L'opinion personnelle} : bannissez \esp{yo creo que}, \esp{en mi opinión} d'une \esp{noticia} : c'est le propre de l'\esp{artículo de opinión}.
\item \textbf{Les majuscules} : en espagnol, seul le premier mot du titre prend une majuscule : \esp{« La juventud de Douala organiza un festival »} (et non chaque mot comme en anglais).
\end{itemize}
\end{attention}

\begin{methode}[Escribir un artículo de prensa en seis pasos]
\begin{enumerate}
\item \textbf{Elegir el tema y el ángulo} : choisir le sujet et le point de vue (quoi de neuf ?).
\item \textbf{Responder las seis preguntas en borrador} : répondre au brouillon aux six questions clés.
\item \textbf{Redactar el título} : rédiger un titre court, nominal ou au présent.
\item \textbf{Escribir el copete} : maximum 30 mots, avec qui/quoi/où/quand.
\item \textbf{Desarrollar el cuerpo} : détails, chiffres, une citation entre guillemets.
\item \textbf{Cerrar con una consecuencia o apertura} : conclure par une conséquence ou une perspective d'avenir.
\end{enumerate}
\end{methode}

\begin{exoresolu}[Transformer des notes en copete]
\textbf{Énoncé.} Voici les notes d'un journaliste de \emph{Cameroon Tribune} : « lycée de Nkolbisson / concours de débat sur la paix / hier / 300 élèves / organisé par le club d'espagnol ». Rédigez le copete (maximum 30 mots).

\tcblower

\textbf{Solution détaillée.} On identifie d'abord les réponses aux quatre questions : ¿Quién? = \esp{trescientos alumnos} ; ¿Qué? = \esp{un concurso de debate sobre la paz} ; ¿Dónde? = \esp{en el liceo de Nkolbisson} ; ¿Cuándo? = \esp{ayer}. On assemble ensuite dans une phrase unique, au passé (\esp{participaron}, \esp{organizó}) :

\esp{« Trescientos alumnos participaron ayer en un concurso de debate sobre la cultura de paz en el liceo de Nkolbisson, organizado por el club de español. »}

« Trois cents élèves ont participé hier à un concours de débat sur la culture de la paix au lycée de Nkolbisson, organisé par le club d'espagnol. » (26 mots : consigne respectée, information complète et immédiate.)
\end{exoresolu}

\begin{savaistu}[La prensa en el mundo hispanohablante y en África]
L'Espagne compte de grands quotidiens comme \emph{El País} ou \emph{El Mundo}, et l'Amérique latine des journaux centenaires comme \emph{La Nación} (Argentine) ou \emph{El Universal} (Mexique). Saviez-vous que la Guinée équatoriale, voisine du Cameroun, est le seul pays d'Afrique subsaharienne dont l'espagnol est langue officielle ? On y publie des journaux en espagnol : un pont linguistique précieux pour les hispanistes camerounais !
\end{savaistu}

\begin{exercice}[À vous de jouer]
\begin{enumerate}
\item Classez ces trois titres selon le genre journalistique (noticia, reportaje, artículo de opinión) : \esp{« Terremoto en Chile: cientos de damnificados »} ; \esp{« Los jóvenes deben tomar el futuro en sus manos »} ; \esp{« Viaje al corazón de los mercados de Douala »}.
\item Rédigez le copete d'un article à partir de ces notes : « équipe de football du lycée de Bafoussam / victoire au tournoi interscolaire / samedi dernier / stade municipal ».
\item Réécrivez ce titre au style nominal espagnol : « Le gouvernement inaugure une nouvelle école à Garoua ».
\end{enumerate}
\end{exercice}

\begin{corrige}[Exercice « À vous de jouer »]
\begin{enumerate}
\item \esp{« Terremoto en Chile... »} = \esp{noticia} (fait, objectif) ; \esp{« Los jóvenes deben... »} = \esp{artículo de opinión} (jugement, « deben ») ; \esp{« Viaje al corazón de los mercados... »} = \esp{reportaje} (enquête, immersion).
\item Copete possible : \esp{« El equipo de fútbol del liceo de Bafoussam ganó el sábado pasado el torneo interescolar en el estadio municipal. »} « L'équipe de football du lycée de Bafoussam a gagné samedi dernier le tournoi interscolaire au stade municipal. »
\item \esp{« Nueva escuela en Garoua »} (titre nominal) ou \esp{« El Gobierno inaugura una nueva escuela en Garoua »} (présent historique).
\end{enumerate}
\end{corrige}

Vous maîtrisez désormais l'architecture de l'article de presse. Dans la section suivante, nous approfondirons la compréhension de l'article modèle et enrichirons votre lexique thématique des médias et de la culture de la paix.

\coursec{Compréhension et lexique thématique de l'article modèle}

Dans la section précédente, nous avons découvert la \cle{structure} de l'article de presse (titre, chapeau, corps, conclusion) et les six questions clés (\esp{quién, qué, dónde, cuándo, por qué, cómo}). Il est temps maintenant de passer de la théorie à la pratique : nous allons \cle{lire attentivement} un article modèle, vérifier qu'il répond bien à toutes les questions, puis en extraire le \cle{vocabulaire} et les \cle{procédés de langue} que vous devrez réutiliser dans votre propre rédaction. Un bon rédacteur est d'abord un bon lecteur.

\courssub{Lecture silencieuse de l'article modèle}

Lisez l'article suivant en silence, une première fois pour le sens général, une deuxième fois en soulignant les mots que vous ne connaissez pas.

\begin{textebox}[Article modèle — 120 palabras]
\textbf{La juventud de Douala organiza un festival por la paz}

Douala, 15 marzo (EFE). — Más de 500 jóvenes de diez asociaciones participaron ayer en el Festival por la Paz, celebrado en el parque Bonanjo. El evento, organizado por la Red Juvenil por la Convivencia, busca « fomentar el diálogo entre comunidades » tras los incidentes de enero. « La paz se construye cada día », declaró la coordinadora, Marie Nkeng. Hubo música, teatro y talleres. Las autoridades locales prometieron apoyo anual.

La próxima edición incluirá a escuelas rurales.
\end{textebox}

\textbf{Glossaire :}

\begin{itemize}
\item \esp{más de 500 jóvenes} : plus de 500 jeunes
\item \esp{participaron} : ont participé (prétérit de \esp{participar})
\item \esp{celebrado en} : qui s'est tenu à / célébré à
\item \esp{busca fomentar} : cherche à encourager
\item \esp{tras los incidentes} : après les incidents
\item \esp{se construye} : se construit (présent, passif réflexe)
\item \esp{declaró} : a déclaré
\item \esp{hubo} : il y a eu (prétérit de \esp{haber})
\item \esp{talleres} : ateliers
\item \esp{prometieron apoyo anual} : ont promis un soutien annuel
\item \esp{la próxima edición incluirá} : la prochaine édition inclura
\end{itemize}

\courssub{Compréhension globale : sujet, angle, sources, ton}

Avant de compter les mots ou de chercher le vocabulaire, un lecteur professionnel se pose quatre questions globales.

\begin{enumerate}
\item \textbf{Le sujet} : de quoi parle l'article ? Ici, d'un \esp{festival por la paz} organisé par des jeunes à Douala.
\item \textbf{L'angle} : quel aspect est privilégié ? L'article ne parle pas des incidents eux-mêmes, mais de la \emph{réponse pacifique} de la jeunesse : c'est un angle \cle{positif} et tourné vers la solution.
\item \textbf{Les sources citées} : qui parle dans l'article ?
\end{enumerate}

\begin{definition}[La fuente : la source]
La \esp{fuente} est l'origine de l'information : un témoin (\esp{testigo}), un document, un communiqué (\esp{comunicado}), une déclaration officielle. Dans notre article, deux sources sont citées : \esp{declaró la coordinadora, Marie Nkeng} (source directe, entre guillemets) et \esp{las autoridades locales prometieron apoyo anual} (source institutionnelle). La mention \esp{(EFE)} indique l'agence de presse qui fournit la dépêche.
\end{definition}

\begin{enumerate}
\setcounter{enumi}{3}
\item \textbf{Le ton général} : l'article est \cle{informatif} et \cle{factuel}. Le rédacteur ne donne pas son opinion ; il rapporte des faits, des chiffres (\esp{más de 500 jóvenes}, \esp{diez asociaciones}) et des paroles citées. Le ton est sobre, légèrement optimiste grâce à la citation finale.
\end{enumerate}

\begin{savaistu}[Les agences de presse]
En Espagne, l'\esp{Agencia EFE} est l'équivalent de l'AFP (Agence France-Presse) : elle rédige des dépêches reprises par tous les journaux. C'est pourquoi nos articles modèles commencent par \esp{Douala, 15 marzo (EFE)}. Au Cameroun, les sources d'information les plus citées sont \emph{Cameroon Tribune} (presse écrite) et la \emph{CRTV} (radio-télévision). En Guinée équatoriale, pays hispanophone voisin, on cite la radio nationale et l'agence officielle de Malabo.
\end{savaistu}

\courssub{Compréhension détaillée : les six réponses}

Vérifions que l'article modèle répond bien aux six questions clés vues dans la section 1. Ce contrôle est un réflexe à acquérir : avant de rendre votre propre article au Bac, relisez-le avec cette grille.

\begin{center}
\begin{tabularx}{\linewidth}{|l|X|X|}
\hline
\rowcolor{popblueL} \textbf{Pregunta} & \textbf{Respuesta en el artículo} & \textbf{Français} \\
\hline
\esp{¿Quién?} & \esp{más de 500 jóvenes de diez asociaciones ; la Red Juvenil ; Marie Nkeng} & qui ? : les jeunes, la coordinatrice \\
\hline
\esp{¿Qué?} & \esp{el Festival por la Paz} & quoi ? : un festival pour la paix \\
\hline
\esp{¿Dónde?} & \esp{en el parque Bonanjo, Douala} & où ? : au parc Bonanjo \\
\hline
\esp{¿Cuándo?} & \esp{ayer ; 15 marzo} & quand ? : hier, le 15 mars \\
\hline
\esp{¿Por qué?} & \esp{fomentar el diálogo tras los incidentes de enero} & pourquoi ? : favoriser le dialogue après les incidents \\
\hline
\esp{¿Cómo?} & \esp{hubo música, teatro y talleres} & comment ? : musique, théâtre, ateliers \\
\hline
\end{tabularx}
\end{center}

\begin{aretenir}[Le réflexe du correcteur]
Au Bac, le correcteur vérifie que votre article répond aux six questions en moins de 120 mots. L'article modèle le prouve : on peut tout dire en deux courts paragraphes si chaque phrase apporte une information nouvelle.
\end{aretenir}

\courssub{Extraction du lexique par champs sémantiques}

\begin{methode}[Classer le vocabulaire d'un texte]
\begin{enumerate}
\item Surligner les mots-clés du texte (noms, verbes, adjectifs porteurs de sens).
\item Les classer en quatre colonnes : \textbf{Médias}, \textbf{Paix}, \textbf{Jeunesse}, \textbf{Événement}.
\item Ajouter la traduction française exacte de chaque mot.
\item Mémoriser par colonne : 5 verbes, 5 noms et 3 adjectifs au minimum.
\end{enumerate}
\end{methode}

\begin{center}
\begin{tabularx}{\linewidth}{|X|X|X|}
\hline
\rowcolor{popgreenL} \textbf{Campo} & \textbf{Español} & \textbf{Français} \\
\hline
Médias & \esp{el titular, el redactor, la fuente, la edición, el artículo, la agencia} & le titre, le rédacteur, la source, l'édition, l'article, l'agence \\
\hline
Paix & \esp{la reconciliación, el diálogo, la convivencia, la no violencia, la tolerancia} & la réconciliation, le dialogue, la coexistence, la non-violence, la tolérance \\
\hline
Jeunesse & \esp{el colectivo, la asociación, el voluntariado, el liderazgo, el empoderamiento} & le collectif, l'association, le bénévolat, le leadership, l'autonomisation \\
\hline
Événement & \esp{el festival, el escenario, la actuación, la asistencia, la organización} & le festival, la scène, la représentation, la participation, l'organisation \\
\hline
\end{tabularx}
\end{center}

\begin{popfigure}
\begin{tikzpicture}[mindmap,level 1 concept/.append style={sibling angle=90,level distance=3.5cm}]
\node[concept,fill=popblue!30, align=center]{ARTÍCULO:\\PAZ Y JUVENTUD}
child[concept color=popgreen!50] {node[concept,align=center] {MEDIOS\\periódico, titular,\\redactor, fuente,\\edición}}
child[concept color=poporange!50] {node[concept,align=center] {PAZ\\reconciliación, diálogo,\\convivencia,\\no violencia,\\tolerancia}}
child[concept color=poppink!60] {node[concept,align=center] {JUVENTUD\\colectivo, asociación,\\voluntariado, liderazgo,\\empoderamiento}}
child[concept color=poppurple!50] {node[concept,align=center] {EVENTO\\festival, escenario,\\actuación, asistencia,\\organización}};
\end{tikzpicture}
\legende{Carte mentale des quatre champs sémantiques de l'article modèle.}
\end{popfigure}

\begin{attention}[Asistir / asistencia : le piège classique]
Ne confondez pas \esp{asistir} (verbe : assister à, être présent) et \esp{asistencia} (nom : la participation, le public). \esp{Asistió el alcalde.} « Le maire a assisté (à la cérémonie). » \esp{La asistencia fue masiva.} « La participation a été massive. » Erreur fréquente : écrire \esp{la asistencia asistió}, ce qui ne veut rien dire. Autre piège : \esp{asistir} ne signifie jamais « aider » ; « aider » se dit \esp{ayudar}.
\end{attention}

\courssub{Les connecteurs logiques de l'article}

Un article de presse n'est pas une liste de faits : les phrases sont reliées par des \cle{connecteurs} qui guident le lecteur. Notre article modèle est trop bref pour tous les contenir, mais voici les quatre connecteurs que le jury attend dans vos copies.

\begin{center}
\begin{tabularx}{\linewidth}{|l|X|X|}
\hline
\rowcolor{poporangeL} \textbf{Conector} & \textbf{Valor lógico} & \textbf{Ejemplo traducido} \\
\hline
\esp{por tanto} & conséquence & \esp{Hubo incidentes; por tanto, los jóvenes reaccionaron.} « Il y a eu des incidents ; par conséquent, les jeunes ont réagi. » \\
\hline
\esp{además} & addition & \esp{Hubo música; además, se organizaron talleres.} « Il y a eu de la musique ; de plus, des ateliers ont été organisés. » \\
\hline
\esp{sin embargo} & opposition & \esp{Llovió; sin embargo, el festival continuó.} « Il a plu ; cependant, le festival a continué. » \\
\hline
\esp{finalmente} & clôture & \esp{Finalmente, las autoridades prometieron apoyo.} « Finalement, les autorités ont promis leur soutien. » \\
\hline
\end{tabularx}
\end{center}

\begin{attention}[Faux amis et calques]
\esp{Por tanto} s'écrit en deux mots et signifie « par conséquent » ; ne le confondez pas avec \esp{por lo tanto} (même sens, plus soutenu) ni avec le français « pour tant ». Attention aussi : « finalement » au sens de « à la fin de l'événement » se traduit \esp{finalmente} ou \esp{al final}, mais jamais \esp{finamente}.
\end{attention}

\courssub{Analyse du registre : la langue objective du journaliste}

Relisons l'article : \esp{participaron}, \esp{busca}, \esp{declaró}, \esp{hubo}, \esp{prometieron}, \esp{incluirá}. Qu'observe-t-on ? Tous les verbes sont à la \textbf{troisième personne} ; le journaliste ne dit jamais \esp{yo} ni \esp{nosotros}. C'est la marque du \cle{registre informatif}.

\begin{propriete}[Le passif réflexe : \esp{se} + verbe à la 3e personne]
L'espagnol remplace souvent le passif français par le \cle{passif réflexe} : \esp{se} + verbe à la 3e personne (singulier ou pluriel selon le sujet).
\begin{itemize}
\item \esp{Se celebró el festival.} « Le festival a eu lieu / a été célébré. »
\item \esp{Se busca fomentar el diálogo.} « On cherche à favoriser le dialogue. »
\item \esp{Se organizaron talleres.} « Des ateliers ont été organisés. »
\item \esp{Se construye la paz cada día.} « La paix se construit chaque jour. »
\end{itemize}
\textbf{Accord} : le verbe s'accorde avec le nom qui suit : \esp{se celebró el festival} (singulier) mais \esp{se organizaron talleres} (pluriel).
\end{propriete}

Autres marques du registre journalistique présentes dans le modèle :
\begin{itemize}
\item \textbf{Tournures impersonnelles} : \esp{hubo música} « il y a eu de la musique » (pas de sujet humain).
\item \textbf{Peu d'adjectifs subjectifs} : le journaliste écrit \esp{más de 500 jóvenes} (fait mesurable), jamais \esp{un festival maravilloso y fantástico} (opinion). Les adjectifs subjectifs (\esp{maravilloso, estupendo, terrible}) appartiennent au commentaire, pas à l'information.
\item \textbf{Paroles citées entre guillemets} : \esp{« La paz se construye cada día », declaró la coordinadora.} Notez l'inversion sujet-verbe après la citation, typique de la presse : \esp{declaró la coordinadora} et non \esp{la coordinadora declaró}.
\end{itemize}

\begin{exoresolu}[Repérage dans l'article modèle]
Énoncé : 1) Relevez dans l'article deux verbes à la 3e personne du prétérit. 2) Identifiez la source directe et la source institutionnelle. 3) Trouvez le connecteur de temps qui ouvre le dernier paragraphe et proposez \esp{sin embargo} dans une phrase sur le même sujet.
\tcblower
Solution : 1) \esp{participaron} (ils ont participé) et \esp{prometieron} (ils ont promis) ; on peut ajouter \esp{declaró} et \esp{hubo}. 2) Source directe : \esp{la coordinadora, Marie Nkeng}, citée entre guillemets ; source institutionnelle : \esp{las autoridades locales}. 3) Le dernier paragraphe s'ouvre sur \esp{la próxima edición} (repère temporel). Avec \esp{sin embargo} : \esp{El parque era pequeño; sin embargo, acogió a más de 500 jóvenes.} « Le parc était petit ; cependant, il a accueilli plus de 500 jeunes. »
\end{exoresolu}

\begin{exercice}[À vous de jouer]
1) Classez ces mots dans les quatre colonnes (Médias, Paix, Jeunesse, Événement) : \esp{titular, taller, voluntariado, tolerancia, escenario, fuente, diálogo, asociación}. 2) Traduisez : « Le festival a été organisé par les jeunes ; cependant, les autorités ont aidé. » 3) Relevez dans l'article modèle les six réponses aux questions clés et recopiez-les en espagnol.
\end{exercice}

\begin{corrige}[Exercice 2.6]
1) Médias : \esp{titular, fuente} ; Paix : \esp{tolerancia, diálogo} ; Jeunesse : \esp{voluntariado, asociación} ; Événement : \esp{taller, escenario}. 2) \esp{El festival fue organizado por los jóvenes; sin embargo, las autoridades ayudaron.} (on accepte le passif réflexe : \esp{se organizó}). 3) Voir le tableau de la sous-section 2.3 : \esp{quién} : \esp{más de 500 jóvenes} ; \esp{qué} : \esp{el Festival por la Paz} ; \esp{dónde} : \esp{parque Bonanjo} ; \esp{cuándo} : \esp{ayer} ; \esp{por qué} : \esp{fomentar el diálogo} ; \esp{cómo} : \esp{música, teatro y talleres}.
\end{corrige}

\begin{aretenir}[Bilan de la section]
Un article de presse informe sans juger : troisième personne, passif réflexe \esp{se}, sources citées, chiffres précis, connecteurs logiques (\esp{por tanto, además, sin embargo, finalmente}). Vous disposez maintenant du lexique et des outils de langue : la section suivante vous guidera pas à pas pour rédiger votre propre article de 120 mots.
\end{aretenir}

\coursec{Production guidée : rédiger son article de presse (120 mots)}

Dans la section précédente, nous avons lu, commenté et exploité lexicalement un article de presse modèle : nous en avons dégagé la structure (titre, chapeau, corps, conclusion) et repéré les réponses aux six questions fondamentales. Il est temps maintenant de passer de la lecture à l'\cle{écriture} : vous allez produire votre propre article de presse, comme on vous le demandera au baccalauréat. Suivez le guide, étape par étape.

\courssub{Choisir son sujet}

La première décision est le choix du sujet. Au Bac, le sujet est souvent imposé, mais il est toujours lié à l'actualité ou à la vie sociale. Entraînons-nous avec des thèmes proches de votre quotidien et du monde hispanophone.

\begin{exemplebox}[Cinq sujets d'actualité pour s'entraîner]
\begin{enumerate}
\item \esp{La CAN 2025 : la selección camerunesa se prepara para el torneo.}
\item \esp{La rentrée scolaire : más de tres millones de alumnos vuelven a las aulas en Camerún.}
\item \esp{El clima : lluvias intensas provocan inundaciones en Douala.}
\item \esp{La cultura : un festival de música urbana reúne a jóvenes de Yaoundé.}
\item \esp{La paz : jornada escolar por la cultura de la paz en un liceo bilingüe.}
\end{enumerate}
\end{exemplebox}

\begin{attention}[Article ou essai ?]
Le piège classique du candidat francophone est d'écrire un \emph{essai d'opinion} (« Je pense que... », « À mon avis... ») au lieu d'un \emph{article factuel}. Un article de presse rapporte des faits, pas des opinions personnelles. Solution : zéro \esp{yo}, zéro \esp{creo que}. Pour introduire une information, utilisez les formules d'objectivité : \esp{se dice} (on dit), \esp{según fuentes oficiales} (selon des sources officielles), \esp{declaró el director} (a déclaré le directeur), \esp{afirmó}, \esp{anunció}, \esp{explicó}.
\end{attention}

\courssub{Remplir la fiche des six questions}

Avant d'écrire la moindre phrase, le journaliste remplit sa fiche d'enquête. C'est la garantie de ne rien oublier. Recopiez et complétez ce tableau pour le sujet que vous avez choisi.

\begin{center}
\begin{tabularx}{\linewidth}{|l|X|X|}
\hline
\rowcolor{popblueL} \textbf{Pregunta} & \textbf{Contenido} & \textbf{Mi respuesta (ejemplo : jornada de la paz)} \\
\hline
\esp{¿Quién?} (qui) & les acteurs & \esp{los alumnos del Liceo Bilingüe y una ONG local} \\
\hline
\esp{¿Qué?} (quoi) & l'événement & \esp{una marcha y talleres por la paz} \\
\hline
\esp{¿Dónde?} (où) & le lieu & \esp{en el patio del liceo, Yaoundé} \\
\hline
\esp{¿Cuándo?} (quand) & la date & \esp{el 21 de septiembre} \\
\hline
\esp{¿Por qué?} (pourquoi) & la cause & \esp{por un clima de tensión en el barrio} \\
\hline
\esp{¿Cómo?} (comment) & le déroulement & \esp{con cantos, discursos y dibujos} \\
\hline
\end{tabularx}
\end{center}

\courssub{Le plan de production en un coup d'œil}

La figure suivante résume tout le parcours de rédaction, du sujet à la conclusion, avec l'exemple de la « Journée de la paix à l'école ».

\begin{popfigure}
\begin{tikzpicture}[node distance=0.8cm]
\node[draw=popblue,rounded corners,fill=popblueL,align=center] (s) {SUJET : « Journée de la paix\\à l'école »};
\node[draw=popgreen,rounded corners,fill=popgreenL,right=1.2cm of s,align=center] (q) {FICHE 6 QUESTIONS\\Quién : élèves / ONG\\Qué : marche / ateliers\\Dónde : lycée Bilingue\\Cuándo : 21 sept.\\Por qué : climat tendu\\Cómo : chants, discours};
\node[draw=popgold,rounded corners,fill=popgoldL,below=of s,align=center] (t) {TÍTULO : « Estudiantes\\marchan por la paz »};
\node[draw=poporange,rounded corners,fill=poporangeL,below=of q,align=center] (c) {COPETE\\(28 palabras)};
\node[draw=poppink,rounded corners,fill=poppinkL,below=of t,align=center] (cu) {CUERPO\\(2 párrafos, 60 pal.)};
\node[draw=poppurple,rounded corners,fill=poppurpleL,below=of c,align=center] (co) {CONCLUSIÓN\\(15 pal.)};
\draw[->,thick,popdark] (s)--(q);
\draw[->,thick,popdark] (q)--(c);
\draw[->,thick,popdark] (s)--(t);
\draw[->,thick,popdark] (t)--(cu);
\draw[->,thick,popdark] (c)--(co);
\draw[->,thick,popdark] (cu)--(co);
\end{tikzpicture}
\legende{De la fiche des six questions à l'article complet : le parcours de rédaction.}
\end{popfigure}

\begin{methode}[Étape par étape : rédiger son article]
\begin{enumerate}
\item Choisir le sujet parmi les cinq proposés (ou celui imposé à l'examen).
\item Remplir la fiche des six questions : c'est votre plan.
\item Écrire le titre. Test de qualité : est-il informatif ? Donne-t-il envie de lire ?
\item Écrire le chapeau (\esp{copete}) en 25-30 mots. Chronométrez-vous : 5 minutes maximum.
\item Écrire le corps : deux paragraphes, avec une donnée chiffrée et une citation.
\item Écrire la conclusion : ouverture, conséquence future ou citation finale.
\item Compter les mots : vous devez être entre 110 et 130.
\item Relire avec la grille d'auto-évaluation proposée à la fin de cette section.
\end{enumerate}
\end{methode}

\courssub{Rédiger le titre : la phrase nominale percutante}

Le titre d'un article de presse espagnol n'est pas une phrase complète comme en français scolaire : c'est souvent une \cle{phrase nominale} (sans verbe conjugué) ou une phrase courte au présent, sans article quand c'est possible. Comparez :

\begin{center}
\begin{tabularx}{\linewidth}{|X|X|}
\hline
\rowcolor{popblueL} \textbf{Phrase complète (français)} & \textbf{Titre de presse (español)} \\
\hline
Les élèves ont organisé une marche pour la paix. & \esp{Estudiantes marchan por la paz} \\
\hline
La rentrée scolaire a commencé hier dans tout le pays. & \esp{Arranca el curso con nuevos retos} \\
\hline
De fortes pluies ont provoqué des inondations à Douala. & \esp{Fuertes lluvias inundan Douala} \\
\hline
Le festival de musique a réuni des milliers de jeunes. & \esp{Un festival reúne a miles de jóvenes} \\
\hline
\end{tabularx}
\end{center}

\begin{propriete}[Règles du titre journalistique]
\begin{itemize}
\item \textbf{Forme} : phrase nominale ou verbe au \cle{présent} de l'indicatif (le présent « historique » donne de l'actualité).
\item \textbf{Longueur} : 4 à 8 mots maximum.
\item \textbf{Suppression} : on omet souvent l'article défini : \esp{Estudiantes marchan...} et non \esp{Los estudiantes marchan...}.
\item \textbf{Interdit} : pas de point final, pas de première personne, pas de point d'interrogation sauf titre volontairement polémique.
\end{itemize}
\end{propriete}

\begin{exoresolu}[Transformer une phrase en titre]
Énoncé : transformez en titre de presse la phrase suivante : \esp{El Gobierno anunció ayer la construcción de cincuenta nuevas aulas en la región del Centro.}
\tcblower
Solution : on supprime l'article, on passe le verbe au présent, on garde l'information essentielle (qui, quoi, où) : \esp{El Gobierno anuncia cincuenta nuevas aulas en el Centro}. Version nominale encore plus courte : \esp{Cincuenta nuevas aulas para la región del Centro}. Les deux sont correctes ; la seconde est la plus « presse ». À éviter : garder le passé simple et l'adverbe \esp{ayer} dans le titre, ce qui l'alourdit inutilement.
\end{exoresolu}

\courssub{Rédiger le chapeau (copete) : 25-30 mots, les 4 W}

Le \esp{copete} (chapeau), appelé aussi \esp{entradilla} dans la presse espagnole, est le petit paragraphe en gras sous le titre. Sa contrainte est stricte : 25 à 30 mots, et il doit répondre aux quatre questions principales : \esp{quién, qué, dónde, cuándo} (le pourquoi et le comment sont réservés au corps).

\begin{exemplebox}[Du sujet au copete : un exemple complet]
Sujet : « Rentrée scolaire 2025 » $\rightarrow$ Titre : \esp{Arranca el curso con nuevos retos} $\rightarrow$ Copete : \esp{Más de tres millones de alumnos regresaron ayer a las aulas en todo Camerún tras las vacaciones, con nuevos libros y horarios.} (Plus de trois millions d'élèves sont retournés hier en classe dans tout le Cameroun après les vacances, avec de nouveaux livres et de nouveaux horaires.) — 21 mots, 4 W couvertes : quién (\esp{alumnos}), qué (\esp{regresaron a las aulas}), dónde (\esp{en todo Camerún}), cuándo (\esp{ayer}).
\end{exemplebox}

\begin{propriete}[La contrainte de longueur : 120 mots ± 10]
L'article complet doit faire \cle{120 mots}, avec une tolérance de 10 : visez donc entre 110 et 130 mots. Comment compter ? Méthode simple : comptez chaque mot, en considérant les mots-outils très courts (articles, prépositions, conjonctions : \esp{el, la, de, en, y, que}) comme négligeables — ou comptez tout et visez 125. Répartition conseillée : titre 5 mots, copete 28 mots, corps 60-70 mots, conclusion 15 mots.
\end{propriete}

\courssub{Développer le corps : données, citation, connecteurs}

Le corps se compose de \cle{deux paragraphes} de 3-4 lignes chacun. Premier paragraphe : les faits détaillés, avec au moins une \cle{donnée chiffrée} (un nombre, un pourcentage, une heure). Deuxième paragraphe : les réactions, avec une \cle{citation}. Remarquez que la presse hispanophone privilégie le \cle{passé simple} (\esp{comenzó, participaron, afirmó}) pour raconter l'événement : c'est le temps du récit journalistique.

\begin{definition}[La citation textuelle]
\esp{Cita textual} : reproduction exacte des paroles d'une source, entre guillemets, suivie du verbe de parole. Exemple : \esp{« La paz es tarea de todos », afirmó el director.} (« La paix est l'affaire de tous », a affirmé le directeur.) Verbes utiles : \esp{afirmar, declarar, explicar, añadir, señalar, comentar}.
\end{definition}

Les paragraphes doivent être reliés par des \cle{connecteurs}. Voici la boîte à outils indispensable :

\begin{center}
\begin{tabularx}{\linewidth}{|X|X|X|}
\hline
\rowcolor{popblueL} \textbf{Fonction} & \textbf{Español} & \textbf{Français} \\
\hline
Ajout & \esp{además, también, asimismo} & de plus, aussi, également \\
\hline
Cause & \esp{porque, debido a, a causa de} & parce que, en raison de \\
\hline
Conséquence & \esp{por eso, por lo tanto, así que} & c'est pourquoi, donc \\
\hline
Opposition & \esp{sin embargo, no obstante, pero} & cependant, néanmoins, mais \\
\hline
Temps & \esp{primero, después, finalmente} & d'abord, ensuite, enfin \\
\hline
Source & \esp{según, de acuerdo con} & selon, conformément à \\
\hline
\end{tabularx}
\end{center}

\begin{textebox}[Modèle de corps d'article : la Journée de la paix]
\esp{La marcha comenzó a las nueve de la mañana en el patio del Liceo Bilingüe de Yaoundé. Más de quinientos alumnos participaron con pancartas y cantos. Según los organizadores, también hubo talleres de diálogo y dibujos sobre la convivencia.}

\esp{« La paz es tarea de todos », afirmó el director del centro. Además, una ONG local ofreció cuadernos a los alumnos. Sin embargo, algunos padres pidieron más actividades durante el año.}
\end{textebox}

Glossaire : \esp{pancartas} = pancartes ; \esp{talleres} = ateliers ; \esp{convivencia} = vivre-ensemble ; \esp{centro} = établissement scolaire ; \esp{cuadernos} = cahiers ; \esp{hubo} = il y eut (passé simple de \esp{haber}). Comptez : ce corps fait environ 55 mots — exactement la cible.

\courssub{Rédiger la conclusion}

La conclusion d'un article de presse n'est pas un résumé : c'est une \cle{ouverture}. Trois formules possibles :
\begin{itemize}
\item la conséquence future : \esp{La próxima jornada tendrá lugar en diciembre.} (La prochaine journée aura lieu en décembre.)
\item l'ouverture vers un enjeu plus large : \esp{Un gesto pequeño que puede transformar todo un barrio.} (Un petit geste qui peut transformer tout un quartier.)
\item la citation finale : \esp{« Volveremos el año que viene », prometieron los alumnos.} (« Nous reviendrons l'année prochaine », ont promis les élèves.)
\end{itemize}

\begin{attention}[Faux amis et calques à éviter]
\begin{itemize}
\item \esp{actualmente} signifie « actuellement, en ce moment » ; pour dire « en fait », utilisez \esp{de hecho} ou \esp{en realidad}. De même, \esp{actual} signifie « actuel », et non « actuel » au sens de « réel » : « le problème réel » se dit \esp{el verdadero problema}.
\item Ne traduisez pas « environ 500 élèves » par un calque du français : dites \esp{unos quinientos alumnos}, \esp{cerca de quinientos alumnos} ou \esp{aproximadamente quinientos alumnos}.
\item « Avoir lieu » se dit \esp{tener lugar} : \esp{La fiesta tendrá lugar mañana.} Ne le confondez pas avec \esp{llevar a cabo} (« mener à bien, réaliser »), qui exige un complément : \esp{llevar a cabo un proyecto}.
\item Attention aux accents du passé simple : \esp{afirmó, participó, regresó, comenzó} portent tous un accent final à la 3\textsuperscript{e} personne du singulier ; sans accent, \esp{afirmo} est un présent (« j'affirme »).
\end{itemize}
\end{attention}

\courssub{Auto-évaluation : la grille du candidat}

Avant de rendre votre copie, vérifiez chaque point de cette grille. Chaque « oui » vaut des points ; chaque « non » est une erreur facile à corriger.

\begin{center}
\begin{tabularx}{\linewidth}{|X|l|}
\hline
\rowcolor{popblueL} \textbf{Critère} & \textbf{Oui / Non} \\
\hline
Mon article a les 4 parties : titre, copete, corps (2 §), conclusion & \\
\hline
Les 6 questions (quién, qué, dónde, cuándo, por qué, cómo) ont une réponse & \\
\hline
Le nombre de mots est entre 110 et 130 & \\
\hline
Registre objectif : aucun \esp{yo}, aucun \esp{creo que} & \\
\hline
Au moins 3 connecteurs différents & \\
\hline
Au moins 1 donnée chiffrée et 1 citation entre guillemets & \\
\hline
Orthographe : accents, \esp{ñ}, concordances des temps vérifiés & \\
\hline
\end{tabularx}
\end{center}

\begin{exercice}[À vous de jouer]
Choisissez l'un des cinq sujets de la liste du début de cette section (par exemple la CAN 2025 ou les inondations de Douala). Remplissez la fiche des six questions, puis rédigez l'article complet en 120 mots ± 10, en suivant les huit étapes de la méthode. Terminez par l'auto-évaluation.
\end{exercice}

\begin{corrige}[Exercice : éléments de correction]
Il n'existe pas une seule bonne réponse, mais tout article réussi contient : un titre nominal ou au présent (\esp{Los Leones Indomables preparan la CAN}), un copete de 25-30 mots couvrant les 4 W, un corps avec un chiffre (\esp{veintitrés jugadores convocados}) et une citation (\esp{« Estamos listos », declaró el entrenador.}), et une conclusion d'ouverture (\esp{El primer partido se jugará en enero.}). Vérifiez l'absence totale de \esp{yo} et la présence d'au moins trois connecteurs.
\end{corrige}

\begin{savaistu}[L'article de presse au Bac]
Au baccalauréat camerounais, l'article de presse de 120 mots tombe très régulièrement en épreuve d'expression écrite. Les sujets récurrents sont : l'environnement, la jeunesse, la technologie, la culture de la paix et l'égalité des genres. En Guinée équatoriale, seul pays hispanophone d'Afrique, la presse écrite en espagnol utilise exactement cette structure titre-chapeau-corps : un excellent modèle à imiter !
\end{savaistu}

Votre article est prêt. La section suivante vous fera découvrir, par l'observation, comment l'espagnol exprime la condition — un outil précieux pour nuancer vos phrases, y compris dans vos conclusions d'articles.

\coursec{Observation et découverte : la condition en espagnol (si-clauses)}

Après avoir rédigé votre article de presse dans la section précédente, nous changeons de registre pour nous concentrer sur un point de grammaire fondamental, omniprésent dans la presse, la littérature et l'oral : \cle{la phrase conditionnelle (si-clause)}. Que vous commentiez un éditorial, traduisiez un extrait de roman ou argumentiez dans un débat, maîtriser les trois types de condition (réalisable, irréalisable présent, irréalisable passé) est indispensable pour obtenir la note maximale au Bac.

\courssub{Corpus d'observation : six phrases-modèles}

Commençons par observer six phrases authentiques. Lisez-les attentivement, repérez le mot \esp{si}, identifiez le temps du verbe qui suit immédiatement \esp{si} (dans la proposition subordonnée) et le temps du verbe principal (dans la proposition principale). Classez-les mentalement en trois groupes selon le sens : \textbf{réel/possible}, \textbf{imaginaire présent}, \textbf{imaginaire passé}.

\begin{textebox}[Corpus d'observation : les 6 phrases-modèles]
1. \esp{Si estudias, aprobarás.} \\
2. \esp{Si estudiara, aprobaría.} \\
3. \esp{Si hubiera estudiado, habría aprobado.} \\
4. \esp{Si vienes mañana, te veo.} \\
5. \esp{Si vinieras mañana, te vería.} \\
6. \esp{Si hubieras venido ayer, te habría visto.}
\end{textebox}

\begin{exoresolu}[Analyse guidée des phrases 1 et 4 (Type 1 : Réalisable)]
\textbf{Énoncé :} Analysez les phrases 1 et 4. Quel temps suit \esp{si} ? Quel temps est à la principale ? Quel est le sens ?
\tcblower
\textbf{Solution détaillée :}
\begin{itemize}
    \item \textbf{Phrase 1 :} \esp{Si estudias (présent indicatif), aprobarás (futur simple indicatif).} Sens : condition \textbf{réelle et possible} dans le futur. « Si tu étudies (réalité possible), tu réussiras. »
    \item \textbf{Phrase 4 :} \esp{Si vienes (présent indicatif), te veo (présent indicatif).} Sens : condition \textbf{réelle et possible} (présent/futur proche). « Si tu viens demain, je te vois. » Le présent à la principale exprime un futur immédiat ou une certitude.
\end{itemize}
\textbf{Constatation :} Dans les deux cas, \esp{si} est suivi de l'\textbf{indicatif présent}. Jamais du futur.
\end{exoresolu}

\begin{exoresolu}[Analyse guidée des phrases 2 et 5 (Type 2 : Irréalisable présent)]
\textbf{Énoncé :} Analysez les phrases 2 et 5. Quel temps suit \esp{si} ? Quel temps est à la principale ? Quel est le sens ?
\tcblower
\textbf{Solution détaillée :}
\begin{itemize}
    \item \textbf{Phrase 2 :} \esp{Si estudiara (subjonctif imparfait), aprobaría (conditionnel simple).} Sens : condition \textbf{irréelle au présent}. Le sujet n'étudie pas (ou pas assez), donc il ne réussira pas. C'est une hypothèse contraire à la réalité actuelle.
    \item \textbf{Phrase 5 :} \esp{Si vinieras (subjonctif imparfait), te vería (conditionnel simple).} Sens : condition \textbf{irréelle au présent/futur}. « Si tu venais (mais tu ne viens pas / c'est improbable), je te verrais. »
\end{itemize}
\textbf{Constatation :} \esp{si} + \textbf{subjonctif imparfait} (souvent en \esp{-ra} : \esp{estudiara, viniera, tuviera, fuera}) $\rightarrow$ principale au \textbf{conditionnel simple}.
\end{exoresolu}

\begin{exoresolu}[Analyse guidée des phrases 3 et 6 (Type 3 : Irréalisable passé)]
\textbf{Énoncé :} Analysez les phrases 3 et 6. Quel temps suit \esp{si} ? Quel temps est à la principale ? Quel est le sens ?
\tcblower
\textbf{Solution détaillée :}
\begin{itemize}
    \item \textbf{Phrase 3 :} \esp{Si hubiera estudiado (subjonctif plus-que-parfait), habría aprobado (conditionnel composé).} Sens : condition \textbf{irréelle au passé}. Le sujet n'a pas étudié (fait passé accompli), il n'a donc pas réussi. On imagine un passé alternatif.
    \item \textbf{Phrase 6 :} \esp{Si hubieras venido (subjonctif plus-que-parfait), te habría visto (conditionnel composé).} Sens : condition \textbf{irréelle au passé}. « Si tu étais venu (hier, mais tu n'es pas venu), je t'aurais vu. »
\end{itemize}
\textbf{Constatation :} \esp{si} + \textbf{subjonctif plus-que-parfait} (\esp{hubiera/hubieras/hubiera/hubiéramos/hubierais/hubieran} + participe passé) $\rightarrow$ principale au \textbf{conditionnel composé} (\esp{habría/habrías/habría/habríamos/habríais/habrían} + participe passé).
\end{exoresolu}

\courssub{Classification en trois types (Typologie officielle)}

Sur la base de l'observation, nous classons les phrases conditionnelles (\cle{si-clauses}) en trois types canoniques. Ce classement est la clé de voûte de la traduction thème/version au Bac.

\begin{popfigure}
\begin{tikzpicture}[node distance=1.2cm and 0.5cm, every node/.style={align=center, font=\small}]
\node[draw, rounded corners, fill=popblueL, thick] (root) {\textbf{SI-CLAUSES : 3 TIPOS}};
\node[draw, rounded corners, fill=popgreenL, below left=of root, xshift=-1cm, align=center] (t1){\textbf{TIPO 1: REALIZABLE}\\\textit{Condition réelle / possible}\\\esp{Si + PRES. IND. $\rightarrow$ FUT. IND. (ou PRES. IND.)}\\\esp{Si llueve, no iremos.}};
\node[draw, rounded corners, fill=popgoldL, below=of root, align=center] (t2){\textbf{TIPO 2: IRREAL PRESENTE}\\\textit{Condition hypothétique / irréelle maintenant}\\\esp{Si + SUBJ. IMPERF. $\rightarrow$ COND. SIMPLE}\\\esp{Si lloviera, no iríamos.}};
\node[draw, rounded corners, fill=poppinkL, below right=of root, xshift=1cm, align=center] (t3){\textbf{TIPO 3: IRREAL PASADO}\\\textit{Condition irréelle dans le passé}\\\esp{Si + SUBJ. PLUSQ. $\rightarrow$ COND. COMPOSÉ}\\\esp{Si hubiera llovido, no habríamos ido.}};
\draw[->, thick, popblue] (root.south west) -- (t1.north east);
\draw[->, thick, poporange] (root.south) -- (t2.north);
\draw[->, thick, poppink] (root.south east) -- (t3.north west);
\end{tikzpicture}
\legende{Cartographie mentale des trois types de conditionnels (Si-clauses).}
\end{popfigure}

\courssub{Terminologie grammaticale : Protase et Apodose}

Pour parler comme un linguiste (et réussir l'oral ou le commentaire de texte), il faut nommer les deux parties de la phrase conditionnelle.

\begin{definition}[Protase et Apodose]
\begin{itemize}
    \item \textbf{La Protase} (du grec \emph{protaxis}, « première proposition ») : c'est la \textbf{proposition subordonnée conditionnelle introduite par \esp{si}}. Elle exprime la condition. \textit{Exemple : \esp{Si estudias...}}
    \item \textbf{L'Apodose} (du grec \emph{apodosis}, « donnée en retour ») : c'est la \textbf{proposition principale}, celle qui exprime la conséquence, le résultat. \textit{Exemple : \esp{...aprobarás.}}
\end{itemize}
L'ordre peut s'inverser sans changer la grammaire : \esp{Aprobarás si estudias.} (Apodose + Protase). La virgule disparaît souvent quand la protase est en seconde position.
\end{definition}

\courssub{La Règle d'Or : Ce qui est INTERDIT après \esp{si}}

C'est le point unique qui fait perdre le plus de points aux francophones. En français, on dit « \emph{Si j'irai...} » (futur) ou « \emph{Si j'aurais su...} » (conditionnel). \textbf{En espagnol, c'est strictement impossible.}

\begin{propriete}[RÈGLE D'OR : Interdits absolus après \esp{si}]
Après la conjonction \esp{si} (introduisant une condition), on ne trouve \textbf{JAMAIS} :
\begin{enumerate}
    \item L'\textbf{Indicatif Futur} (\esp{iré, tendrás, vendrá, habrá}...). $\rightarrow$ \textbf{Incorrect :} \esp{*Si yo iré...}
    \item Le \textbf{Conditionnel} (simple ou composé : \esp{iría, habría, habría ido}...). $\rightarrow$ \textbf{Incorrect :} \esp{*Si yo tendría...}, \esp{*Si yo habría sabido...}
\end{enumerate}
\textbf{Seuls trois temps sont autorisés après \esp{si} :}
\begin{enumerate}
    \item \textbf{Indicatif Présent} (Type 1 : Réel).
    \item \textbf{Subjonctif Imparfait} (Type 2 : Irréel présent). Formes en \esp{-ra} (\esp{cantara, comiera, viviera, fuera, tuviera}) ou \esp{-se} (\esp{cantase, comiese, viviese, fuese, tuviese}) — les deux sont corrects, \esp{-ra} est plus fréquent à l'oral et dans la presse.
    \item \textbf{Subjonctif Plus-que-parfait} (Type 3 : Irréel passé). \esp{hubiera/hubiese + participe passé}.
\end{enumerate}
\end{propriete}

\courssub{Tableau récapitulatif comparatif Français / Espagnol}

Le piège principal vient de l'interférence du français. Ce tableau doit être mémorisé par cœur.

\begin{center}
\begin{tabularx}{\linewidth}{|X|X|X|X|}
\hline
\rowcolor{popblueL}
\textbf{TYPE} & \textbf{FRANÇAIS (Structure)} & \textbf{ESPAGNOL (Structure)} & \textbf{EXEMPLE (Es / Fr)} \\
\hline
\textbf{1. Réalisable} & \emph{Si} + \textbf{Présent} $\rightarrow$ \textbf{Futur} (ou Présent) & \esp{Si} + \textbf{Présent Ind.} $\rightarrow$ \textbf{Futur Ind.} (ou Présent Ind.) & \esp{Si tienes tiempo, vendrás.} / « Si tu as le temps, tu viendras. » \\
\hline
\textbf{2. Irréel Présent} & \emph{Si} + \textbf{Imparfait} $\rightarrow$ \textbf{Conditionnel Présent} & \esp{Si} + \textbf{Subj. Imparfait} $\rightarrow$ \textbf{Conditionnel Simple} & \esp{Si tuvieras tiempo, vendrías.} / « Si tu avais le temps, tu viendrais. » \\
\hline
\textbf{3. Irréel Passé} & \emph{Si} + \textbf{Plus-que-parfait} $\rightarrow$ \textbf{Conditionnel Passé} & \esp{Si} + \textbf{Subj. Plus-que-parfait} $\rightarrow$ \textbf{Conditionnel Composé} & \esp{Si hubieras tenido tiempo, habrías venido.} / « Si tu avais eu le temps, tu serais venu. » \\
\hline
\end{tabularx}
\end{center}

\courssub{Exemples traduits variés (Thème / Version)}

Entraînez-vous à passer de l'un à l'autre instantanément. Notez l'usage du \esp{voseo} (Amérique) ou \esp{tuteo} (Espagne) selon le contexte ; ici nous utilisons la norme péninsulaire standard (\esp{tú/usted}) exigée au Bac camerounais.

\begin{exemplebox}[Série Type 1 : Réalisable (Probabilité forte / Certitude future)]
\begin{itemize}
    \item \esp{Si llueve mañana, no iremos a la playa.} « S'il pleut demain, nous n'irons pas à la plage. »
    \item \esp{Si terminas tus deberes, puedes jugar.} « Si tu finis tes devoirs, tu peux jouer. » (Présent à la principale = permission/avenir proche).
    \item \esp{Si el precio baja, lo compraré.} « Si le prix baisse, je l'achèterai. »
    \item \esp{Si no estudias, suspenderás el examen.} « Si tu n'étudies pas, tu rateras l'examen. » (Avertissement).
\end{itemize}
\end{exemplebox}

\begin{exemplebox}[Série Type 2 : Irréel Présent (Hypothèse contraire à la réalité actuelle)]
\begin{itemize}
    \item \esp{Si yo fuera tú, no lo haría.} « Si j'étais toi, je ne le ferais pas. » (\esp{ser} $\rightarrow$ \esp{fuera} ; \emph{verbe irrégulier clé}).
    \item \esp{Si tuviera dinero, compraría ese coche.} « Si j'avais de l'argent, j'achèterais cette voiture. » (\esp{tener} $\rightarrow$ \esp{tuviera}).
    \item \esp{Si supiera la verdad, te la diría.} « Si je savais la vérité, je te la dirais. » (\esp{saber} $\rightarrow$ \esp{supiera}).
    \item \esp{Si lloviera ahora, nos mojaríamos.} « S'il pleuvait maintenant, nous nous mouillerions. »
\end{itemize}
\end{exemplebox}

\begin{exemplebox}[Série Type 3 : Irréel Passé (Regret, reproche, passé alternatif)]
\begin{itemize}
    \item \esp{Si hubieras estudiado, habrías aprobado.} « Si tu avais étudié, tu aurais réussi. »
    \item \esp{Si no hubiera llovido, habríamos ido al campo.} « S'il n'avait pas plu, nous serions allés à la campagne. »
    \item \esp{Si me hubieras avisado, habría venido.} « Si tu m'avais prévenu, je serais venu. »
    \item \esp{Si el equipo hubiera marcado, habrían ganado.} « Si l'équipe avait marqué, ils auraient gagné. »
\end{itemize}
\end{exemplebox}

\courssub{Pièges du Bac : Erreurs fréquentes des francophones (Boîte \texttt{attention})}

Ces trois erreurs représentent 80\% des fautes de conditionnel en thème/version. Ancrez les corrections dans votre mémoire à long terme.

\begin{attention}[Top 3 des erreurs fatales après \esp{si}]
\textbf{Erreur n°1 : Le « Futur français » calqué.}
\begin{itemize}
    \item \textcolor{red}{$\times$ \esp{Si yo iré a Yaoundé, te veré.}} (Calque de « Si j'irai... »)
    \item \textcolor{popgreen}{$\checkmark$ \esp{Si voy a Yaoundé, te veré.}} (Indicatif présent \esp{voy} après \esp{si}).
\end{itemize}

\textbf{Erreur n°2 : Le « Conditionnel français » calqué (Type 2).}
\begin{itemize}
    \item \textcolor{red}{$\times$ \esp{Si yo tendría dinero, compraría la casa.}} (Calque de « Si j'aurais... »)
    \item \textcolor{popgreen}{$\checkmark$ \esp{Si tuviera dinero, compraría la casa.}} (Subjonctif imparfait \esp{tuviera} après \esp{si}).
\end{itemize}

\textbf{Erreur n°3 : Le « Conditionnel passé français » calqué (Type 3).}
\begin{itemize}
    \item \textcolor{red}{$\times$ \esp{Si yo habría sabido, habría venido.}} (Calque de « Si j'aurais su... »)
    \item \textcolor{popgreen}{$\checkmark$ \esp{Si hubiera sabido, habría venido.}} (Subjonctif plus-que-parfait \esp{hubiera sabido} après \esp{si}).
\end{itemize}

\textbf{Autre piège :} \esp{Si} + \textbf{Indicatif Imparfait} n'existe \textbf{PAS} en espagnol standard pour la condition. \esp{*Si estudiaba...} (Incorrect pour une hypothèse). On dit \esp{Si estudiara...} (Subj. Imperf.).
\end{attention}

\courssub{Nuances avancées : \esp{Si} + Présent + Présent (Loi générale) et la concession (\esp{Aunque})}

Deux cas particuliers enrichissent votre expression écrite (article de presse, essai) :

\begin{propriete}[Cas particuliers utiles pour la rédaction]
\begin{enumerate}
    \item \textbf{Loi générale / Vérité scientifique (Type 1 variant) :} \esp{Si} + \textbf{Présent Indicatif} $\rightarrow$ \textbf{Présent Indicatif}.
    \textit{Exemple :} \esp{Si calientas el agua a 100 °C, hierve.} « Si tu chauffes l'eau à 100 °C, elle bout. » (Pas de futur, c'est une loi physique).
    \item \textbf{La concession ( = « Bien que », « Même si ») :} En espagnol standard, on utilise \textbf{\esp{Aunque}} (et non \esp{Si}) pour exprimer la concession.
    \begin{itemize}
        \item \textbf{Fait réel (concession réelle) :} \esp{Aunque} + \textbf{Indicatif Présent}.
        \textit{Exemple :} \esp{Aunque es caro, lo compraré.} « \textbf{Bien qu'il soit} cher (réel), je l'achèterai. »
        \item \textbf{Hypothèse (concession hypothétique) :} \esp{Aunque} + \textbf{Subjonctif Présent}.
        \textit{Exemple :} \esp{Aunque sea caro, lo compraré.} « \textbf{Même s'il était} cher (hypothétique), je l'achèterai. »
    \end{itemize}
    \textit{Attention :} L'usage de \esp{Si} pour « même si » (\esp{Si es caro...}) est possible dans la langue courante pour une concession réelle (indicatif), mais \esp{Aunque} est la norme cultive attendue au Bac.
\end{enumerate}
\end{propriete}

\courssub{Entraînement résolu : Thème / Version mixte}

Appliquons la règle complète sur des phrases de niveau Bac (sujets de société, Cameroun, Espagne).

\begin{exoresolu}[Thème : Traduisez en espagnol]
\textbf{Énoncé :} Traduisez les phrases suivantes en respectant la typologie des si-clauses.
\begin{enumerate}
    \item Si le gouvernement investit dans l'éducation, le pays se développera. (Réalisable)
    \item Si j'étais ministre, je construirais plus d'écoles. (Irréel présent)
    \item Si les jeunes avaient voté, les résultats auraient changé. (Irréel passé)
    \item Bien qu'il pleuve (fait réel), le match aura lieu. (Concession réelle)
    \item Si tu avais étudié l'espagnol plus tôt, tu le parlerais couramment maintenant. (Mixte : Passé $\rightarrow$ Présent)
\end{enumerate}
\tcblower
\textbf{Solution détaillée :}
\begin{enumerate}
    \item \esp{Si el gobierno invierte en la educación, el país se desarrollará.} (Type 1 : \esp{invierte} Pres. Ind. $\rightarrow$ \esp{se desarrollará} Fut. Ind.). \textit{Note : \esp{invertir en} + nom.}
    \item \esp{Si yo fuera ministro, construiría más escuelas.} (Type 2 : \esp{fuera} Subj. Imperf. \emph{ser} $\rightarrow$ \esp{construiría} Cond. Simple). \textit{Attention : \esp{ministro} sans article après \esp{ser} pour la profession.}
    \item \esp{Si los jóvenes hubieran votado, los resultados habrían cambiado.} (Type 3 : \esp{hubieran votado} Subj. Plusq. $\rightarrow$ \esp{habrían cambiado} Cond. Composé).
    \item \esp{Aunque llueve, el partido tendrá lugar.} (Concession réelle : \esp{Aunque} + \textbf{Indicatif Présent} \esp{llueve} $\rightarrow$ Futur \esp{tendrá}). \textit{Variante possible : \esp{Si llueve...} (Indicatif), mais \esp{Aunque} est préférable.}
    \item \esp{Si hubieras estudiado español antes, lo hablarías con fluidez ahora.} (Phrase \textbf{mixte} : Protase Type 3 \esp{hubieras estudiado} / Apodose Type 2 \esp{hablarías}. Le conditionnel simple \esp{hablarías} exprime le résultat \textbf{actuel} d'une cause \textbf{passée}).
\end{enumerate}
\end{exoresolu}

\courssub{Le saviez-vous ? : L'influence de l'anglais et l'oral familier}

\begin{savaistu}[L'anglicisme \esp{*Si habría...*} : une peste moderne]
Dans l'espagnol très influencé par l'anglais (surtout aux USA, Porto Rico, et de plus en plus sur les réseaux sociaux), on entend souvent la structure calquée sur \emph{« If I would have..., I would have... »} :
\begin{itemize}
    \item \textcolor{red}{Incorrect :} \esp{*Si yo habría sabido, habría ido.*} (Double conditionnel).
    \item \textcolor{red}{Incorrect :} \esp{*Si lloviera, no habría ido.*} (Mélange Subj. Imperf / Cond. Composé).
\end{itemize}
\textbf{Rappel strict pour le Bac (écrit et oral) :} La norme cultive panhispanique (RAE) exige \textbf{impérativement} le subjonctif après \esp{si} pour l'irréel. \esp{Si hubiera sabido...} / \esp{Si lloviera...}. Ne vous laissez pas contaminer par les séries Netflix doublées ou les tweets ! La « \cle{Règle d'Or} » (boîte Propriété 4.4) est votre seule boussole.
\end{savaistu}

\courssub{Exercices d'application (À faire sur cahier)}

\begin{exercice}[Classification et Conjugaison]
Mettez le verbe entre parenthèses au temps convenable (Ind. Présent, Subj. Imperf, Subj. Plusq., Futur, Cond. Simple, Cond. Composé). Identifiez le Type (1, 2 ou 3).
\begin{enumerate}
    \item Si yo \_\_\_\_\_\_\_\_ (tener) tiempo, \_\_\_\_\_\_\_\_ (leer) ese libro. (Contexte : J'ai du temps libre cet après-midi).
    \item Si yo \_\_\_\_\_\_\_\_ (tener) tiempo, \_\_\_\_\_\_\_\_ (leer) ese libro. (Contexte : Je suis débordé, c'est impossible aujourd'hui).
    \item Si ayer \_\_\_\_\_\_\_\_ (hacer) buen tiempo, \_\_\_\_\_\_\_\_ (salir) al parque.
    \item Si \_\_\_\_\_\_\_\_ (tú / estudiar) más, \_\_\_\_\_\_\_\_ (aprobar) el examen de selectividad.
    \item \_\_\_\_\_\_\_\_ (Aunque) \_\_\_\_\_\_\_\_ (ser) difícil, lo intentaré. (Contexte : C'est difficile, c'est un fait).
\end{enumerate}
\end{exercice}

\begin{corrige}[Exercice 4.11]
\begin{enumerate}
    \item \esp{Si yo \textbf{tengo} tiempo, \textbf{leeré} ese libro.} (Type 1 : Réel possible). \textit{Ind. Pres. $\rightarrow$ Fut.}
    \item \esp{Si yo \textbf{tuviera} tiempo, \textbf{leería} ese libro.} (Type 2 : Irréel présent). \textit{Subj. Imperf. \esp{tuviera} $\rightarrow$ Cond. Simple.}
    \item \esp{Si ayer \textbf{hubiera hecho} buen tiempo, \textbf{habríamos salido} al parque.} (Type 3 : Irréel passé). \textit{Subj. Plusq. \esp{hubiera hecho} $\rightarrow$ Cond. Composé \esp{habríamos salido}.}
    \item \esp{Si \textbf{estudiaras} más, \textbf{aprobarías} el examen.} (Type 2 : Conseil hypothétique / Irréel présent). \textit{Subj. Imperf. \esp{estudiaras} $\rightarrow$ Cond. Simple.}
    \item \esp{\textbf{Aunque} \textbf{es} difícil, lo intentaré.} (Concession réelle : \esp{Aunque} + \textbf{Indicatif Présent} \esp{es} $\rightarrow$ Indicatif Futur \esp{intentaré}). \textit{Si c'était une hypothèse (\emph{« Même s'il était difficile »}), on écrirait : \esp{Aunque sea difícil...} (Subj. Présent).}
\end{enumerate}
\end{corrige}

\begin{methode}[Comment ne JAMAIS se tromper sur un si-clause (Check-list Bac)]
\begin{enumerate}
    \item \textbf{Identifiez le sens :} Est-ce \textbf{réel/possible} (Type 1) ? \textbf{Imaginaire maintenant} (Type 2) ? \textbf{Imaginaire dans le passé} (Type 3) ?
    \item \textbf{Choisissez le temps de la PROTASE (après \esp{si}) :}
    \begin{itemize}
        \item Type 1 $\rightarrow$ \textbf{Indicatif Présent}.
        \item Type 2 $\rightarrow$ \textbf{Subjonctif Imparfait} (\esp{-ra} ou \esp{-se}).
        \item Type 3 $\rightarrow$ \textbf{Subjonctif Plus-que-parfait} (\esp{hubiera} + Part. Passé).
    \end{itemize}
    \item \textbf{Choisissez le temps de l'APODOSE (principale) :}
    \begin{itemize}
        \item Type 1 $\rightarrow$ \textbf{Futur} (ou Présent).
        \item Type 2 $\rightarrow$ \textbf{Conditionnel Simple}.
        \item Type 3 $\rightarrow$ \textbf{Conditionnel Composé} (\esp{habría} + Part. Passé).
    \end{itemize}
    \item \textbf{Vérifiez l'interdiction :} \esp{Si} + Futur ? \esp{Si} + Conditionnel ? $\rightarrow$ \textbf{FOUTE !} Recommencez à l'étape 1.
    \item \textbf{Vérifiez les irréguliers du Subjonctif Imparfait (formes en \esp{-ra}) :}
    \begin{center}
    \begin{tabular}{|l|l|l|l|l|}
    \hline
    \rowcolor{popblueL}
    \textbf{Infinitif} & \textbf{1ère pers. sg.} & \textbf{Infinitif} & \textbf{1ère pers. sg.} & \textbf{Infinitif} \\
    \hline
    \esp{ser / ir} & \esp{fuera} & \esp{haber} & \esp{hubiera} & \esp{saber} \\
    \esp{tener} & \esp{tuviera} & \esp{poder} & \esp{pudiera} & \esp{querer} \\
    \esp{venir} & \esp{viniera} & \esp{poner} & \esp{pusiera} & \esp{decir} \\
    \esp{traer} & \esp{trajera} & \esp{caber} & \esp{cupiera} & \esp{haber (aux)} \\
    \hline
    \end{tabular}
    \end{center}
    \textit{Pour le Subj. Plus-que-parfait : ajouter \esp{hubiera} + participe passé (\esp{hubiera sido, hubiera tenido, hubiera sabido...}).}
\end{enumerate}
\end{methode}

\begin{aretenir}[Transition vers la suite]
Cette section pose les fondations grammaticales incontournables. La section 5 « Règle complète, tableaux de formes et comparaison FR/ES » consolidera ces acquis par des tableaux de conjugaison complets des verbes irréguliers au subjonctif imparfait et plus-que-parfait, et la section 6 proposera un entraînement intensif thème/version sur corpus de presse.
\end{aretenir}

\coursec{Règle complète, tableaux de formes et comparaison FR/ES}

Après avoir observé, en section 4, le fonctionnement des phrases conditionnelles à travers des exemples authentiques tirés de la presse hispanophone, nous allons maintenant fixer les connaissances sous forme de règles rigoureuses, de tableaux de conjugaison et de comparaisons systématiques avec le français. C'est l'étape indispensable pour réussir les exercices de thème/version du Baccalauréat et pour rédiger des articles de presse où l'hypothèse et la projection sont fréquentes (\emph{« Si el gobierno invierte más, la educación mejorará »}).

\courssub{Les trois types de phrases conditionnelles : tableau de synthèse}

En espagnol, comme en français, une phrase conditionnelle comporte deux propositions : la \cle{protase} (introduite par \esp{si}, condition) et l'\cle{apodose} (conséquence). La combinaison des temps de ces deux propositions définit trois types distincts, que le tableau suivant résume de manière exhaustive.

\begin{popfigure}
\begin{tikzpicture}[scale=0.75, every node/.style={font=\small}]
\node[draw=popblue, thick, fill=popblueL, rounded corners=8pt, inner sep=10pt, align=center] (tab){
\begin{tabular}{|>{\centering\arraybackslash}m{2.2cm}|>{\centering\arraybackslash}m{2.5cm}|>{\centering\arraybackslash}m{2.5cm}|>{\centering\arraybackslash}m{2.8cm}|>{\centering\arraybackslash}m{3.5cm}|}
\hline
\rowcolor{popblue}
\textcolor{white}{\textbf{TIPO}} & \textcolor{white}{\textbf{SI (PROTASE)}} & \textcolor{white}{\textbf{APODOSE}} & \textcolor{white}{\textbf{VALEUR}} & \textcolor{white}{\textbf{EJEMPLO}} \\
\hline
\textbf{1. Realizable} & Présent de l'indicatif & Futur simple / Présent indicatif / Impératif & Possible, probable, future & \esp{Si estudias, aprobarás.} \\
\hline
\textbf{2. Irreal Presente} & Subjonctif imparfait & Conditionnel simple & Hypothétique, peu probable, présent/futur & \esp{Si estudiara, aprobaría.} \\
\hline
\textbf{3. Irreal Pasado} & Subjonctif plus-que-parfait & Conditionnel composé & Contre-factuel, regret, passé & \esp{Si hubieras estudiado, habrías aprobado.} \\
\hline
\end{tabular}
};
\node[below=0.2cm of tab, text=popdark, font=\small\itshape] {Tableau de synthèse des trois types de \emph{si-clauses} en espagnol.};
\end{tikzpicture}
\legende{Les trois types de conditionnels en espagnol : forme, valeur et exemple type.}
\end{popfigure}

\begin{definition}[Valeurs des trois types]
\begin{itemize}
    \item \textbf{Type 1 (Réalisable) :} La condition est envisagée comme \textbf{réelle et possible} au moment de l'énonciation. L'apodose exprime une conséquence certaine ou une prévision. \esp{Si mañana llueve, no iremos al campo.} « S'il pleut demain, nous n'irons pas à la campagne. »
    \item \textbf{Type 2 (Irréel présent) :} La condition est \textbf{imaginaire, hypothétique ou contraire à la réalité présente}. On exprime un souhait, un conseil ou une conjecture. \esp{Si yo fuera tú, aceptaría la oferta.} « Si j'étais toi, j'accepterais l'offre. »
    \item \textbf{Type 3 (Irréel passé) :} La condition \textbf{n'a pas été réalisée dans le passé} (contre-factuel passé). On exprime un regret ou une critique rétrospective. \esp{Si hubiéramos sabido la verdad, no habríamos firmado.} « Si nous avions su la vérité, nous n'aurions pas signé. »
\end{itemize}
\end{definition}

\courssub{Formation des temps nécessaires : subjonctif et conditionnel}

Pour construire correctement les types 2 et 3, il faut maîtriser parfaitement la formation de quatre temps : le subjonctif imparfait (deux séries), le subjonctif plus-que-parfait, le conditionnel simple et le conditionnel composé.

\textbf{Subjonctif imparfait : les deux séries (\texttt{-ra} / \texttt{-se})}\par

Le radical du subjonctif imparfait se construit sur la \textbf{3\textsuperscript{e} personne du pluriel du Prétérit Indéfini (PIS)}. On supprime la finale \texttt{-ron} et on ajoute les terminaisons.

\begin{propriete}[Formation du subjonctif imparfait : deux séries correctes]
\textbf{Radical :} 3\textsuperscript{e} pers. plur. PIS sans \texttt{-ron} (\esp{hablaron} $\rightarrow$ \esp{habl-}, \esp{comieron} $\rightarrow$ \esp{comier-}, \esp{vivieron} $\rightarrow$ \esp{vivier-}).

\textbf{Série 1 (en \texttt{-ra}) :} \texttt{-ra, -ras, -ra, -ramos, -rais, -ran} \\
\textbf{Série 2 (en \texttt{-se}) :} \texttt{-se, -ses, -se, -semos, -seis, -sen}

\begin{center}
\begin{tabular}{|l|c|c|c|}
\hline
\rowcolor{poporangeL}
\textbf{Verbe} & \textbf{Radical} & \textbf{Série \texttt{-ra} (Amérique / usage courant)} & \textbf{Série \texttt{-se} (Espagne / littéraire)} \\
\hline
\esp{Hablar} & \esp{habl-} & \esp{hablara, hablaras, hablara, habláramos, hablarais, hablaran} & \esp{hablase, hablases, hablase, hablásemos, hablaseis, hablasen} \\
\hline
\esp{Comer} & \esp{comier-} & \esp{comiera, comieras, comiera, comiéramos, comierais, comieran} & \esp{comiese, comieses, comiese, comiésemos, comieseis, comiesen} \\
\hline
\esp{Vivir} & \esp{vivier-} & \esp{viviera, vivieras, viviera, viviéramos, vivierais, vivieran} & \esp{viviese, vivieses, viviese, viviésemos, vivieseis, viviesen} \\
\hline
\end{tabular}
\end{center}

\emph{Remarque fondamentale :} \textbf{Les deux séries sont correctes et interchangeables} dans la norme cultuelle. Au Baccalauréat camerounais, les deux sont acceptées. La série en \texttt{-ra} prédomine en Amérique latine et à l'oral ; la série en \texttt{-se} est plus fréquente en Espagne péninsulaire et à l'écrit soigné.
\end{propriete}

\begin{attention}[Accents graphiques au subjonctif imparfait]
Ne jamais oublier l'accent sur la \textbf{1\textsuperscript{re} personne du pluriel} (\esp{habláramos / hablásemos, comiéramos / comiésemos, viviéramos / viviésemos}) et sur la \textbf{2\textsuperscript{e} personne du pluriel} de la série \texttt{-se} (\esp{hablaseis, comieseis, vivieseis}). L'absence d'accent change le temps (ex: \esp{hablaramos} = présent indicatif / passé composé).
\end{attention}

\textbf{Subjonctif plus-que-parfait}\par

Il se forme avec l'auxiliaire \esp{haber} au subjonctif imparfait + le \textbf{participe passé} du verbe principal. Les deux séries de l'imparfait s'appliquent à l'auxiliaire.

\begin{propriete}[Subjonctif plus-que-parfait]
\textbf{Formule :} \esp{hubiera / hubiese} + \textbf{participio} (\texttt{-ado, -ido}).

\begin{center}
\begin{tabular}{|l|c|c|}
\hline
\rowcolor{popgreenL}
\textbf{Sujet} & \textbf{Série \texttt{-ra}} & \textbf{Série \texttt{-se}} \\
\hline
\esp{Yo} & \esp{hubiera hablado / comido / vivido} & \esp{hubiese hablado / comido / vivido} \\
\hline
\esp{Tú} & \esp{hubieras hablado...} & \esp{hubieses hablado...} \\
\hline
\esp{Él/Ella/Ud.} & \esp{hubiera hablado...} & \esp{hubiese hablado...} \\
\hline
\esp{Nosotros} & \esp{hubiéramos hablado...} & \esp{hubiésemos hablado...} \\
\hline
\emph{Vosotros} & \esp{hubierais hablado...} & \esp{hubieseis hablado...} \\
\hline
\esp{Ellos/Uds.} & \esp{hubieran hablado...} & \esp{hubiesen hablado...} \\
\hline
\end{tabular}
\end{center}
\end{propriete}

\textbf{Conditionnel simple : réguliers et irréguliers}\par

Le conditionnel simple s'obtient en ajoutant les terminaisons de l'imparfait de l'indicatif (\texttt{-ía, -ías, -ía, -íamos, -íais, -ían}) à l'\textbf{infinitif} (pour les réguliers) ou au \textbf{radical irrégulier} (identique au futur simple).

\begin{propriete}[Verbes irréguliers au conditionnel (mêmes radicaux qu'au futur)]
Ces 9 verbes « radicaux » sont \textbf{impératifs} à connaître par cœur pour le Bac. Ils ne prennent \textbf{jamais} de \texttt{d} supplémentaire (erreurs fréquentes : *\esp{podría} $\neq$ *\esp{podrría}, *\esp{querría} $\neq$ *\esp{querería}).

\begin{center}
\begin{tabular}{|l|l|l|l|}
\hline
\rowcolor{poppurpleL}
\textbf{Infinitif} & \textbf{Radical irrégulier} & \textbf{1\textsuperscript{re} pers. sg (Yo)} & \textbf{Traduction (Je...)} \\
\hline
\esp{Tener} & \esp{tendr-} & \esp{Tendría} & aurais \\
\hline
\esp{Poner} & \esp{pondr-} & \esp{Pondría} & mettrais \\
\hline
\esp{Salir} & \esp{saldr-} & \esp{Saldría} & sortirais \\
\hline
\esp{Venir} & \esp{vendr-} & \esp{Vendría} & viendrais \\
\hline
\esp{Hacer} & \esp{har-} & \esp{Haría} & ferais \\
\hline
\esp{Querer} & \esp{querr-} & \esp{Querría} & voudrais \\
\hline
\esp{Caber} & \esp{cabr-} & \esp{Cabría} & tiendrais (place) \\
\hline
\esp{Saber} & \esp{sabr-} & \esp{Sabría} & saurais \\
\hline
\esp{Poder} & \esp{podr-} & \esp{Podría} & pourrais \\
\hline
\end{tabular}
\end{center}

\textbf{Exemples en conditionnel :} \esp{Si pudiera, lo haría.} « Si je pouvais, je le ferais. » \esp{Si tuviera dinero, compraría un coche.} « Si j'avais de l'argent, j'achèterais une voiture. »
\end{propriete}

\textbf{Conditionnel composé}\par

Il exprime l'antériorité dans l'hypothèse (apodose du type 3).

\begin{propriete}[Conditionnel composé]
\textbf{Formule :} \esp{Habría} (conditionnel de \esp{haber}) + \textbf{Participe passé}.

\begin{center}
\begin{tabular}{|l|l|}
\hline
\rowcolor{poppinkL}
\textbf{Sujet} & \textbf{Forme} \\
\hline
\esp{Yo} & \esp{habría hablado / comido / vivido} \\
\hline
\esp{Tú} & \esp{habrías hablado...} \\
\hline
\esp{Él/Ella/Ud.} & \esp{habría hablado...} \\
\hline
\esp{Nosotros} & \esp{habríamos hablado...} \\
\hline
\emph{Vosotros} & \esp{habríais hablado...} \\
\hline
\esp{Ellos/Uds.} & \esp{habrían hablado...} \\
\hline
\end{tabular}
\end{center}
\end{propriete}

\courssub{Comparaison systématique Français / Espagnol : tableau de correspondance et pièges}

Le passage du français à l'espagnol est source d'erreurs systématiques pour les francophones (calques, confusion des temps). Le tableau suivant aligne les structures françaises standards et leurs équivalents espagnols corrects, en signalant les pièges à éviter absolument.

\begin{center}
\begin{tabularx}{\linewidth}{|>{\raggedright\arraybackslash}X|>{\centering\arraybackslash}X|>{\raggedright\arraybackslash}X|}
\hline
\rowcolor{popblue}
\textbf{Français (Structure standard)} & \textbf{Espagnol (Type \& Formes)} & \textbf{Piège à éviter (Erreur fréquente)} \\
\hline
\textbf{Si + Présent indicatif $\rightarrow$ Futur simple} \\
\emph{« Si j'ai le temps, je viendrai. »} & \textbf{Type 1} : \esp{Si + Présent ind. $\rightarrow$ Futur simple} \\
\esp{Si tengo tiempo, vendré.} & \textcolor{popdark}{\textbf{*Si tendré tiempo, vendré.}} (Futur après \esp{si} \textbf{interdit}) \\
\hline
\textbf{Si + Imparfait indicatif $\rightarrow$ Conditionnel présent} \\
\emph{« Si j'avais le temps, je viendrais. »} & \textbf{Type 2} : \esp{Si + Subj. imparfait $\rightarrow$ Cond. simple} \\
\esp{Si tuviera tiempo, vendría.} & \textcolor{popdark}{\textbf{*Si tenía tiempo, vendría.}} (Indicatif imparfait après \esp{si} = habitude, pas hypothèse) \\
& \textcolor{popdark}{\textbf{*Si tuviera tiempo, vendré.}} (Mélange Type 2/1) \\
\hline
\textbf{Si + Plus-que-parfait indicatif $\rightarrow$ Conditionnel passé} \\
\emph{« Si j'avais eu le temps, je serais venu. »} & \textbf{Type 3} : \esp{Si + Subj. plus-que-parfait $\rightarrow$ Cond. composé} \\
\esp{Si hubiera tenido tiempo, habría venido.} & \textcolor{popdark}{\textbf{*Si había tenido tiempo, habría venido.}} (Indicatif plus-que-parfait après \esp{si} incorrect) \\
& \textcolor{popdark}{\textbf{*Si hubiera tenido tiempo, vine.}} (Indicatif passé dans apodose) \\
\hline
\textbf{Si + Présent (vérité générale) $\rightarrow$ Présent} \\
\emph{« Si on chauffe l'eau, elle bout. »} & \textbf{Type 1 (valeur générale)} : \esp{Si + Présent $\rightarrow$ Présent} \\
\esp{Si calientas agua, hierve.} & \textcolor{popdark}{\textbf{*Si calientas agua, hervirá.}} (Futur inutile pour loi physique) \\
\hline
\end{tabularx}
\end{center}

\begin{exemplebox}[Trois phrases pivots : Thème / Version]
\begin{center}
\begin{tabularx}{\linewidth}{|>{\raggedright\arraybackslash}X|>{\centering\arraybackslash}X|>{\raggedright\arraybackslash}X|}
\hline
\rowcolor{popgoldL}
\textbf{Français} & \textbf{Espagnol (Type)} & \textbf{Analyse grammaticale} \\
\hline
\textbf{1.} « Si j'ai le temps, je viendrai. » & \esp{Si tengo tiempo, vendré.} & \textbf{Type 1} : \esp{tengo} (Présent ind.) $\rightarrow$ \esp{vendré} (Futur simple, irrégulier \esp{venir}). \\
\hline
\textbf{2.} « Si j'avais le temps, je viendrais. » & \esp{Si tuviera tiempo, vendría.} & \textbf{Type 2} : \esp{tuviera} (Subj. imperf. \esp{tener}, radical \esp{tuv-}) $\rightarrow$ \esp{vendría} (Cond. simple, irrégulier \esp{venir}). \\
\hline
\textbf{3.} « Si j'avais eu le temps, je serais venu. » & \esp{Si hubiera tenido tiempo, habría venido.} & \textbf{Type 3} : \esp{hubiera tenido} (Subj. plus-que-parfait) $\rightarrow$ \esp{habría venido} (Cond. composé). \\
\hline
\end{tabularx}
\end{center}
\end{exemplebox}

\courssub{Cas particuliers et nuances avancées}

Au-delà des trois types canoniques, l'espagnol présente des constructions spécifiques que le candidat de Terminale doit reconnaître et savoir employer, notamment à l'écrit (article de presse, essai) et à l'oral (exposé).

\textbf{Si + Présent indicatif : valeur d'habitude ou de vérité générale}\par

Lorsque la condition exprime une \textbf{loi physique, une habitude répétée} ou une \textbf{vérité générale}, les deux verbes sont au \textbf{présent de l'indicatif}. Ce n'est pas une hypothèse future, mais une constance.

\begin{attention}[Piège majeur : Ne pas traduire par le conditionnel !]
\begin{itemize}
    \item \esp{Si caliento agua a 100\,°C, hierve.} « Si je chauffe de l'eau à 100\,°C, elle bout. » $\rightarrow$ Loi physique. \textbf{Pas} *\esp{herviría}.
    \item \esp{Si llueve, no salgo.} « S'il pleut, je ne sors pas / Chaque fois qu'il pleut, je ne sors pas. » $\rightarrow$ Habitude. \textbf{Pas} *\esp{no saldría}.
    \item \emph{Test de substitution :} On peut remplacer \esp{si} par \esp{cuando} (quand / chaque fois que) sans changer le sens. \esp{Cuando llueve, no salgo.}
\end{itemize}
\end{attention}

\textbf{Inversion de l'apodose et de la protase}\par

L'ordre des propositions est libre en espagnol (comme en français). La virgule disparaît généralement si l'apodose précède.

\begin{exemplebox}[Inversion]
\begin{itemize}
    \item \esp{Vendré si tengo tiempo.} « Je viendrai si j'ai le temps. » — Type 1.
    \item \esp{Me iría si pudiera.} « Je partirais si je pouvais. » — Type 2.
    \item \esp{Habríamos ganado si hubiéramos jugado mejor.} « Nous aurions gagné si nous avions joué mieux. » — Type 3.
\end{itemize}
\end{exemplebox}

\textbf{Omission de \texttt{si} par inversion du sujet (style littéraire / formel)}\par

Dans la langue écrite soignée (presse, littérature, discours officiel), on peut supprimer \esp{si} en inversant le verbe et le sujet dans la protase. Cela marque un registre soutenu.

\begin{propriete}[Ellipse de \texttt{si} + Inversion (Registre formel)]
\begin{itemize}
    \item \textbf{Type 2 :} \esp{Si yo fuera rico...} $\rightarrow$ \esp{\textbf{Fuera yo} rico...} « Fusse-je riche... ».
    \item \textbf{Type 3 :} \esp{Si hubieras venido...} $\rightarrow$ \esp{\textbf{Hubieras venido}...} / \esp{\textbf{Hubieras} tú venido...} « Eussiez-vous venu... ».
\end{itemize}
\emph{Attention :} Cette structure n'existe qu'avec le subjonctif (types 2 et 3), jamais avec l'indicatif (type 1).
\end{propriete}

\begin{savaistu}[Histoire de la condition : l'ancien \texttt{si} + Futur indicatif]
Au Moyen Âge et au Siècle d'Or (XVI\textsuperscript{e}-XVII\textsuperscript{e} s.), l'espagnol utilisait encore le \textbf{futur indicatif après \texttt{si}} pour exprimer une condition future réalisée (\esp{Si vendrá, le diré} = « S'il vient, je le lui dirai »). Cette construction a disparu de la norme moderne au profit du \textbf{présent indicatif} (Type 1). On la trouve encore dans certains proverbes figés (\esp{Si quieres la paz, prepara la guerra} — mais ici \esp{quieres} est présent) ou dans la littérature classique (Cervantès, Lope de Vega). Aujourd'hui, \esp{si vendrá} après \esp{si} est une \textbf{faute grave}.
\end{savaistu}

\courssub{Méthode : traduire une phrase conditionnelle (Thème / Version)}

\begin{methode}[Traduire une phrase avec \texttt{si} en 4 étapes]
\begin{enumerate}
    \item \textbf{Identifier les temps français :} Repérer le temps de la proposition \texttt{si} (protase) et de la principale (apodose). Sont-ils : Présent/Futur ? Imparfait/Conditionnel ? Plus-que-parfait/Conditionnel passé ?
    \item \textbf{Choisir le type espagnol correspondant :}
        \begin{itemize}
            \item Présent / Futur $\rightarrow$ \textbf{Type 1} (\esp{Si + Présent Ind. $\rightarrow$ Futur}).
            \item Imparfait / Conditionnel $\rightarrow$ \textbf{Type 2} (\esp{Si + Subj. Imperf. $\rightarrow$ Cond. Simple}).
            \item Plus-que-parfait / Conditionnel passé $\rightarrow$ \textbf{Type 3} (\esp{Si + Subj. Plus-que-parfait $\rightarrow$ Cond. Composé}).
        \end{itemize}
    \item \textbf{Conjuguer correctement les deux verbes :} Vérifier les irréguliers (radicaux conditionnel/futur, radicaux subjonctif imparfait sur base PIS, participes irréguliers : \esp{escrito, puesto, visto, hecho, dicho, muerto, roto, vuelto, abierto, cubierto}).
    \item \textbf{Vérifier les interdits absolus :} \textbf{Jamais} de Futur ni de Conditionnel après \esp{si}. \textbf{Jamais} d'Indicatif (Imparfait, Plus-que-parfait) après \esp{si} pour l'irréel (Types 2 et 3).
\end{enumerate}
\end{methode}

\courssub{Application : texte d'entraînement et exercice résolu}

Pour consolider la règle, lisons un court article de presse original (style \emph{El País} / \emph{La Nación}) qui intègre naturellement les trois types de conditionnels. Ce texte sert de modèle pour votre propre rédaction (Section 3) et de support pour l'exercice résolu.

\begin{textebox}[Article de presse modèle : « La educación en la era digital »]
\textbf{¿Claves para el futuro? La brecha digital en las aulas camerunesas}

\textbf{YAUNDÉ.} Si el gobierno \textbf{invierte} masivamente en infraestructura tecnológica, la brecha digital entre las zonas urbanas y rurales \textbf{se reducirá} notablemente en los próximos cinco años. Esta es la conclusión principal del informe presentado ayer por el Ministerio de Educación Básica (MINEDUB) en colaboración con la UNESCO.

Actualmente, muchas escuelas primarias en el Extremo Norte o en el Este carecen de electricidad estable. \textbf{Si tuvieran} acceso a internet de alta velocidad, los estudiantes \textbf{podrían} consultar bibliotecas virtuales y seguir cursos a distancia, como ya ocurre en Douala o Yaoundé. Los expertos insisten: \textbf{si no se actúa} ahora, una generación entera \textbf{quedará} rezagada frente a los desafíos de la economía del conocimiento.

El ministro declaró: « \textbf{Si hubiéramos empezado} este programa hace una década, hoy \textbf{habríamos formado} a miles de jóvenes en competencias digitales. No podemos perder más tiempo. » La comunidad educativa espera que el plan de acción, dotado de 15 mil millones de francos CFA, se ejecute sin retrasos burocráticos. \textbf{Si se cumplen} los plazos, el próximo curso escolar \textbf{traerá} tabletas y conectividad a 500 escuelas piloto.
\end{textebox}

\begin{definition}[Glossaire : Mots difficiles du texte]
\begin{itemize}
    \item \esp{La brecha digital} : La fracture numérique.
    \item \esp{El Extremo Norte / el Este} : Régions du Cameroun (Extrême-Nord, Est).
    \item \esp{Quedar(se) rezagado} : Être laissé pour compte, prendre du retard.
    \item \esp{La economía del conocimiento} : L'économie de la connaissance.
    \item \esp{Dotado de} : Doté de, pourvu de.
    \item \esp{Retrasos burocráticos} : Retards bureaucratiques / lenteurs administratives.
    \item \esp{Cumplir los plazos} : Respecter les délais / échéances.
\end{itemize}
\end{definition}

\begin{exoresolu}[Thème conditionnel : du français vers l'espagnol]
\textbf{Énoncé :} Traduisez les phrases suivantes en espagnol en justifiant le type employé.

1. Si le ministre vient demain, nous lui poserons la question. \\
2. Si j'étais à ta place, j'accepterais la bourse d'études. \\
3. Si les professeurs avaient eu des tablettes, les élèves auraient appris plus vite. \\
4. Si on chauffe le cacao, il fond. (Vérité générale). \\
5. S'il avait plu, le match n'aurait pas eu lieu.

\tcblower

\textbf{Solution détaillée :}

\begin{enumerate}
    \item \esp{Si el ministro \textbf{viene} mañana, le \textbf{haremos} la pregunta.}
    \begin{itemize}
        \item \textbf{Analyse FR :} \texttt{Si + Présent (vient) $\rightarrow$ Futur simple (poserons)}.
        \item \textbf{Type ES :} \textbf{Type 1 (Réalisable)}. \esp{Si + Présent Ind. $\rightarrow$ Futur Simple}.
        \item \textbf{Conjugaison :} \esp{venir} (présent \esp{viene} diphtongue) $\rightarrow$ \esp{viene}. \esp{Hacer} (irrégulier futur/cond. \esp{har-}) $\rightarrow$ \esp{haremos}.
    \end{itemize}

    \item \esp{Si yo \textbf{fuera} tú, \textbf{aceptaría} la beca de estudios.}
    \begin{itemize}
        \item \textbf{Analyse FR :} \texttt{Si + Imparfait (étais) $\rightarrow$ Conditionnel présent (accepterais)}.
        \item \textbf{Type ES :} \textbf{Type 2 (Irréel présent)}. \esp{Si + Subj. Imperf. $\rightarrow$ Cond. Simple}.
        \item \textbf{Conjugaison :} \esp{Ser} (Subj. Imperf. : \esp{fuera} / \esp{fuese} — radical \esp{fuer-} sur \esp{fueron}). \esp{Aceptar} (régulier cond. \esp{aceptaría}). \emph{Note :} \esp{yo} explicite pour lever l'ambiguïté \esp{fuera} (1\textsuperscript{re}/3\textsuperscript{e} pers.).
    \end{itemize}

    \item \esp{Si los profesores \textbf{hubieran tenido} tabletas, los estudiantes \textbf{habrían aprendido} más rápido.}
    \begin{itemize}
        \item \textbf{Analyse FR :} \texttt{Si + Plus-que-parfait (avaient eu) $\rightarrow$ Conditionnel passé (auraient appris)}.
        \item \textbf{Type ES :} \textbf{Type 3 (Irréel passé)}. \esp{Si + Subj. Plus-que-parfait $\rightarrow$ Cond. Composé}.
        \item \textbf{Conjugaison :} \esp{Haber} (Subj. Imperf. \esp{hubieran}) + \esp{tenido} (Part. Passé irrégulier de \esp{tener}). \esp{Habrían} + \esp{aprendido} (régulier \texttt{-ido}).
    \end{itemize}

    \item \esp{Si \textbf{calientas} el cacao, \textbf{se derrite}.} (ou \esp{Si calientas cacao, se derrite.})
    \begin{itemize}
        \item \textbf{Analyse FR :} \texttt{Si + Présent (chauffe) $\rightarrow$ Présent (fond)}. Vérité scientifique / habitude.
        \item \textbf{Type ES :} \textbf{Type 1 (Valeur générale)}. \esp{Si + Présent Ind. $\rightarrow$ Présent Ind.}
        \item \textbf{Piège évité :} Pas de futur (*\esp{se derretirá}), pas de conditionnel (*\esp{se derretiría}). \esp{Derretirse} (pronominal).
    \end{itemize}

    \item \esp{Si \textbf{hubiera llovido}, el partido \textbf{no se habría jugado}.} (ou \esp{hubiese llovido} / \esp{no habría tenido lugar}).
    \begin{itemize}
        \item \textbf{Analyse FR :} \texttt{Si + Plus-que-parfait (avait plu) $\rightarrow$ Conditionnel passé (n'aurait pas eu lieu)}.
        \item \textbf{Type ES :} \textbf{Type 3 (Irréel passé)}.
        \item \textbf{Conjugaison :} \esp{Llover} (verbe diphtonguant \texttt{o $\rightarrow$ ue} au présent, mais régulier au participe : \esp{llovido}). Auxiliaire \esp{haber} au subj. plus-que-parfait \esp{hubiera}. \esp{Jugar} (verbe \texttt{u $\rightarrow$ ue} : \esp{juego}), voix passive réflexive \esp{se jugó} / \esp{se habría jugado}.
    \end{itemize}
\end{enumerate}
\end{exoresolu}

\begin{exercice}[Entraînement personnel : Thème / Version mixte]
Traduisez en espagnol (Thème) ou en français (Version) en identifiant le type.

\textbf{Thème :}
1. Si nous gagnons le match, nous serons champions. \\
2. Si j'avais su la nouvelle, je t'aurais appelé. \\
3. Si tu étudies bien, tu passeras l'examen. \\
4. Si j'étais riche, j'achèterais une maison à Kribi. \\
5. Si le bus arrive à l'heure, nous ne serons pas en retard.

\textbf{Version :}
6. \esp{Si tienes hambre, come algo.} \\
7. \esp{Si tuviera alas, volaría.} \\
8. \esp{Si hubiéramos salido antes, no habríamos perdido el tren.} \\
9. \esp{Si mezclas azul y amarillo, obtienes verde.} \\
10. \esp{Si el presidente firmara el decreto, entraría en vigor mañana.}
\end{exercice}

\begin{corrige}[Exercice n° 5.6]
\textbf{Thème :}
1. \esp{Si ganamos el partido, seremos campeones.} (Type 1 : \esp{ganamos} Prés. / \esp{seremos} Fut. irrégulier \esp{ser}).
2. \esp{Si hubiera sabido la noticia, te habría llamado.} (Type 3 : \esp{hubiera sabido} Subj. Plusq. / \esp{habría llamado} Cond. Composé. \esp{Saber} part. irrég. \esp{sabido}).
3. \esp{Si estudias bien, aprobarás el examen.} (Type 1 : \esp{estudias} Prés. / \esp{aprobarás} Fut.).
4. \esp{Si fuera rico, compraría una casa en Kribi.} (Type 2 : \esp{fuera} Subj. Imperf. \esp{ser} / \esp{compraría} Cond. Simple).
5. \esp{Si el autobús llega a tiempo, no llegaremos tarde.} (Type 1 : \esp{llega} Prés. / \esp{llegaremos} Fut. \emph{Verbe \esp{llegar} intr. \esp{llegar tarde} = arriver en retard}).

\textbf{Version :}
6. « Si tu as faim, mange quelque chose. » (Type 1 valeur impérative/conseil : \esp{Si + Présent $\rightarrow$ Impératif / Présent}).
7. « Si j'avais des ailes, je volerais. » (Type 2 : \esp{tuviera} Subj. Imperf. \esp{tener} / \esp{volaría} Cond. Simple).
8. « Si nous étions partis plus tôt, nous n'aurions pas raté le train. » (Type 3 : \esp{hubiéramos salido} Subj. Plusq. / \esp{habríamos perdido} Cond. Composé).
9. « Si tu mélanges bleu et jaune, tu obtiens du vert. » (Type 1 vérité générale : \esp{mezclas} Prés. / \esp{obtienes} Prés. diphtongue \texttt{e $\rightarrow$ ie}).
10. « Si le président signait le décret, il entrerait en vigueur demain. » (Type 2 : \esp{firmara} Subj. Imperf. / \esp{entraría} Cond. Simple. Hypothèse peu probable ou protocolaire).
\end{corrige}

\coursec{Entraînement traduction (thème/version) et pièges du Bac}

Nous venons d'établir dans la section précédente la règle complète des trois types de conditionnelles et les tableaux de conjugaison. Il reste maintenant à transformer cette connaissance en \cle{réflexe de traduction} : c'est exactement ce que l'épreuve du Baccalauréat exige. Cette dernière partie est donc un véritable \emph{entraînement de compétition} : thème, version, exercices mixtes, analyse d'erreurs et stratégie d'auto-correction.

\begin{savaistu}[La traduction au Bac camerounais]
Au Baccalauréat A, l'épreuve de traduction (thème et version) porte généralement sur cinq phrases, dont au moins deux contiennent une condition. Les verbes les plus fréquemment testés sont : \esp{tener, hacer, poder, querer, saber, venir, ir, haber}. Ce sont précisément les verbes dont le futur, le conditionnel et le subjonctif imparfait sont irréguliers (\esp{tendría, haría, podría, querría, sabría, vendría, iría, hubiera}). Celui qui maîtrise ces huit verbes sécurise la moitié des points de traduction !
\end{savaistu}

\courssub{La stratégie Bac : traduire une phrase avec « si » en cinq étapes}

Avant de plonger dans les exercices, mémorisons la démarche. Une phrase conditionnelle ne se traduit pas mot à mot : elle se \cle{décode} puis se \cle{reconstruit} selon une formule.

\begin{methode}[La stratégie en cinq étapes]
\begin{enumerate}
\item \textbf{Repérer le marqueur temporel} : \esp{ayer} (hier), \esp{ahora} (maintenant), \esp{mañana} (demain), \esp{el año pasado} (l'année dernière), \esp{ya} (déjà). C'est lui qui trahit le type de condition.
\item \textbf{Identifier le type en français} : réel (futur probable), hypothétique au présent (conditionnel), hypothétique au passé (conditionnel passé).
\item \textbf{Appliquer la formule espagnole} : Type 1 : \esp{si + presente → futuro} ; Type 2 : \esp{si + subjuntivo imperfecto → condicional} ; Type 3 : \esp{si + subjuntivo pluscuamperfecto → condicional compuesto}.
\item \textbf{Vérifier les irréguliers} : \esp{tendría, haría, podría, hubiera, hubiese}.
\item \textbf{Relire} : jamais de futur ni de conditionnel directement après \esp{si} !
\end{enumerate}
\end{methode}

\begin{popfigure}
\begin{tikzpicture}[node distance=0.9cm, every node/.style={align=center}]
\node[draw,rounded corners,fill=popblueL,font=\bfseries] (t) {STRATÉGIE BAC : TRADUIRE SI};
\node[draw,rounded corners,fill=popgreenL,below=of t, align=center] (s1){1. Repérer le marqueur temporel\\ (ayer, ahora, mañana, ya)};
\node[draw,rounded corners,fill=popgoldL,below=of s1, align=center] (s2){2. Identifier le type FR\\ (réel / hypothétique présent / passé)};
\node[draw,rounded corners,fill=poporangeL,below=of s2, align=center] (s3){3. Appliquer la formule ES\\ Type 1 : Si + Pres $\rightarrow$ Fut\\ Type 2 : Si + Subj Imperf $\rightarrow$ Cond\\ Type 3 : Si + Subj Plusq $\rightarrow$ Cond Comp};
\node[draw,rounded corners,fill=poppinkL,below=of s3, align=center] (s4){4. Vérifier irréguliers\\ (tendría, haría, hubiera)};
\node[draw,rounded corners,fill=poppurpleL,below=of s4] (s5) {5. Relire : pas de si + futur/cond};
\draw[->,thick,popdark] (t)--(s1);
\draw[->,thick,popdark] (s1)--(s2);
\draw[->,thick,popdark] (s2)--(s3);
\draw[->,thick,popdark] (s3)--(s4);
\draw[->,thick,popdark] (s4)--(s5);
\end{tikzpicture}
\legende{Organigramme de décision : traduire une conditionnelle en cinq étapes.}
\end{popfigure}

\courssub{Thème (français → espagnol) : dix phrases progressives}

\begin{exemplebox}[Trois modèles de thème corrigés]
\begin{enumerate}
\item « Si tu finis tes devoirs, tu sortiras. » → \esp{Si terminas los deberes, saldrás.} (Type 1 : présent → futur régulier de \esp{salir})
\item « Si j'étais riche, j'achèterais une maison. » → \esp{Si fuera rico, compraría una casa.} (Type 2 : subjonctif imparfait → conditionnel)
\item « S'ils avaient étudié, ils auraient réussi. » → \esp{Si hubieran estudiado, habrían aprobado.} (Type 3 : subjonctif plus-que-parfait → conditionnel passé)
\end{enumerate}
\end{exemplebox}

\begin{exercice}[Thème : traduis en espagnol]
\textbf{Type 1 (réalisable)} :
\begin{enumerate}
\item Si tu viens demain à Yaoundé, nous visiterons le marché central.
\item S'il pleut ce soir, nous n'irons pas au stade.
\item Si tu étudies beaucoup, tu réussiras l'examen.
\end{enumerate}

\textbf{Type 2 (irréalisable au présent)} :
\begin{enumerate}
\item Si j'avais de l'argent, je ferais un voyage en Espagne.
\item Si nous pouvions, nous aiderions nos voisins.
\item Si elle savait la vérité, elle serait triste.
\item Si tu voulais, tu pourrais apprendre l'espagnol.
\end{enumerate}

\textbf{Type 3 (irréalisable au passé)} :
\begin{enumerate}
\item Si vous aviez écouté le professeur, vous auriez compris la règle.
\item Si j'étais venu hier, j'aurais vu le match.
\end{enumerate}
\end{exercice}

\begin{corrige}[Exercice de thème]
\begin{enumerate}
\item \esp{Si vienes mañana a Yaundé, visitaremos el mercado central.} (présent irrégulier de \esp{venir} : \esp{vienes} ; futur régulier \esp{visitaremos})
\item \esp{Si llueve esta noche, no iremos al estadio.} (futur irrégulier de \esp{ir} : \esp{iremos})
\item \esp{Si estudias mucho, aprobarás el examen.} (futur régulier)
\item \esp{Si tuviera dinero, haría un viaje a España.} (subj. imparfait de \esp{tener} : \esp{tuviera} ; conditionnel irrégulier de \esp{hacer} : \esp{haría})
\item \esp{Si pudiéramos, ayudaríamos a nuestros vecinos.} (attention à l'accent : \esp{pudiéramos})
\item \esp{Si ella supiera la verdad, estaría triste.} (subj. imparfait irrégulier de \esp{saber} : \esp{supiera} ; conditionnel irrégulier de \esp{estar} : \esp{estaría})
\item \esp{Si quisieras, podrías aprender español.} (subj. imparfait de \esp{querer} : \esp{quisieras} ; conditionnel irrégulier de \esp{poder} : \esp{podrías})
\item \esp{Si hubierais escuchado al profesor, habríais comprendido la regla.} (ou \esp{hubieseis} ; participe passé de \esp{escuchar} : \esp{escuchado})
\item \esp{Si hubiera venido ayer, habría visto el partido.} (participes irréguliers : \esp{venido}, \esp{visto})
\end{enumerate}
\end{corrige}

\courssub{Version (espagnol → français) : dix phrases, dont les formes piégeuses}

\begin{exemplebox}[Trois modèles de version corrigés]
\begin{enumerate}
\item \esp{Si llueve, el partido se cancela.} → « S'il pleut, le match est annulé. »
\item \esp{Si lloviera, el partido se cancelaría.} → « S'il pleuvait, le match serait annulé. »
\item \esp{Si hubiera llovido, el partido se habría cancelado.} → « S'il avait plu, le match aurait été annulé. »
\end{enumerate}
Observe la gradation : le même contenu, trois degrés de réalité, trois formules.
\end{exemplebox}

\begin{exercice}[Version : traduis en français]
\begin{enumerate}
\item \esp{Si tienes hambre, come algo.}
\item \esp{Si hace buen tiempo mañana, jugaremos al fútbol.}
\item \esp{Si yo fuera presidente, construiría más escuelas.}
\item \esp{Si me lo pidieras, te lo daría.}
\item \esp{Si hubiéramos salido antes, no habríamos perdido el autobús.}
\item \esp{Si lo hubiera sabido, no habría venido.}
\item \esp{De tener dinero, compraría un coche.} (inversion littéraire, \esp{si} omis)
\item \esp{De haberlo sabido, habría actuado de otra manera.} (inversion littéraire au passé)
\item \esp{Si se trabaja duro, se consiguen buenos resultados.} (passif avec \esp{se})
\item \esp{Si se hubieran tomado medidas, se habría evitado la crisis.} (passif avec \esp{se} au Type 3)
\end{enumerate}
\end{exercice}

\begin{corrige}[Exercice de version]
\begin{enumerate}
\item « Si tu as faim, mange quelque chose. » (Type 1 avec impératif en apodose, fréquent)
\item « S'il fait beau demain, nous jouerons au football. »
\item « Si j'étais président, je construirais plus d'écoles. »
\item « Si tu me le demandais, je te le donnerais. » (attention à l'ordre des pronoms : \esp{me lo} = « tu me le »)
\item « Si nous étions partis plus tôt, nous n'aurions pas raté le bus. »
\item « Si je l'avais su, je ne serais pas venu. »
\item « Si j'avais de l'argent, j'achèterais une voiture. » La tournure \esp{de + infinitivo} remplace \esp{si + subjuntivo} dans la langue soutenue : \esp{de tener} = \esp{si tuviera}.
\item « Si je l'avais su, j'aurais agi autrement. » \esp{De haberlo sabido} = \esp{si lo hubiera sabido}.
\item « Si l'on travaille dur, on obtient de bons résultats. » Le \esp{se} passif/impersonnel se rend par « on » ou par un passif français.
\item « Si des mesures avaient été prises, la crise aurait été évitée. » Double \esp{se} passif : traduire par deux passifs français.
\end{enumerate}
\end{corrige}

\begin{attention}[La tournure littéraire \esp{de + infinitif}]
Dans la presse et la littérature, on rencontre \esp{de + infinitivo} à la place de la proposition en \esp{si} : \esp{De no llover, iremos al campo} = « S'il ne pleut pas, nous irons à la campagne ». Au passé : \esp{de haber + participio} = \esp{si + subjuntivo pluscuamperfecto}. Repère cette tournure en version : elle vaut une conditionnelle !
\end{attention}

\courssub{Exercices mixtes : complète avec le bon temps (type Bac)}

\begin{exoresolu}[Choix multiple type Bac]
\textbf{Énoncé.} Pour chaque phrase, choisis la bonne forme (a, b ou c) et justifie par le type de condition.

1. \esp{Si tengo tiempo, te \_\_\_\_\_.} a) \esp{ayudo} b) \esp{ayudaré} c) \esp{ayudaría}

2. \esp{Si \_\_\_\_\_ más, habrías aprobado.} a) \esp{estudias} b) \esp{estudiarías} c) \esp{hubieras estudiado}

3. \esp{Si fuera tú, no \_\_\_\_\_ eso.} a) \esp{haré} b) \esp{haría} c) \esp{hice}

4. \esp{Si \_\_\_\_\_ el tren, llegaremos a tiempo.} a) \esp{cogemos} b) \esp{cogeremos} c) \esp{cogeríamos}

\tcblower
\textbf{Solution détaillée.}
\begin{enumerate}
\item b) \esp{ayudaré}. Type 1 : après \esp{si + presente}, l'apodose standard est au futur. La forme a) \esp{ayudo} (présent) est possible pour exprimer une habitude ou une vérité générale (\esp{Si tengo tiempo, te ayudo} = « Quand j'ai le temps, je t'aide »), mais ici le contexte suggère une promesse ponctuelle : futur.
\item c) \esp{hubieras estudiado}. L'apodose \esp{habrías aprobado} (conditionnel passé) impose le Type 3 : subjonctif plus-que-parfait dans la protase.
\item b) \esp{haría}. Type 2 : \esp{si fuera} (subj. imparfait) appelle le conditionnel ; \esp{hacer} a un conditionnel irrégulier : \esp{haría}.
\item a) \esp{cogemos}. Type 1 : jamais de futur (b) ni de conditionnel (c) après \esp{si} ; l'apodose \esp{llegaremos} confirme le futur.
\end{enumerate}
\end{exoresolu}

\begin{exercice}[À toi : complète avec la forme correcte]
\begin{enumerate}
\item \esp{Si \_\_\_\_\_ (venir) mañana, te veré.}
\item \esp{Si yo \_\_\_\_\_ (tener) tiempo, iría contigo.}
\item \esp{Si ellos \_\_\_\_\_ (estudiar), no habrían suspendido.}
\item \esp{Si \_\_\_\_\_ (hacer) calor, abriremos las ventanas.}
\item \esp{Si me lo \_\_\_\_\_ (decir) antes, lo habría sabido.}
\end{enumerate}
\end{exercice}

\begin{corrige}[Exercice mixte]
\begin{enumerate}
\item \esp{vienes} (Type 1, présent irrégulier de \esp{venir})
\item \esp{tuviera} ou \esp{tuviese} (Type 2, subjonctif imparfait de \esp{tener})
\item \esp{hubieran estudiado} (Type 3)
\item \esp{hace} (Type 1, présent de \esp{hacer} impersonnel ; jamais \esp{hará} après \esp{si})
\item \esp{hubieras dicho} (Type 3 ; participe irrégulier de \esp{decir} : \esp{dicho})
\end{enumerate}
\end{corrige}

\courssub{Analyse d'erreurs types : les fautes qui coûtent des points}

\begin{attention}[Erreur n°1 : \esp{si habría}]
Faute reine des francophones, par calque du français « si j'aurais » (familier) ou par attraction du conditionnel : \esp{*Si habría dinero, compraría un coche} est FAUX. Correct : \esp{Si tuviera dinero, compraría un coche.} Règle absolue : après \esp{si}, jamais de conditionnel, jamais de futur.
\end{attention}

\begin{attention}[Erreur n°2 : \esp{fuera} ou \esp{fuese} pour \esp{ser/ir}]
Les verbes \esp{ser} et \esp{ir} partagent le même subjonctif imparfait : \esp{fuera/fueras/fuera...} ou \esp{fuese/fueses/fuese...}. Les deux séries existent et sont correctes : \esp{Si yo fuera rico} = \esp{Si yo fuese rico} = « Si j'étais riche » ; \esp{Si fuera a Douala} = \esp{Si fuese a Douala} = « Si j'allais à Douala ». Seul le contexte décide entre \esp{ser} et \esp{ir}. Au Bac, les deux séries sont acceptées — mais ne mélange pas les deux dans la même phrase.
\end{attention}

\begin{attention}[Erreur n°3 : confusion \esp{tuviera/tuviese}]
\esp{Tuviera} et \esp{tuviese} sont les deux formes du subjonctif imparfait de \esp{tener} : elles sont parfaitement synonymes et interchangeables (\esp{si tuviera} = \esp{si tuviese}). L'erreur consiste à écrire une forme hybride (\esp{*tuviese} mal accentué, \esp{*teniera} calqué sur l'infinitif). La racine est \esp{tuv-} (du prétérit \esp{tuve}), jamais \esp{ten-}.
\end{attention}

\begin{attention}[Erreur n°4 : la vérité générale]
Quand la phrase exprime une loi physique ou une vérité permanente, l'espagnol utilise \esp{si + presente, presente} : \esp{Si calientas hielo, se derrite.} = « Si tu chauffes de la glace, elle fond. » Ne mets pas le futur (\esp{*se derretirá}) : ce serait une prédiction, pas une loi. En français aussi, le présent s'impose.
\end{attention}

\begin{center}
\begin{tabularx}{\linewidth}{|X|X|X|}
\hline
\rowcolor{popblueL} \textbf{Faute produite} & \textbf{Correction} & \textbf{Cause} \\
\hline
\esp{*Si habría tiempo...} & \esp{Si tuviera tiempo...} & Conditionnel interdit après \esp{si} \\
\hline
\esp{*Si vendrás, te veré} & \esp{Si vienes, te veré} & Futur interdit après \esp{si} \\
\hline
\esp{*Si teniera dinero} & \esp{Si tuviera dinero} & Racine irrégulière \esp{tuv-} \\
\hline
\esp{*Si hubieras venir} & \esp{Si hubieras venido} & Participe passé requis (\esp{venido}) \\
\hline
\esp{*Si estudiarías, aprobarías} & \esp{Si estudiaras, aprobarías} & Calque du français « si tu étudierais » \\
\hline
\end{tabularx}
\end{center}

\courssub{Auto-correction guidée : la grille d'analyse d'erreur}

\begin{methode}[Grille d'auto-correction]
Pour chaque phrase ratée, remplis une ligne de la grille suivante :
\begin{enumerate}
\item \textbf{Phrase} : recopie la phrase fautive.
\item \textbf{Type visé} : 1 (réel), 2 (hypothétique présent) ou 3 (hypothétique passé).
\item \textbf{Temps de la protase} : temps attendu / temps que tu as écrit.
\item \textbf{Temps de l'apodose} : temps attendu / temps que tu as écrit.
\item \textbf{Erreur} : nomme-la (calque, irrégulier oublié, accord, participe).
\item \textbf{Correction} : réécris la phrase entière correctement.
\end{enumerate}
Exemple rempli : « Si j'aurais su, je serais venu » → Type 3 ; protase attendue : \esp{hubiera sabido}, écrite : \esp{*habría sabido} ; apodose : \esp{habría venido} (correcte) ; erreur : conditionnel après \esp{si} ; correction : \esp{Si lo hubiera sabido, habría venido.}
\end{methode}

\begin{experience}[Atelier d'auto-correction en binômes]
Échange ta production de l'exercice de thème avec ton voisin. Chacun corrige la copie de l'autre en remplissant la grille d'analyse pour chaque faute relevée. Puis comparez vos grilles : êtes-vous d'accord sur le type visé ? Sur la cause de l'erreur ? Terminez en réécrivant ensemble les trois phrases les plus difficiles au tableau, en annonçant à voix haute la formule utilisée : « Type 2 : \esp{si + subjuntivo imperfecto → condicional} ».
\end{experience}

\begin{aretenir}[L'essentiel de la section]
\begin{itemize}
\item Stratégie en cinq étapes : marqueur temporel → type → formule → irréguliers → relecture.
\item Jamais de futur ni de conditionnel après \esp{si}.
\item \esp{Fuera/fuese}, \esp{tuviera/tuviese} : deux séries correctes, une seule à la fois.
\item \esp{De + infinitivo} = tournure littéraire équivalente de la conditionnelle.
\item Le \esp{se} passif se traduit par « on » ou par un passif français.
\item Les huit verbes à maîtriser absolument : \esp{tener, hacer, poder, querer, saber, venir, ir, haber}.
\end{itemize}
\end{aretenir}

\newpage

\coursec{Activité d'intégration}

\coursec{Leçon E29 : Redacción : artículo de prensa (120 palabras) ; traducción : la condition réalisable et irréalisable}

\courssub{Objectifs de l'activité}
\begin{itemize}
    \item Mobiliser les compétences de \textbf{compréhension de l'écrit}, \textbf{expression écrite} et \textbf{traduction} dans une tâche unique et authentique.
    \item Rédiger un \textbf{article de presse} structuré (titular, copete, cuerpo, conclusión) d'environ \textbf{120 mots} en espagnol, à partir d'un communiqué de presse en français.
    \item Traduire correctement trois types de \textbf{phrases conditionnelles} (réalisable, irréalisable présent, irréalisable passé) extraites du communiqué.
    \item Travailler en équipe (collaboration, relecture par les pairs, évaluation critériée) pour simuler le fonctionnement d'une rédaction journalistique.
\end{itemize}

\courssub{Situation}
\begin{textebox}[Communiqué de presse brut — Mairie de Bafoussam, 15 mars 2025]
\textbf{POUR DIFFUSION IMMÉDIATE}

\textbf{MARCHE DES JEUNES POUR LA PAIX À BAFOUSSAM : 500 PARTICIPANTS, UN ÉLAN NATIONAL}

Bafoussam, le 15 mars 2025. — Ce samedi, la place de la Fédération a vibré au rythme de la « Marche des Jeunes pour la Paix », organisée par le Conseil National de la Jeunesse (CNJ) en partenariat avec la Mairie de Bafoussam. Plus de \textbf{500 jeunes} venus des sept arrondissements du département de la Mifi ont défilé dans le calme et la discipline, brandissant des banderoles aux messages forts : « La paix se construit », « Jeunesse unie, Cameroun fort », « Non à la violence, oui au dialogue ».

La cérémonie d'ouverture a été présidée par \textbf{Son Excellence Roger Nkodo Dang}, Maire de la ville, qui a salué « l'engagement citoyen d'une jeunesse qui refuse la fatalité des conflits ». Il a annoncé la création d'un \textbf{Fonds Municipal d'Appui aux Initiatives de Paix} doté de 10 millions FCFA pour l'année 2025.

L'après-midi a été consacrée à des \textbf{ateliers de médiation} animés par des facilitateurs formés par le Centre Carter et l'ONG locale « Paix et Développement ». Les thèmes abordés : résolution non-violente des conflits, leadership éthique, et prévention du discours de haine sur les réseaux sociaux.

\textbf{Extraits clés du discours du Maire (à traduire pour l'édition hispanophone) :}
\begin{enumerate}
    \item « \emph{Si les jeunes s'engagent pleinement, la paix avancera dans nos quartiers.} »
    \item « \emph{Si le maire avait financé ce programme l'an dernier, nous aurions eu plus d'ateliers dès aujourd'hui.} »
    \item « \emph{Si nous avions plus de moyens financiers, nous ferions plus pour la formation des médiateurs.} »
\end{enumerate}

\textit{Contact presse :} Cellule de Communication, Mairie de Bafoussam — \texttt{presse@mairie-bafoussam.cm} — +237 6XX XX XX XX.
\end{textebox}

\emph{Contexte pédagogique :} Vous êtes journalistes stagiaires au sein de la rédaction hispanophone du journal numérique \textbf{« Horizonte Camerún »}. Le rédacteur en chef vous confie l'édition spéciale « Cultura de Paz ». Votre matière première : ce communiqué officiel en français. Votre mission : produire la version espagnole pour vos lecteurs d'Espagne, d'Amérique latine et de Guinée équatoriale.

\courssub{Consignes}
\emph{Durée totale : 45 minutes. Travail en groupes de 4 élèves (rôles suggérés : Rédacteur en chef / Secrétaire de rédaction / Traducteur / Correcteur / Metteur en page).}

\begin{enumerate}
    \item \textbf{Analyse du document source (5 min).} Lisez le communiqué. Identifiez les \textbf{5 W} (Quién, Qué, Dónde, Cuándo, Por qué, Cómo) et les \textbf{3 phrases conditionnelles} à traduire. Repérez le vocabulaire clé (marche, paix, médiation, maire, budget, ateliers, jeunesse).
    
    \item \textbf{Rédaction de l'article en espagnol (25 min).} Rédigez l'article complet (\textbf{120 mots ± 10\%}) structuré en 4 parties obligatoires :
    \begin{itemize}
        \item \textbf{Titular} (accrocheur, infinitif ou nominal, max 8-10 mots).
        \item \textbf{Copete / Entrada} (1-2 phrases : résumé complet répondant aux 5 W).
        \item \textbf{Cuerpo} (2-3 paragraphes : détails, citations indirectes ou directes, données chiffrées, déroulement).
        \item \textbf{Conclusión} (ouverture, perspective, citation finale ou appel à l'action).
    \end{itemize}
    \emph{Contrainte :} Intégrez naturellement \textbf{au moins une phrase conditionnelle} (réalisable ou irréalisable) issue de votre traduction (étape 3) dans le corps ou la conclusion de l'article.
    
    \item \textbf{Traduction ciblée des 3 conditions (10 min).} Traduisez les 3 phrases du discours du Maire (listées dans le communiqué) en espagnol, en choisissant la structure \textbf{si + indicatif / subjonctif} appropriée. Écrivez les 3 phrases sur une fiche à part pour l'affichage.
    
    \item \textbf{Relecture croisée et affichage (5 min).} Échangez votre production avec un autre groupe. Vérifiez : structure (4 parties), nombre de mots, orthographe (accents, ñ, ¿¡), concordance des temps (si-clauses), registre journalistique (passif réfléchi \emph{se celebró}, impersonnel \emph{se anunció}, vocabulaire précis). Notez vos observations sur une grille d'évaluation fournie.
    
    \item \textbf{Vote du « Meilleur Article » (5 min).} Lecture à haute voix des 2-3 meilleurs articles. Vote à main levée selon 4 critères : \textbf{Structure} (25\%), \textbf{Lexique / Registre} (25\%), \textbf{Grammaire / Conditions} (25\%), \textbf{Originalité / Fluidité} (25\%).
\end{enumerate}

\courssub{Ressources à mobiliser}
\begin{propriete}[Structure type de l'article de presse (Repaso)]
\begin{center}
\begin{tabularx}{\linewidth}{|X|X|X|}
\hline
\rowcolor{popblueL} \textbf{Parte} & \textbf{Función} & \textbf{Pistas lingüísticas} \\
\hline
\textbf{Titular} & Informar y atraer & Sustantivos, infinitivos, presente histórico (\emph{Bafoussam vibra por la paz}). Sin verbo conjugado a veces. \\
\hline
\textbf{Copete} & Responder a las 5 W & Pretérito perfecto / indefinido (\emph{ha tenido lugar, participaron}). Voz pasiva refleja (\emph{se celebró}). \\
\hline
\textbf{Cuerpo} & Desarrollar: datos, testimonios, causas, consecuencias & Pretérito indefinido / imperfecto. Discurso indirecto (\emph{el alcalde anunció que...}). Conectores (\emph{además, por tanto, sin embargo}). \\
\hline
\textbf{Conclusión} & Cerrar con perspectiva & Futuro, condicional, subjuntivo (\emph{se espera que, si hay fondos, se hará más}). \\
\hline
\end{tabularx}
\end{center}
\end{propriete}

\begin{propriete}[La condition (Si-clauses) : Tableau récapitulatif]
\begin{center}
\begin{tabular}{|l|l|l|l|}
\hline
\rowcolor{popblueL} \textbf{Tipo} & \textbf{Si (Protasis)} & \textbf{Principal (Apodosis)} & \textbf{Ejemplo (Contexto lección)} \\
\hline
\textbf{Realizable} (Posible/Futuro) & \textbf{Presente de indicativo} & \textbf{Futuro simple} / Presente / Imperativo & \esp{Si los jóvenes \textbf{se comprometen}, la paz \textbf{avanzará}.} \\
\hline
\textbf{Irrealizable Presente} (Hipótesis actual) & \textbf{Subjuntivo imperfecto} (\emph{-ra/-se}) & \textbf{Condicional simple} & \esp{Si \textbf{tuviéramos} más medios, \textbf{haríamos} más.} \\
\hline
\textbf{Irrealizable Pasado} (Contrario a la realidad) & \textbf{Subjuntivo pluscuamperfecto} & \textbf{Condicional compuesto} & \esp{Si el alcalde \textbf{hubiera financiado}, \textbf{habríamos tenido} más talleres.} \\
\hline
\end{tabular}
\end{center}
\end{propriete}

\begin{attention}[Pièges fréquents pour francophones (Bac)]
\begin{itemize}
    \item \textbf{Si + présent indicatif $\neq$ Futur en français.} \esp{Si llueve, no voy} (Si \emph{il pleut}, je n'irai pas). Le français utilise le présent dans les deux propositions ; l'espagnol met le futur dans la principale.
    \item \textbf{Subjonctif imparfait : forme \emph{-ra} vs \emph{-se}.} Les deux sont correctes (\emph{si tuviera / si tuviese}), mais \emph{-ra} est plus fréquent à l'oral et dans la presse. \textbf{Ne pas confondre avec le conditionnel français} (*\emph{si tendría} $\rightarrow$ \textbf{incorrect}).
    \item \textbf{Plus-que-parfait du subjonctif :} \esp{habría + participio} (conditionnel composé) en principale ; \esp{hubiera + participio} en \emph{si}. Attention à l'accord du participe (\emph{habríamos tenido}).
    \item \textbf{Faux ami :} \esp{Realizable} = « réalisable / possible », pas « réaliste ». \esp{Actual} = « actuel », pas « réel ».
    \item \textbf{Orthographe :} \esp{avanzará} (accent sur la dernière syllabe), \esp{habríamos} (h muet, accent sur i), \esp{tuviéramos} (accent sur i pour la diphtongue).
\end{itemize}
\end{attention}

\begin{definition}[Lexique utile : Paz, Juventud, Medios de comunicación]
\begin{multicols}{2}
\begin{itemize}
    \item \esp{la marcha (por la paz)} — la marche (pour la paix)
    \item \esp{el desfile / la manifestación} — le défilé / la manifestation
    \item \esp{la pancarta / el cartel} — la banderole / l'affiche
    \item \esp{el alcalde / la alcaldesa} — le maire / la mairesse
    \item \esp{el concejo / el ayuntamiento} — le conseil municipal / la mairie
    \item \esp{el fondo municipal} — le fonds municipal
    \item \esp{la mediación / el mediador} — la médiation / le médiateur
    \item \esp{la resolución de conflictos} — la résolution de conflits
    \item \esp{el discurso de odio} — le discours de haine
    \item \esp{la cultura de paz} — la culture de paix
    \item \esp{comprometerse (con)} — s'engager (pour)
    \item \esp{financiar / sufragar} — financer
    \item \esp{dotar (de)} — doter (de)
    \item \esp{el presupuesto / los medios} — le budget / les moyens
\end{itemize}
\end{multicols}
\end{definition}

\courssub{Proposition de production et corrigé commenté}
\begin{exoresolu}[Production modèle — Grupo « Horizonte Joven »]
\textbf{Article de presse (118 mots)} \\
\begin{textebox}[Artículo : « Bafoussam vibra por la paz: 500 jóvenes marchan por el diálogo »]
\textbf{Bafoussam vibra por la paz: 500 jóvenes marchan por el diálogo}

\textbf{La plaza de la Federación acogió ayer la «Marcha por la Paz», donde más de 500 jóvenes de los siete distritos de la Mifi exigieron el fin de la violencia. El alcalde Roger Nkodo Dang inauguró el acto y anunció un fondo municipal de 10 millones de francos CFA para iniciativas pacíficas en 2025.}

\textbf{Durante la jornada, facilitadores del Centro Carter dirigieron talleres de mediación sobre liderazgo ético y prevención del discurso de odio en redes. El edil subrayó: «\esp{Si los jóvenes se comprometen, la paz avanzará en nuestros barrios}».}

\textbf{No obstante, persisten retos.} \esp{Si tuviéramos más recursos, formaríamos a más mediadores}, recordó el alcalde, quien lamentó: «\esp{Si el ayuntamiento hubiera financiado este programa el año pasado, habríamos contado con más talleres hoy}». La juventud camerunesa demuestra que la paz se construye día a día.
\end{textebox}

\tcblower
\textbf{Commentaire détaillé des choix de langue :}

\begin{enumerate}
    \item \textbf{Titular} : Nominal, concis (9 mots), verbe d'action \emph{vibra} (présent historique), lieu + chiffre + sujet + finalité. Efficace pour le web.
    \item \textbf{Copete} : \emph{acogió} (indéfini, événement ponctuel), \emph{exigieron} (indéfini, action accomplie), \emph{anunció} (indéfini), \emph{iniciativas pacíficas} (lexique précis). Les 5 W sont couverts : Qui (500 jóvenes), Qué (Marcha), Dónde (plaza Federación), Cuándo (ayer), Por qué (exigir fin violencia), Cómo (défilé + ateliers).
    \item \textbf{Cuerpo (1er §)} : \emph{dirigieron} (indéfini), \emph{sobre} (thème), \emph{subrayó} (verbe de parole), citation directe intégrée. \textbf{Condition 1 (Réalisable)} : \esp{Si los jóvenes se comprometen (Pres. Ind.), la paz avanzará (Futur)}. Correspond à « Si les jeunes s'engagent, la paix avancera ». \emph{Comprometen} (pronominal), \emph{avanzará} (accent sur dernière syllabe).
    \item \textbf{Cuerpo (2e §) / Conclusión fusionnée} : Transition \emph{No obstante}. \textbf{Condition 2 (Irréalisable présent)} : \esp{Si tuviéramos (Subj. Imperf.), formaríamos (Cond.)}. Correspond à « Si nous avions plus de moyens, nous ferions plus ». \emph{Tuviéramos} (accent sur i pour diphtongue), \emph{formaríamos} (accent sur i). \textbf{Condition 3 (Irréalisable passé)} : \esp{Si el ayuntamiento hubiera financiado (Subj. Plus-que-parf.), habríamos contado (Cond. Composé)}. Correspond à « Si le maire avait financé... ». \emph{Haya} $\rightarrow$ \emph{hubiera} (forme -ra choisie), \emph{financiado} (participe), \emph{habríamos contado} (auxiliaire \emph{haber} + participe). \emph{Contado con} (locution : disposer de).
    \item \textbf{Registre} : Passif réfléchi absent ici mais \emph{se celebró} possible. Vocabulaire précis : \emph{edil} (synonyme journalistique de alcalde), \emph{distritos} (arrondissements), \emph{redes} (réseaux sociaux).
    \item \textbf{Comptage} : 118 mots (dans la fourchette 108-132).
\end{enumerate}

\textbf{Traductions isolées des 3 conditions (Fiche à part) :}
\begin{enumerate}
    \item \esp{Si los jóvenes se comprometen, la paz avanzará en nuestros barrios.} « Si les jeunes s'engagent, la paix avancera dans nos quartiers. »
    \item \esp{Si tuviéramos más recursos, formaríamos a más mediadores.} « Si nous avions plus de moyens, nous formerions plus de médiateurs. »
    \item \esp{Si el ayuntamiento hubiera financiado este programa el año pasado, habríamos contado con más talleres hoy.} « Si la mairie avait financé ce programme l'an dernier, nous aurions eu plus d'ateliers aujourd'hui. »
\end{enumerate}
\end{exoresolu}

\courssub{Bilan de l'activité}
\begin{aretenir}[Ce que cette activité a permis de valider]
\begin{itemize}
    \item \textbf{Compétence « Écrire » (Redacción) :} Passage du communiqué brut (français, style administratif) à l'article journalistique (espagnol, style médiatique) : sélection de l'information, hiérarchisation (titular/copete/cuerpo/conclusión), reformulation (discours indirect), respect de la contrainte de longueur (120 mots).
    \item \textbf{Compétence « Traduire » (Traducción) :} Identification du type de condition (réalisable / irréalisable présent / passé) $\rightarrow$ choix de la structure grammaticale espagnole adéquate (Indicatif/Futur ; Subj. Imperf./Cond. ; Subj. Plus-que-parf./Cond. Composé) $\rightarrow$ production de phrases correctes, accentuées, ponctuées.
    \item \textbf{Compétence « Grammaire en contexte » :} L'intégration \textbf{obligatoire} des 3 conditions dans l'article a forcé la manipulation naturelle des \emph{si-clauses}, évitant l'exercice de conjugaison décontextualisé.
    \item \textbf{Compétence « Culture / Citoyenneté » :} Ancrage camerounais réel (Bafoussam, Mifi, CNJ, FCFA, ONG locales) + ouverture hispanophone (lecteurs Espagne/Amérique Latine/Guinée Équatoriale). Lexique de la paix, de la médiation, de la gouvernance locale.
    \item \textbf{Compétence « Travailler ensemble / Évaluer » :} Rôles définis, relecture par les pairs (grille critériée), argumentation pour le vote, acceptation de la critique constructive.
\end{itemize}
\end{aretenir}

\begin{methode}[Pour réinvestir seul(e) à la maison (Devoirs / Bac blanc)]
\begin{enumerate}
    \item Choisissez un fait d'actualité camerounaise récent (sport, culture, environnement, éducation). Rédigez un \textbf{communiqué de presse fictif en français} (10 lignes) contenant \textbf{3 conditions} (1 réalisable, 1 irréalisable présent, 1 irréalisable passé).
    \item Transformez-le en \textbf{article de presse en espagnol (120 mots)} avec les 4 parties.
    \item Traduisez les 3 conditions sur une fiche annexe.
    \item Relisez-vous avec la \textbf{check-list} : \cle{Structure 4 parties} ? \cle{120 mots ±10} ? \cle{5 W dans copete} ? \cle{3 conditions correctes + intégrées} ? \cle{Accents / ñ / ¿¡} ? \cle{Registre journalistique (passif réfléchi, vocabulaire précis)} ?
\end{enumerate}
\end{methode}

\begin{savaistu}[Le saviez-vous ? La presse hispanophone au Cameroun et en Guinée équatoriale]
La Guinée équatoriale est le seul pays africain où l'espagnol est langue officielle. Son journal de référence, \emph{La Gaceta de Guinea Ecuatorial}, et des médias numériques comme \emph{Diario Rombe} ou \emph{Asonga} traitent de l'actualité nationale et sous-régionale (CEMAC). Au Cameroun, bien que l'espagnol ne soit pas officiel, la proximité géographique et économique (frontière, commerce, pétroliers, étudiants à Salamanque/Malabo) crée une demande croissante d'informations en espagnol. Rédiger un article en espagnol sur un événement camerounais (comme la marche de Bafoussam), c'est pratiquer un \textbf{journalisme de pont} : rendre visible l'Afrique francophone dans l'espace médiatique hispanophone mondial. C'est une compétence professionnelle recherchée (ONG, ambassades, médias internationaux, entreprises de la zone CEMAC).
\end{savaistu}

\newpage

\coursec{Méthodes et exercices résolus}

\courssub{Fiche méthode 1 : Analyser et structurer un article de presse}

\begin{methode}[Déconstruire un article de presse en 5 étapes]
\begin{enumerate}
\item \textbf{Repérer le titre} (\esp{titular}) : court, accrocheur, souvent au présent ou sous forme nominale. Il annonce le sujet sans tout dire.
\item \textbf{Lire le chapeau} (\esp{entradilla}) : 2 ou 3 phrases, souvent en caractères gras ou en italique, qui résument l'essentiel et répondent déjà aux questions clés.
\item \textbf{Identifier les 6 questions clés} dans le corps (\esp{cuerpo}) : \esp{¿quién?} (qui), \esp{¿qué?} (quoi), \esp{¿dónde?} (où), \esp{¿cuándo?} (quand), \esp{¿por qué?} (pourquoi), \esp{¿cómo?} (comment). Souligner dans le texte la phrase qui répond à chacune.
\item \textbf{Repérer la conclusion} (\esp{cierre}) : ouverture, conséquence attendue, citation d'un acteur.
\item \textbf{Dégager le point de vue} : article informatif (ton neutre, passé simple, troisième personne) ou article d'opinion (première personne, connecteurs comme \esp{en mi opinión}, \esp{sin embargo}).
\end{enumerate}
\textbf{Réflexes} : un article de presse s'écrit à la troisième personne, au passé simple (\esp{indefinido}) pour les faits et à l'imparfait pour le décor ; les citations sont introduites par \esp{según}, \esp{declaró}, \esp{afirmó}, \esp{explicó}.
\end{methode}

\begin{exoresolu}[Analyse d'un article modèle]
Lisez l'article suivant et répondez en espagnol : a) ¿Cuál es el titular? b) ¿Qué información da la entradilla? c) ¿Quién, qué, dónde, cuándo? d) ¿Qué anuncia el cierre?

\begin{textebox}[Prensa escolar]
\textbf{Un torneo de fútbol une a los barrios de Yaoundé}

\emph{Más de doscientos jóvenes participaron el sábado pasado en un torneo de fútbol organizado por el colegio bilingüe de Mvog-Ada para promover la cultura de la paz.}

El sábado por la mañana, el campo del barrio acogió la final de un torneo que duró tres semanas. Los equipos de seis barrios de la capital compitieron ante un público entusiasta. Según el director del colegio, «el deporte es un lenguaje que une a los jóvenes más allá de sus diferencias». El próximo año se organizará una edición nacional del torneo, anunciaron los responsables.
\end{textebox}
\tcblower
\textbf{a) Titular} : \esp{«Un torneo de fútbol une a los barrios de Yaoundé»}. Il est court, au présent (\esp{une}), et annonce le sujet : le sport comme facteur d'unité.

\textbf{b) Entradilla} : elle donne le nombre de participants (\esp{más de doscientos jóvenes}), la date (\esp{el sábado pasado}), le lieu (\esp{colegio bilingüe de Mvog-Ada}) et le but (\esp{promover la cultura de la paz}). C'est un résumé complet de l'événement.

\textbf{c) Questions clés} :
\begin{itemize}
\item \esp{¿Quién?} : más de doscientos jóvenes y los equipos de seis barrios.
\item \esp{¿Qué?} : un torneo de fútbol por la cultura de la paz.
\item \esp{¿Dónde?} : en Yaoundé, en el campo del barrio de Mvog-Ada.
\item \esp{¿Cuándo?} : el sábado pasado, durante tres semanas.
\end{itemize}

\textbf{d) Cierre} : il annonce une ouverture vers l'avenir : \esp{«El próximo año se organizará una edición nacional»} (futur simple \esp{se organizará} pour un projet).

\textbf{Conclusion} : l'article respecte la structure classique : titular, entradilla en italique, cuerpo avec citation (\esp{según el director}), cierre au futur. Ton informatif, troisième personne, passé simple pour les faits (\esp{participaron}, \esp{acogió}, \esp{compitieron}).
\end{exoresolu}

\courssub{Fiche méthode 2 : Rédiger un article de presse de 120 mots}

\begin{methode}[Écrire son article en 6 étapes]
\begin{enumerate}
\item \textbf{Choisir un fait précis} (un événement local : match, fête scolaire, marché, campagne de propreté) et noter au brouillon les réponses aux 6 questions : \esp{quién, qué, dónde, cuándo, por qué, cómo}.
\item \textbf{Rédiger le titular} : 5 à 8 mots, présent ou forme nominale, sans point final.
\item \textbf{Rédiger l'entradilla} : 2 phrases qui résument tout (qui, quoi, où, quand, pourquoi).
\item \textbf{Développer el cuerpo} : 3 ou 4 phrases au passé simple pour les faits, imparfait pour le décor (\esp{hacía sol}, \esp{la gente esperaba}), plus une citation introduite par \esp{según} ou \esp{declaró}.
\item \textbf{Conclure} : une phrase d'ouverture au futur (\esp{el próximo año se organizará...}).
\item \textbf{Compter les mots} (viser 110--130) et relire : accents, accords, temps, ponctuation \esp{¿...? ¡...!}.
\end{enumerate}
\textbf{Lexique utile} : \esp{acoger} (accueillir), \esp{organizar}, \esp{participar en}, \esp{el público}, \esp{los organizadores}, \esp{con motivo de} (à l'occasion de), \esp{según testigos} (selon des témoins), \esp{promover la paz}.
\end{methode}

\begin{exoresolu}[Rédaction type Bac]
\textbf{Sujet} : \esp{Redacta un artículo de prensa (unas 120 palabras) sobre una jornada de limpieza organizada en tu barrio o en tu colegio. No olvides el titular, la entradilla, el cuerpo y el cierre.}
\tcblower
\textbf{Étape 1 — Plan au brouillon} : quién = los alumnos de Terminale ; qué = jornada de limpieza ; dónde = mercado de Mokolo, Yaoundé ; cuándo = el sábado pasado ; por qué = luchar contra las enfermedades ; cómo = con escobas, bolsas y guantes.

\textbf{Étape 2 — Article rédigé} :

\esp{\textbf{Los jóvenes limpian el mercado de Mokolo}}

\esp{\emph{El sábado pasado, cincuenta alumnos del liceo de Yaoundé organizaron una jornada de limpieza en el mercado de Mokolo para luchar contra las enfermedades.}}

\esp{Desde las ocho de la mañana, los jóvenes recogieron basura, limpiaron las calles y sensibilizaron a los vendedores. Hacía mucho calor, pero el público ayudaba con entusiasmo. Según la directora, «la higiene es responsabilidad de todos». Los comerciantes ofrecieron agua y frutas a los participantes. El próximo mes, los alumnos repetirán la operación en otro barrio de la capital.}

\textbf{Étape 3 — Vérifications} :
\begin{itemize}
\item Comptage : environ 95 mots hors titular et entradilla ; total environ 120 mots : conforme.
\item Temps : faits au passé simple (\esp{recogieron}, \esp{limpiaron}, \esp{ofrecieron}) ; décor à l'imparfait (\esp{hacía}, \esp{ayudaba}) ; ouverture au futur (\esp{repetirán}).
\item Citation introduite par \esp{según la directora} avec guillemets.
\item Accents vérifiés : \esp{limpieza}, \esp{higiene}, \esp{operación}, \esp{sensibilizaron}.
\end{itemize}
\textbf{Conclusion} : la structure en quatre blocs est respectée, les 6 questions sont traitées dès l'entradilla, le ton est informatif : l'article est prêt.
\end{exoresolu}

\courssub{Fiche méthode 3 : Traduire la condition réalisable (si + présent / futur)}

\begin{methode}[La condition réalisable en 3 réflexes]
\begin{enumerate}
\item \textbf{Identifier le type} : la condition est-elle possible, probable ? (« Si tu travailles, tu réussiras ») : c'est une condition \cle{réalisable}.
\item \textbf{Appliquer la correspondance} : le français \emph{si + présent} se traduit par \esp{si + presente de indicativo} ; la conséquence au futur français se traduit par le futur espagnol.
\begin{center}
\begin{tabularx}{\linewidth}{|X|X|}\hline
\rowcolor{popblueL} Français & Español \\ \hline
Si tu travailles, tu réussiras. & \esp{Si trabajas, aprobarás.} \\ \hline
S'il pleut, nous resterons à la maison. & \esp{Si llueve, nos quedaremos en casa.} \\ \hline
\end{tabularx}
\end{center}
\item \textbf{Vérifier l'interdiction absolue} : jamais de futur ni de conditionnel directement après \esp{si} : on ne dit pas \esp{si trabajarás}. Après \esp{si}, toujours le présent de l'indicatif (ou le subjonctif imparfait/plus-que-parfait pour l'irréel).
\end{enumerate}
\textbf{Variante} : la conséquence peut être à l'impératif : \esp{Si tienes tiempo, ven a verme.} (Si tu as le temps, viens me voir.)
\end{methode}

\begin{exoresolu}[Thème : condition réalisable]
Traduisez en espagnol : 1) « Si les élèves écoutent la radio, ils apprendront beaucoup de vocabulaire. » 2) « Si tu viens à Douala demain, nous visiterons le port ensemble. »
\tcblower
\textbf{1) Analyse} : condition possible → \esp{si + presente} ; conséquence au futur → futur simple espagnol.
\begin{itemize}
\item « les élèves écoutent » : \esp{los alumnos escuchan} (présent, 3\up{e} pers. plur.).
\item « ils apprendront » : futur de \esp{aprender} : \esp{aprenderán} (accent sur la dernière syllabe).
\item « beaucoup de vocabulaire » : \esp{mucho vocabulario} (\esp{vocabulario} est masculin singulier).
\end{itemize}
\textbf{Traduction} : \esp{Si los alumnos escuchan la radio, aprenderán mucho vocabulario.}

\textbf{2) Analyse} : même schéma. « demain » = \esp{mañana} (avec tilde sur le ñ). « nous visiterons » : futur de \esp{visitar} : \esp{visitaremos}. « ensemble » : \esp{juntos}.
\begin{itemize}
\item Attention : « si tu viens » = \esp{si vienes} (présent de \esp{venir}, irrégulier : vengo, vienes, viene...).
\end{itemize}
\textbf{Traduction} : \esp{Si vienes a Douala mañana, visitaremos el puerto juntos.}

\textbf{Conclusion} : dans les deux cas, jamais de futur après \esp{si} : c'est le présent de l'indicatif qui porte la condition, le futur porte la conséquence.
\end{exoresolu}

\courssub{Fiche méthode 4 : Traduire l'irréalisable au présent (si + subj. imparfait + conditionnel)}

\begin{methode}[L'irréel du présent en 4 étapes]
\begin{enumerate}
\item \textbf{Identifier} : le français « si + imparfait » avec conditionnel dans l'autre proposition (« Si j'avais de l'argent, je voyagerais ») exprime une condition \cle{irréalisable au présent} (contraire à la réalité actuelle).
\item \textbf{Traduire la subordonnée} : \esp{si + imperfecto de subjuntivo} en \esp{-ra} : \esp{si tuviera}, \esp{si fuera}, \esp{si pudiera}.
\item \textbf{Traduire la principale} : \esp{condicional simple} : \esp{viajaría}, \esp{compraría}, \esp{sería}.
\begin{center}
\begin{tabularx}{\linewidth}{|X|X|}\hline
\rowcolor{popblueL} Français & Español \\ \hline
Si j'avais de l'argent, je voyagerais. & \esp{Si tuviera dinero, viajaría.} \\ \hline
Si j'étais toi, je me tairais. & \esp{Si yo fuera tú, me callaría.} \\ \hline
\end{tabularx}
\end{center}
\item \textbf{Mémoriser les formes clés} du subjonctif imparfait (radical de la 3\up{e} pers. plur. du passé simple + \esp{-ra}) : \esp{tener → tuviera} ; \esp{ser/ir → fuera} ; \esp{poder → pudiera} ; \esp{hacer → hiciera} ; \esp{querer → quisiera} ; \esp{haber → hubiera}.
\end{enumerate}
\textbf{Réflexe} : l'imparfait français après « si » ne se traduit JAMAIS par l'imparfait de l'indicatif espagnol : c'est le piège majeur du Bac.
\end{methode}

\begin{exoresolu}[Thème : irréalisable au présent]
Traduisez : 1) « Si nous avions plus de temps, nous lirions la presse espagnole tous les jours. » 2) « Si j'étais journaliste, j'écrirais des articles sur la paix au Cameroun. »
\tcblower
\textbf{1) Analyse} : « si nous avions » = imparfait français de condition irréelle → \esp{si tuviéramos} (subjonctif imparfait de \esp{tener} : radical \esp{tuv-} + \esp{-iéramos}, accent sur la pénultième). « nous lirions » = conditionnel de \esp{leer} : \esp{leeríamos}. « tous les jours » : \esp{todos los días} (accent sur \esp{días}).

\textbf{Traduction} : \esp{Si tuviéramos más tiempo, leeríamos la prensa española todos los días.}

\textbf{2) Analyse} : « si j'étais » → \esp{si fuera} (subjonctif imparfait irrégulier de \esp{ser} ; identique à celui d'\esp{ir}). « j'écrirais » → conditionnel de \esp{escribir} : \esp{escribiría}. « des articles sur la paix » : \esp{artículos sobre la paz} ; « au Cameroun » : \esp{en Camerún} (accent final).

\textbf{Traduction} : \esp{Si fuera periodista, escribiría artículos sobre la paz en Camerún.}

\textbf{Conclusion} : le couple \esp{si + imperfecto de subjuntivo / condicional} est la signature de l'irréel présent. Vérifier chaque accent : \esp{tuviéramos}, \esp{leeríamos}, \esp{escribiría}, \esp{días}, \esp{Camerún}.
\end{exoresolu}

\courssub{Fiche méthode 5 : Traduire l'irréalisable au passé (si + subj. plus-que-parfait + conditionnel passé)}

\begin{methode}[L'irréel du passé en 4 étapes]
\begin{enumerate}
\item \textbf{Identifier} : le français « si + plus-que-parfait » avec conditionnel passé (« Si tu m'avais prévenu, je serais venu ») exprime un regret sur le passé : condition \cle{irréalisable au passé}.
\item \textbf{Traduire la subordonnée} : \esp{si + pluscuamperfecto de subjuntivo} : \esp{hubiera} + participe passé : \esp{si hubieras avisado}.
\item \textbf{Traduire la principale} : \esp{condicional compuesto} : \esp{habría} + participe passé : \esp{habría venido}.
\begin{center}
\begin{tabularx}{\linewidth}{|X|X|}\hline
\rowcolor{popblueL} Français & Español \\ \hline
Si tu m'avais prévenu, je serais venu. & \esp{Si me hubieras avisado, habría venido.} \\ \hline
S'il avait plu, le match aurait été annulé. & \esp{Si hubiera llovido, el partido se habría anulado.} \\ \hline
\end{tabularx}
\end{center}
\item \textbf{Contrôler les participes irréguliers} : \esp{hecho} (fait), \esp{dicho} (dit), \esp{escrito} (écrit), \esp{visto} (vu), \esp{puesto} (mis), \esp{vuelto} (revenu), \esp{abierto} (ouvert), \esp{muerto} (mort).
\end{enumerate}
\textbf{Réflexe} : \esp{habría} prend toujours un accent (conditionnel de \esp{haber}) ; ne jamais confondre avec \esp{había} (imparfait).
\end{methode}

\begin{exoresolu}[Thème et version : irréalisable au passé]
\textbf{A. Thème} — Traduisez : « Si le journaliste avait vérifié ses sources, il n'aurait pas écrit cet article faux. »

\textbf{B. Version} — Traduisez en français : \esp{Si hubiéramos salido antes, habríamos llegado a tiempo al estadio.}
\tcblower
\textbf{A. Analyse} : « avait vérifié » = plus-que-parfait français de regret → \esp{hubiera verificado} (subjonctif plus-que-parfait). « il n'aurait pas écrit » = conditionnel passé négatif → \esp{no habría escrito} (participe irrégulier de \esp{escribir} : \esp{escrito}). « cet article faux » : \esp{ese artículo falso} ; attention au faux ami : « faux » (mensonger) = \esp{falso}, correct ici.

\textbf{Traduction} : \esp{Si el periodista hubiera verificado sus fuentes, no habría escrito ese artículo falso.}

\textbf{B. Analyse de la version} : \esp{hubiéramos salido} = subjonctif plus-que-parfait de \esp{salir} → « si nous étions partis » ; \esp{antes} = « plus tôt » ; \esp{habríamos llegado} = conditionnel passé → « nous serions arrivés » ; \esp{a tiempo} = « à l'heure » (faux ami : \esp{a tiempo} ne signifie pas « à temps partiel »).

\textbf{Traduction} : « Si nous étions partis plus tôt, nous serions arrivés à l'heure au stade. »

\textbf{Conclusion} : dans les deux sens de traduction, repérer d'abord le couple de temps : \esp{hubiera + participe} côté \esp{si}, \esp{habría + participe} côté conséquence. La correspondance avec le français est régulière : plus-que-parfait / conditionnel passé.
\end{exoresolu}

\courssub{Fiche méthode 6 : Choisir le bon schéma de condition (transformation)}

\begin{methode}[L'arbre de décision des si-clauses]
Face à une phrase conditionnelle, poser deux questions :
\begin{enumerate}
\item \textbf{La condition est-elle réalisable ?} Oui → \esp{si + presente} / futur ou présent ou impératif. Exemple : \esp{Si estudias, aprobarás.}
\item \textbf{Elle est irréalisable : à quel moment ?}
\begin{itemize}
\item Par rapport au \textbf{présent} → \esp{si + imperfecto de subjuntivo} / \esp{condicional simple}. Exemple : \esp{Si estudiaras, aprobarías.} (Si tu travaillais, tu réussirais.)
\item Par rapport au \textbf{passé} → \esp{si + pluscuamperfecto de subjuntivo} / \esp{condicional compuesto}. Exemple : \esp{Si hubieras estudiado, habrías aprobado.} (Si tu avais travaillé, tu aurais réussi.)
\end{itemize}
\item \textbf{Transformer} une phrase d'un type à l'autre : garder le même verbe et ne changer que les deux temps, en respectant le couple.
\item \textbf{Interdiction} : après \esp{si}, jamais de présent de subjonctif, jamais de futur, jamais de conditionnel.
\end{enumerate}
\end{methode}

\begin{popfigure}
\begin{center}
\begin{tikzpicture}[
  node distance=0.9cm,
  q/.style={rectangle, rounded corners, draw=popblue, fill=popblueL, align=center, font=\small, inner sep=5pt},
  r/.style={rectangle, rounded corners, draw=popgreen, fill=popgreenL, align=center, font=\small, inner sep=5pt},
  i/.style={rectangle, rounded corners, draw=poporange, fill=poporangeL, align=center, font=\small, inner sep=5pt},
  p/.style={rectangle, rounded corners, draw=poppurple, fill=poppurpleL, align=center, font=\small, inner sep=5pt},
  fleche/.style={-{Stealth[length=2.5mm]}, thick, popdark}
]
\node[q, align=center] (depart){Phrase avec \esp{si} :\\ la condition est-elle réalisable ?};
\node[r, below left=0.9cm and -1.2cm of depart, align=center] (reel){OUI : réel\\ \esp{si + presente}\\ futur / présent / impératif};
\node[q, below right=0.9cm and -1.2cm of depart, align=center] (irreel){NON : irréel\\ à quel moment ?};
\node[i, below left=0.9cm and -0.6cm of irreel, align=center] (pres){Présent\\ \esp{si + imp. subj.}\\ \esp{+ condicional}};
\node[p, below right=0.9cm and -0.6cm of irreel, align=center] (pas){Passé\\ \esp{si + hubiera + part.}\\ \esp{+ habría + part.}};
\draw[fleche] (depart) -- (reel);
\draw[fleche] (depart) -- (irreel);
\draw[fleche] (irreel) -- (pres);
\draw[fleche] (irreel) -- (pas);
\end{tikzpicture}
\end{center}
\legende{Arbre de décision : les trois schémas de la condition en espagnol.}
\end{popfigure}

\begin{exoresolu}[Transformations type Bac]
\textbf{Consigne} : transformez les phrases selon le modèle. Modèle : \esp{Si tengo dinero, compraré el periódico} (irréel du présent) → \esp{Si tuviera dinero, compraría el periódico.}

1) \esp{Si llueve mañana, no iremos al mercado.} (irréel du présent)

2) \esp{Si trabajas más, ganarás más dinero.} (irréel du passé)

3) \esp{Si me escuchas, entenderás la lección.} (irréel du présent)
\tcblower
\textbf{1)} Verbe de la subordonnée : \esp{llueve} (présent) → subjonctif imparfait : \esp{lloviera}. Verbe de la principale : \esp{iremos} (futur) → conditionnel : \esp{iríamos}. Le mot \esp{mañana} devient incohérent avec l'irréel présent : on le supprime ou on le remplace par \esp{ahora}.

\textbf{Solution} : \esp{Si lloviera, no iríamos al mercado.} (S'il pleuvait, nous n'irions pas au marché.)

\textbf{2)} \esp{trabajas} → subjonctif plus-que-parfait : \esp{hubieras trabajado} ; \esp{ganarás} → conditionnel passé : \esp{habrías ganado}.

\textbf{Solution} : \esp{Si hubieras trabajado más, habrías ganado más dinero.} (Si tu avais travaillé davantage, tu aurais gagné plus d'argent.)

\textbf{3)} \esp{escuchas} → \esp{escucharas} ; \esp{entenderás} → \esp{entenderías}.

\textbf{Solution} : \esp{Si me escucharas, entenderías la lección.} (Si tu m'écoutais, tu comprendrais la leçon.)

\textbf{Conclusion} : une transformation réussie ne touche que les deux temps verbaux ; le sens change (réel → regret) mais le vocabulaire reste identique. Toujours vérifier que le couple de temps est complet : subjonctif après \esp{si}, conditionnel dans l'autre proposition.
\end{exoresolu}

\begin{attention}[Les 5 erreurs qui coûtent des points au Bac]
\begin{enumerate}
\item \textbf{Le futur ou le conditionnel après} \esp{si} : écrire \esp{si tendré dinero} ou \esp{si tendría dinero} est une faute grave. Après \esp{si} : présent d'indicatif (réel), subjonctif imparfait (irréel présent) ou subjonctif plus-que-parfait (irréel passé). Jamais autre chose.
\item \textbf{Traduire « si + imparfait » par l'imparfait de l'indicatif} : « Si j'avais de l'argent » ne se traduit pas par \esp{si tenía dinero} mais par \esp{si tuviera dinero}. Le calque du français est l'erreur la plus fréquente des candidats.
\item \textbf{Oublier les accents des temps} : \esp{compraría}, \esp{habría}, \esp{tuviéramos}, \esp{habríamos}, \esp{leería} : le conditionnel et les 1\up{re}/2\up{e} pers. plur. du subjonctif imparfait prennent toujours un accent. Sans accent, le mot change de sens : \esp{habria} n'existe pas, \esp{había} est un imparfait.
\item \textbf{Les faux amis de la presse} : \esp{la prensa} = la presse ; \esp{un artículo} = un article ; « un éditorial » = \esp{un editorial} ; « une actualité » (un fait d'actualité) = \esp{una noticia}, tandis que \esp{la actualidad} désigne l'actualité en général ; \esp{avisar} = prévenir, informer (et non « aviser » au sens de réfléchir).
\item \textbf{L'article de presse sans structure} : au Bac, un texte continu sans titular, sans entradilla et sans citation perd des points de présentation même si la langue est correcte. Toujours afficher les quatre blocs, écrire à la troisième personne et compter les mots (110--130 pour 120 demandés).
\end{enumerate}
\end{attention}

\newpage

\coursec{Exercices}

\courssub{Je vérifie mes connaissances}

\begin{exercice}[QCM : la structure de l'article de presse]
Pour chaque question, choisis la bonne réponse (a, b ou c).

\begin{enumerate}
\item Le \esp{chapeau} d'un article de presse s'appelle en espagnol :
\begin{enumerate}
\item \esp{titular}
\item \esp{entradilla} (ou \esp{copete})
\item \esp{conclusión}
\end{enumerate}
\item Les questions clés auxquelles doit répondre un article sont :
\begin{enumerate}
\item \esp{¿quién?, ¿qué?, ¿dónde?, ¿cuándo?, ¿por qué?, ¿cómo?}
\item \esp{¿quién?, ¿cuánto?, ¿dónde?, ¿adónde?}
\item \esp{¿qué?, ¿de quién?, ¿para qué?}
\end{enumerate}
\item Le titre d'un article doit être :
\begin{enumerate}
\item long et détaillé
\item court, clair et accrocheur
\item une phrase complète au subjonctif
\end{enumerate}
\item Dans une condition irréalisable au présent, après \esp{si}, on emploie :
\begin{enumerate}
\item le présent de l'indicatif
\item le futur simple
\item l'imparfait du subjonctif
\end{enumerate}
\end{enumerate}
\end{exercice}

\begin{exercice}[Vrai ou faux ? Justifie ta réponse]
Dis si les affirmations suivantes sont vraies ou fausses, puis justifie en une phrase.

\begin{enumerate}
\item On peut écrire \esp{Si yo iré mañana, te aviso.}
\item \esp{Si tuviera dinero, compraría ese libro} exprime une condition irréalisable au présent.
\item Le corps de l'article (\esp{cuerpo}) développe les faits et cite des témoins.
\item Après \esp{si} suivi du plus-que-parfait du subjonctif, on trouve le conditionnel passé dans la proposition principale.
\item En espagnol, comme en français, on peut mettre le futur après \esp{si} dans une condition réalisable.
\end{enumerate}
\end{exercice}

\begin{exercice}[Vocabulaire des médias et de la paix]
Complète chaque phrase avec le mot qui convient, choisi dans la liste : \esp{titular, fuente, redactor, edición, entradilla, reconciliación, tolerancia, convivencia, voluntariado, no violencia}.

\begin{enumerate}
\item El \_\_\_\_\_\_\_\_\_\_\_\_ del periódico anuncia la visita del ministro a Yaoundé.
\item El periodista no reveló el nombre de su \_\_\_\_\_\_\_\_\_\_\_\_.
\item La \_\_\_\_\_\_\_\_\_\_\_\_ resume el artículo en dos o tres líneas, justo debajo del título.
\item El \_\_\_\_\_\_\_\_\_\_\_\_ escribió un reportaje sobre la Journée de la Paix en el liceo.
\item La \_\_\_\_\_\_\_\_\_\_\_\_ entre los dos pueblos fue posible gracias al diálogo.
\item Los alumnos participaron en una jornada de \_\_\_\_\_\_\_\_\_\_\_\_ para limpiar el barrio.
\item La \_\_\_\_\_\_\_\_\_\_\_\_ y el respeto mutuo son la base de la cultura de paz.
\end{enumerate}
\end{exercice}

\begin{exercice}[Conjugaison : les temps après \esp{si}]
Conjugue le verbe entre parenthèses au temps qui convient.

\begin{enumerate}
\item Si \_\_\_\_\_\_\_\_\_\_\_\_ (venir, tú) mañana, iremos al estadio.
\item Si yo \_\_\_\_\_\_\_\_\_\_\_\_ (tener) dinero, compraría ese libro.
\item Si ellos \_\_\_\_\_\_\_\_\_\_\_\_ (saber) la verdad, no habrían partido.
\item Si \_\_\_\_\_\_\_\_\_\_\_\_ (llover) mañana, no iremos al parque.
\item Si nosotros \_\_\_\_\_\_\_\_\_\_\_\_ (estudiar) más, habríamos aprobado el examen.
\item Si el gobierno \_\_\_\_\_\_\_\_\_\_\_\_ (invertir) en educación, el país se desarrollará.
\end{enumerate}
\end{exercice}

\courssub{Je m'entraîne}

\begin{exercice}[Traduction : la condition réalisable et irréalisable]
Traduis les phrases suivantes en espagnol.

\begin{enumerate}
\item Si tu viens demain, nous irons au stade.
\item Si j'avais de l'argent, j'achèterais ce livre.
\item S'ils avaient su la vérité, ils ne seraient pas partis.
\item Si le gouvernement investit dans l'éducation, le pays se développera.
\item Si nous avions étudié, nous aurions réussi l'examen.
\end{enumerate}
\end{exercice}

\begin{exercice}[Version : de l'espagnol vers le français]
Traduis les phrases suivantes en français.

\begin{enumerate}
\item \esp{Si estudias, aprobarás.}
\item \esp{Si estudiaras, aprobarías.}
\item \esp{Si hubieras estudiado, habrías aprobado.}
\item \esp{Si no llueve, saldremos.}
\item \esp{Si no lloviera, saldríamos.}
\end{enumerate}
\end{exercice}

\begin{exercice}[Texte à trous : un article de presse]
Complète l'article suivant avec les mots de la liste : \esp{titular, celebró, según, testigos, convivencia, conclusión, entradilla, participaron}.

\begin{textebox}[Un artículo incompleto]
\textbf{El Día de la Paz en el Liceo de Yaoundé}

Ayer, el Liceo de Yaoundé \_\_\_\_\_\_\_\_\_\_\_\_ (1) el Día Internacional de la Paz con una ceremonia emotiva. \_\_\_\_\_\_\_\_\_\_\_\_ (2) el director del establecimiento, más de quinientos alumnos \_\_\_\_\_\_\_\_\_\_\_\_ (3) en las actividades organizadas por el club de español.

Los \_\_\_\_\_\_\_\_\_\_\_\_ (4) afirmaron que los discursos sobre la tolerancia y la \_\_\_\_\_\_\_\_\_\_\_\_ (5) fueron muy aplaudidos. En \_\_\_\_\_\_\_\_\_\_\_\_ (6), la jornada mostró que los jóvenes cameruneses creen en el diálogo y en la no violencia.
\end{textebox}
\end{exercice}

\begin{exercice}[Appariements : mots et champs lexicaux]
Associe chaque mot à son champ lexical : \textbf{Médias}, \textbf{Paix}, \textbf{Jeunesse} ou \textbf{Événement}.

\begin{center}
\begin{tabularx}{\linewidth}{|X|X|}\hline
\rowcolor{popblueL} \textbf{Palabra} & \textbf{Campo léxico} \\ \hline
\esp{titular} & \\ \hline
\esp{reconciliación} & \\ \hline
\esp{redactor} & \\ \hline
\esp{tolerancia} & \\ \hline
\esp{fuente} & \\ \hline
\esp{convivencia} & \\ \hline
\esp{edición} & \\ \hline
\esp{no violencia} & \\ \hline
\esp{copete} & \\ \hline
\esp{voluntariado} & \\ \hline
\end{tabularx}
\end{center}
\end{exercice}

\courssub{Je me prépare au Bac}

\begin{exercice}[Compréhension de texte et questions de langue]
Lis le texte suivant, puis réponds aux questions.

\begin{textebox}[La paz empieza en la escuela]
El pasado viernes, el Colegio Bilingüe de Douala celebró la Journée de la Paix con una jornada dedicada al diálogo y a la tolerancia. Desde las ocho de la mañana, los alumnos de todas las clases participaron en debates, obras de teatro y partidos de fútbol amistosos.

Según la directora, la idea nació de los propios estudiantes: «Si los jóvenes aprenden a escucharse, construirán un país más unido», declaró emocionada. El club de español presentó una obra sobre la reconciliación entre dos barrios enfrentados, que fue muy aplaudida por el público.

Un periodista del periódico local cubrió el evento y entrevistó a varios testigos. «Si hubiera más iniciativas como esta, habría menos conflictos en nuestras ciudades», afirmó un profesor de historia.

Al final de la jornada, los alumnos plantaron un árbol de la paz en el patio del colegio, símbolo de su compromiso con la no violencia y la convivencia.
\end{textebox}

\textbf{Questions de compréhension} (réponds en espagnol, avec des phrases complètes) :
\begin{enumerate}
\item ¿Qué celebró el Colegio Bilingüe de Douala y cuándo?
\item ¿Quién tuvo la idea de la jornada, según la directora?
\item ¿Qué presentó el club de español?
\item ¿Qué hicieron los alumnos al final de la jornada y qué simboliza ese gesto?
\end{enumerate}

\textbf{Questions de langue} :
\begin{enumerate}
\item Releva dans le texte deux phrases conditionnelles et indique, pour chacune, le type de condition (réalisable, irréalisable au présent, irréalisable au passé). Justifie avec les temps employés.
\item Trouve dans le texte trois mots du champ lexical des médias et trois mots du champ lexical de la paix.
\end{enumerate}
\end{exercice}

\begin{exercice}[Thème et version : la condition]
\textbf{A. Thème.} Traduis en espagnol :
\begin{enumerate}
\item Si tu viens demain, nous irons au stade.
\item Si j'avais de l'argent, j'achèterais ce livre.
\item S'ils avaient su la vérité, ils ne seraient pas partis.
\item Si le gouvernement investit dans l'éducation, le pays se développera.
\item Si nous avions étudié, nous aurions réussi l'examen.
\end{enumerate}

\textbf{B. Version.} Traduis en français :
\begin{enumerate}
\item \esp{Si estudias, aprobarás.}
\item \esp{Si estudiaras, aprobarías.}
\item \esp{Si hubieras estudiado, habrías aprobado.}
\item \esp{Si no llueve, saldremos.}
\item \esp{Si no lloviera, saldríamos.}
\end{enumerate}

\textbf{C. Grammaire.} Complète avec le verbe au temps qui convient :
\begin{enumerate}
\item \esp{Si yo \_\_\_\_\_\_\_\_\_\_\_\_ (saber) la verdad, te lo \_\_\_\_\_\_\_\_\_\_\_\_ (decir).}
\item \esp{Si llueve mañana, no \_\_\_\_\_\_\_\_\_\_\_\_ (ir) al parque.}
\item \esp{Si hubiéramos \_\_\_\_\_\_\_\_\_\_\_\_ (llegar) antes, \_\_\_\_\_\_\_\_\_\_\_\_ (ver) el inicio del partido.}
\end{enumerate}

\textbf{D. Correction d'erreurs.} Corrige les phrases suivantes et justifie ta correction :
\begin{enumerate}
\item \esp{Si yo iré, te aviso.}
\item \esp{Si tuviera dinero, compraré el libro.}
\item \esp{Si hubieras venido, iríamos al cine.}
\item \esp{Si hago esto, lo haría mal.}
\end{enumerate}
\end{exercice}

\begin{exercice}[Expression écrite : article de presse]
\textbf{Sujet.} Tu es rédacteur au journal scolaire de ton établissement. Rédige, \textbf{en espagnol}, un article de presse d'environ \textbf{120 mots} sur la Journée de la Paix célébrée dans ton lycée.

\textbf{Consignes obligatoires :}
\begin{itemize}
\item Respecte la structure de l'article de presse : \esp{titular}, \esp{entradilla} (chapeau), \esp{cuerpo}, \esp{conclusión}.
\item Réponds aux questions clés : \esp{¿quién?, ¿qué?, ¿dónde?, ¿cuándo?, ¿por qué?, ¿cómo?}
\item Emploie au moins \textbf{deux phrases conditionnelles} : une condition réalisable (\esp{si} + présent + futur) et une condition irréalisable (\esp{si} + imparfait du subjonctif + conditionnel).
\item Utilise au moins cinq mots du lexique des médias ou de la paix étudiés dans la leçon.
\item Soigne l'orthographe, les accents et la ponctuation espagnole (¿ ? ¡ !).
\end{itemize}
\end{exercice}

\newpage

\coursec{Corrigés des exercices}

\coursec{Leçon E29 : Redacción — Artículo de prensa (120 palabras) ; Traducción : la condición realizable e irrealizable}

\courssub{L'article de presse : structure et questions clés}

\begin{definition}[Structure de l'article de presse en espagnol]
Un article de presse (\esp{artículo de prensa}) suit une structure codifiée en quatre parties obligatoires :
\begin{enumerate}
\item \textbf{El titular} (le titre) : court, accrocheur, informatif. Il résume l'information principale.
\item \textbf{La entradilla} ou \textbf{el copete} (le chapeau) : premier paragraphe, juste sous le titre. Il répond aux questions essentielles (qui, quoi, où, quand) en 2--3 lignes et donne envie de lire la suite.
\item \textbf{El cuerpo} (le corps) : développement des faits, détails, citations de témoins (\esp{testigos}), sources (\esp{fuentes}), données, contexte. Paragraphes courts.
\item \textbf{La conclusión} (la conclusion) : dernier paragraphe, souvent introduit par \esp{En conclusión}, \esp{Por último}, \esp{En definitiva}. Il ouvre une perspective, donne une conséquence ou une citation finale.
\end{enumerate}
\end{definition}

\begin{propriete}[Les six questions clés du journalisme (\esp{Las 6 W})]
Tout article doit répondre explicitement ou implicitement à ces questions dans le chapeau et le corps :
\begin{center}
\begin{tabularx}{\linewidth}{|X|X|X|}\hline
\rowcolor{popblueL} \textbf{Question} & \textbf{Espagnol} & \textbf{Français} \\ \hline
Qui ? & \esp{¿Quién?} & Personnes concernées \\ \hline
Quoi ? & \esp{¿Qué?} & Événement, action \\ \hline
Où ? & \esp{¿Dónde?} & Lieu \\ \hline
Quand ? & \esp{¿Cuándo?} & Moment, date \\ \hline
Pourquoi ? & \esp{¿Por qué?} & Cause, motif \\ \hline
Comment ? & \esp{¿Cómo?} & Manière, circonstances \\ \hline
\end{tabularx}
\end{center}
\end{propriete}

\begin{methode}[Rédiger un article de presse (120 mots)]
\begin{enumerate}
\item \textbf{Analyser le sujet} : identifier l'événement, le public cible, le ton (formel, objectif).
\item \textbf{Répondre aux 6 questions} : noter les réponses brèves (mots-clés) avant d'écrire.
\item \textbf{Rédiger le titre} : nominal, sans verbe conjugué de préférence, percutant.
\item \textbf{Rédiger le chapeau} : 2 phrases max, présent ou passé récent, synthèse totale.
\item \textbf{Développer le corps} : 2--3 paragraphes. Utiliser le passé simple (\esp{pretérito indefinido}) pour les faits accomplis. Intégrer au moins une citation directe (\esp{«...»}) d'un témoin ou d'une autorité.
\item \textbf{Rédiger la conclusion} : 1 phrase, ouverture ou bilan.
\item \textbf{Compter les mots} : viser 110--130 mots. Vérifier accents, \esp{¿ ¡}, accords, concordance des temps.
\end{enumerate}
\end{methode}

\courssub{Compréhension et lexique thématique de l'article modèle}

\begin{textebox}[Modelo de artículo : Jornada de Paz en el Liceo de Yaoundé]
\textbf{El Liceo de Yaoundé celebra el Día de la Paz}

Ayer, los alumnos del Liceo de Yaoundé celebraron el Día Internacional de la Paz con debates, canciones y un partido de fútbol amistoso en el patio del establecimiento.

Según el director, más de seiscientos estudiantes participaron en las actividades organizadas por el club de español. Los discursos sobre la tolerancia y la reconciliación fueron muy aplaudidos por el público. «Si los jóvenes dialogan, construirán un país más unido», declaró una profesora. Un alumno añadió: «Si hubiera más jornadas como esta, habría menos conflictos en nuestros barrios».

En conclusión, esta jornada demostró que la juventud camerunesa cree en la no violencia y en la convivencia. El periódico escolar publicará un reportaje completo en su próxima edición.
\end{textebox}

\begin{definition}[Glossaire : mots difficiles du texte modèle]
\begin{itemize}
\item \esp{celebraron} (pret. indef. 3e pl. de \esp{celebrar}) : ont célébré / fêté.
\item \esp{patio} : cour (de récréation).
\item \esp{Según} : selon (préposition + indicatif).
\item \esp{seiscientos} : six cents (accord en genre : \esp{seiscientas} si fém.).
\item \esp{discursos} : discours.
\item \esp{aplaudidos} : applaudis (part. passé accordé avec \esp{discursos}).
\item \esp{diálogo} : dialogue.
\item \esp{reconciliación} : réconciliation.
\item \esp{no violencia} : non-violence.
\item \esp{convivencia} : coexistence, vie en commun.
\item \esp{reportaje} : reportage.
\item \esp{próxima edición} : prochaine édition / parution.
\end{itemize}
\end{definition}

\begin{exercice}[Compréhension et repérage lexical -- Article modèle]
\begin{enumerate}
\item Identifiez dans le texte le \esp{titular}, la \esp{entradilla}, le \esp{cuerpo} et la \esp{conclusión}.
\item Répondez en espagnol par des phrases complètes :
\begin{enumerate}
\item ¿Quiénes celebraron el Día de la Paz y dónde?
\item ¿Cuántos estudiantes participaron según el director?
\item ¿Qué actividades organizó el club de español?
\item ¿Qué gesto simbólico realizaron los alumnos al final? (Indice : \esp{plantar un árbol}).
\end{enumerate}
\item Relevez dans le texte : 3 mots du champ lexical des \textbf{médias} et 3 mots du champ lexical de la \textbf{paix}.
\item Identifiez les deux phrases conditionnelles du texte. Pour chacune, précisez : type (réalisable / irréalisable), temps après \esp{si}, temps dans la principale.
\end{enumerate}
\end{exercice}

\begin{corrige}[Exercice 2 -- Compréhension et repérage lexical]
\begin{enumerate}
\item \textbf{Structure} : \esp{Titular} = « El Liceo de Yaoundé celebra el Día de la Paz » ; \esp{Entradilla} = « Ayer, los alumnos... del establecimiento » ; \esp{Cuerpo} = les deux paragraphes centraux (débutant par « Según el director... » jusqu'à « ...en nuestros barrios ») ; \esp{Conclusión} = « En conclusión, esta jornada... próxima edición ».
\item \textbf{Réponses} :
\begin{enumerate}
\item \esp{Los alumnos del Liceo de Yaoundé celebraron el Día de la Paz en el patio del establecimiento.}
\item \esp{Más de seiscientos estudiantes participaron según el director.}
\item \esp{El club de español organizó debates, canciones y un partido de fútbol amistoso.}
\item \esp{Los alumnos plantaron un árbol de la paz en el patio; simboliza su compromiso con la no violencia y la convivencia.}
\end{enumerate}
\item \textbf{Champs lexicaux} : Médias -- \esp{periódico, reportaje, edición, director} (aussi \esp{club} comme structure médiatique scolaire) ; Paix -- \esp{tolerancia, reconciliación, no violencia, convivencia, diálogo}.
\item \textbf{Conditions} :
\begin{enumerate}
\item \esp{«Si los jóvenes dialogan, construirán un país más unido»} : \textbf{réalisable} -- \esp{si} + présent indicatif (\esp{dialogan}) + futur simple (\esp{construirán}).
\item \esp{«Si hubiera más jornadas como esta, habría menos conflictos en nuestros barrios»} : \textbf{irréalisable au présent} -- \esp{si} + imparfait subjonctif (\esp{hubiera} de \esp{hay/haber}) + conditionnel présent (\esp{habría}).
\end{enumerate}
\end{enumerate}
\end{corrige}

\courssub{Production guidée : rédiger son article de presse (120 mots)}

\begin{experience}[Atelier rédaction : « La paix dans mon établissement »]
\textbf{Consigne} : En binôme, rédigez un article de presse d'environ 120 mots pour le journal scolaire.
\textbf{Sujet} : Une journée « Culture de la paix » organisée dans votre lycée (Yaoundé, Douala, Bafoussam...).
\textbf{Contraintes obligatoires} :
\begin{itemize}
\item Structure complète : \esp{titular, entradilla, cuerpo, conclusión}.
\item Répondre aux 6 questions (\esp{¿Quién? ¿Qué? ¿Dónde? ¿Cuándo? ¿Por qué? ¿Cómo?}).
\item Inclure \textbf{une condition réalisable} (\esp{si} + présent + futur) et \textbf{une condition irréalisable} (\esp{si} + imparfait subj. + conditionnel).
\item Employer au moins 5 mots du lexique médias/paix vus en classe.
\item Respecter la ponctuation espagnole (\esp{¿ ? ¡ ! « »}).
\end{itemize}
\textbf{Déroulement} : 10 min préparation (plan + lexique) / 20 min rédaction / 10 min relecture croisée (grille d'auto-évaluation) / 5 min mise en commun.
\end{experience}

\begin{aretenir}[Grille d'auto-évaluation de l'article (20 pts)]
\begin{center}
\begin{tabularx}{\linewidth}{|X|c|}\hline
\rowcolor{popblueL} \textbf{Critère} & \textbf{Points} \\ \hline
Structure respectée (4 parties identifiables) & 4 \\ \hline
Réponses aux 6 questions clés & 3 \\ \hline
Condition réalisable correcte (\esp{si} + prés. + fut.) & 3 \\ \hline
Condition irréalisable correcte (\esp{si} + imp. subj. + cond.) & 3 \\ \hline
Lexique médias/paix (5 mots min., genre/nb corrects) & 3 \\ \hline
Correction linguistique (accords, temps, accents, \esp{¿ ¡}) & 3 \\ \hline
Volume (110--130 mots) et présentation soignée & 1 \\ \hline
\end{tabularx}
\end{center}
\end{aretenir}

\courssub{Observation et découverte : la condition en espagnol (si-clauses)}

\begin{textebox}[Phrases d'observation extraites de l'article modèle]
\begin{enumerate}
\item \esp{Si los jóvenes dialogan, construirán un país más unido.}
\item \esp{Si hubiera más jornadas como esta, habría menos conflictos en nuestros barrios.}
\item \esp{Si hubieran participado más padres, la jornada habría sido más completa.} (phrase ajoutée pour l'observation du passé)
\end{enumerate}
\end{textebox}

\begin{exoresolu}[Observation guidée : les trois types de condition]
\begin{enumerate}
\item \textbf{Phrase 1} : \esp{Si los jóvenes dialogan, construirán...}
\begin{itemize}
\item Temps après \esp{si} : \esp{dialogan} = \textbf{présent de l'indicatif}.
\item Temps dans la principale : \esp{construirán} = \textbf{futur simple}.
\item Sens : condition \textbf{réalisable} (possible, probable). « Si les jeunes dialoguent (réalité future possible), ils construiront... »
\end{itemize}
\item \textbf{Phrase 2} : \esp{Si hubiera más jornadas..., habría menos conflictos...}
\begin{itemize}
\item Temps après \esp{si} : \esp{hubiera} = \textbf{imparfait du subjonctif} (verbe \esp{hay/haber}).
\item Temps dans la principale : \esp{habría} = \textbf{conditionnel présent}.
\item Sens : condition \textbf{irréalisable au présent} (hypothétique, contraire à la réalité actuelle). « S'il y avait plus de journées... (mais il n'y en a pas), il y aurait moins de conflits... »
\end{itemize}
\item \textbf{Phrase 3} : \esp{Si hubieran participado más padres, la jornada habría sido más completa.}
\begin{itemize}
\item Temps après \esp{si} : \esp{hubieran participado} = \textbf{plus-que-parfait du subjonctif} (\esp{hubieran} + part. passé).
\item Temps dans la principale : \esp{habría sido} = \textbf{conditionnel passé} (\esp{habría} + part. passé).
\item Sens : condition \textbf{irréalisable au passé} (regret, hypothèse passée non réalisée). « Si plus de parents avaient participé (mais ils n'ont pas participé), la journée aurait été plus complète. »
\end{itemize}
\end{enumerate}
\end{exoresolu}

\courssub{Règle complète, tableaux de formes et comparaison FR/ES}

\begin{propriete}[La condition en espagnol : les trois schémas canoniques]
\textbf{Règle d'or} : \textbf{Jamais de futur, jamais de conditionnel immédiatement après \esp{si}.} Le mot \esp{si} « fige » le temps : il appelle l'indicatif (réel) ou le subjonctif (irréel), jamais le futur/conditionnel.

\begin{center}
\begin{tabularx}{\linewidth}{|X|X|X|X|}\hline
\rowcolor{popblueL} \textbf{Type} & \textbf{Proposition subordonnée (si...)} & \textbf{Proposition principale} & \textbf{Exemple} \\ \hline
\textbf{Réalisable} (futur possible) & \esp{si} + \textbf{Présent indicatif} & \textbf{Futur simple} (ou présent, impératif) & \esp{Si estudias, aprobarás.} \\ \hline
\textbf{Irréalisable présent} (hypothétique maintenant) & \esp{si} + \textbf{Imparfait subjonctif} & \textbf{Conditionnel présent} & \esp{Si estudiara, aprobaría.} \\ \hline
\textbf{Irréalisable passé} (regret, passé non réalisé) & \esp{si} + \textbf{Plus-que-parfait subjonctif} & \textbf{Conditionnel passé} & \esp{Si hubieras estudiado, habrías aprobado.} \\ \hline
\end{tabularx}
\end{center}
\end{propriete}

\begin{definition}[Formation de l'imparfait du subjonctif (rappel)]
Deux séries d'endes \textbf{équivalentes} (la série \textbf{-ra} est la plus fréquente en Amérique et en Guinée équatoriale ; la série \textbf{-se}
\end{definition}


\newpage

\coursec{Fiche bilan}

\begin{popfigure}
\begin{tikzpicture}[
    mindmap,
    concept color=popblue,
    level 1/.style={level distance=3.5cm, sibling angle=360/6},
    level 2/.style={level distance=2.5cm},
    every node/.style={font=\small, align=center},
    branch1/.style={concept color=poporange, grow=30},
    branch2/.style={concept color=popgreen, grow=90},
    branch3/.style={concept color=poppurple, grow=150},
    branch4/.style={concept color=poppink, grow=210},
    branch5/.style={concept color=popgold, grow=270},
    branch6/.style={concept color=popdark, grow=330}
]
\node[concept, text=white, scale=1.2, align=center]{Leçon E29\\Article \& Condition}
    child[branch1] {node[concept, text=white, align=center]{Structure\\Article\\Titre, Chapeau,\\Corps, Conclusion}}
    child[branch2] {node[concept, text=white, align=center]{Les 5 W\\¿Quién? ¿Qué?\\¿Dónde? ¿Cuándo?\\¿Por qué? ¿Cómo?}}
    child[branch3] {node[concept, text=white, align=center]{Condition\\Réalisable\\Si + Présent\\→ Futur}}
    child[branch4] {node[concept, text=white, align=center]{Condition\\Irréalisable Présent\\Si + Subj. Imparfait\\→ Conditionnel}}
    child[branch5] {node[concept, text=white, align=center]{Condition\\Irréalisable Passé\\Si + Subj. PQP\\→ Cond. Passé}}
    child[branch6] {node[concept, text=white, align=center]{Lexique Médias\\Prensa, titular,\\redacción, fuente,\\entrevista, paz}};
\end{tikzpicture}
\legende{Carte mentale : Article de presse et Si-clauses (Leçon E29)}
\end{popfigure}

\begin{bilan}

\end{bilan}

\courssub{Lexique clé : El artículo de prensa}

\begin{center}
\begin{tabularx}{\linewidth}{|X|X|X|}
\hline
\rowcolor{popblueL}
\textbf{Español} & \textbf{Français} & \textbf{Contexte / Exemple} \\
\hline
el titular & le titre (une ligne) & \emph{El titular debe ser breve y llamativo.} \\
el subtítulo / el bajada & le chapeau / sous-titre & \emph{El subtítulo resume la noticia en una frase.} \\
el lead / la entrada & l'attaque (premier paragraphe) & \emph{El lead responde a las 5 W.} \\
el cuerpo & le corps de l'article & \emph{Desarrolla la información con detalles.} \\
la conclusión / el cierre & la conclusion & \emph{Cierra con una cita o consecuencia futura.} \\
la redacción & la rédaction / l'écriture & \emph{La redacción periodística es objetiva.} \\
la fuente & la source & \emph{Citar las fuentes es obligatorio.} \\
el periodista & le journaliste & \emph{El periodista investiga y verifica.} \\
la objetividad & l'objectivité & \emph{Evita adjetivos calificativos subjetivos.} \\
la cultura de la paz & la culture de la paix & \emph{Tema transversal del Módulo 5.} \\
\hline
\end{tabularx}
\end{center}

\courssub{Règle de grammaire à connaître par cœur : Las oraciones condicionales (Si-clauses)}

\begin{center}
\begin{tabularx}{\linewidth}{|>{\bfseries}X|X|X|X|}
\hline
\rowcolor{popblueL}
\textbf{Type} & \textbf{Structure (Si-clause $\rightarrow$ Principal)} & \textbf{Valeur / Probabilité} & \textbf{Exemple traduit} \\
\hline
\textbf{Réalisable} & \esp{Si + Indicatif Présent $\rightarrow$ Futur Simple} & Action possible, probable dans le futur. & \esp{Si llueve, no iremos.} « S'il pleut, nous n'irons pas. » \\
\hline
\textbf{Irréalisable Présent} & \esp{Si + Subjonctif Imparfait $\rightarrow$ Conditionnel Présent} & Hypothèse contraire à la réalité présente, peu probable. & \esp{Si tuviera dinero, compraría el libro.} « Si j'avais de l'argent, j'achèterais le livre. » \\
\hline
\textbf{Irréalisable Passé} & \esp{Si + Subjonctif Plus-que-parfait $\rightarrow$ Conditionnel Passé} & Regret, reproche : action non réalisée dans le passé. & \esp{Si hubieras estudiado, habrías aprobado.} « Si tu avais étudié, tu aurais réussi. » \\
\hline
\end{tabularx}
\end{center}

\courssub{Méthodes express}

\begin{methode}[Rédiger un article (120 mots)]
Trouver un angle d'attaque (actualité locale : Yaoundé, Douala, établissement).
Rédiger un \textbf{titre} accrocheur (nominal, présent de narration).
Écrire le \textbf{chapeau} (1 phrase : Qui ? Quoi ? Où ? Quand ?).
Développer le \textbf{corps} (2-3 paragraphes : Comment ? Pourquoi ? Citations, chiffres).
Conclure par une \textbf{ouverture} (conséquence, avis d'expert, appel à l'action).
Compter les mots (\textasciitilde 120 $\pm$ 10\%). Vérifier accords, accents, \emph{ser/estar}, subjonctif.
\end{methode}

\begin{methode}[Traduire une phrase avec \emph{Si}]
Identifier le temps du verbe français dans la proposition principale (futur, conditionnel, conditionnel passé).
Déduire le type de condition (réalisable / irréalisable présent / passé).
Appliquer la règle espagnole correspondante (Indicatif / Subj. Imparfait / Subj. PQP).
Attention : \textbf{Jamais de conditionnel ni de futur après \emph{si}} en espagnol (sauf \emph{si} = « s'il arrivait que » rare).
Vérifier la concordance des temps (si passé $\rightarrow$ principal passé).
\end{methode}



\begin{aretenir}[Checklist avant le Bac]
$\square$ Je connais la structure en 4 parties de l'article de presse (Titre, Chapeau, Corps, Conclusion).\\
$\square$ Je sais répondre aux 6 questions clés (¿Quién? ¿Qué? ¿Dónde? ¿Cuándo? ¿Por qué? ¿Cómo?) dans le lead.\\
$\square$ J'écris des phrases courtes, objectives, au style direct (pas de « je pense que »).\\
$\square$ Je respecte la consigne de longueur (120 mots $\pm$ 10\%) et je sais compter mes mots.\\
$\square$ \textbf{Condition réalisable :} \esp{Si + Présent Indicatif $\rightarrow$ Futur} (ex: \esp{Si estudias, aprobarás}).\\
$\square$ \textbf{Condition irréalisable présent :} \esp{Si + Subj. Imparfait (-ra/-se) $\rightarrow$ Conditionnel} (ex: \esp{Si estudiara, aprobaría}).\\
$\square$ \textbf{Condition irréalisable passé :} \esp{Si + Subj. Plus-que-parfait $\rightarrow$ Conditionnel Passé} (ex: \esp{Si hubieras estudiado, habrías aprobado}).\\
$\square$ Je ne mets \textbf{jamais} de conditionnel ni de futur après \emph{si} (piège n°1).\\
$\square$ Je distingue \emph{si} (condition) de \emph{sí} (oui) et de \emph{si} (si musical).\\
$\square$ Je maîtrise les irréguliers au Subj. Imparfait (\esp{tuviera, fuera, dijera, pudiera, quisiera, supiera}).\\
$\square$ Je sais former le Subj. Plus-que-parfait (\esp{hubiera + participio} / \esp{hubiese + participio}).\\
$\square$ En version, je traduis le conditionnel français par le conditionnel espagnol, et le futur par le futur.
\end{aretenir}

\findecours

\end{document}
